% Generated mechanically from frozen input p12/ko/build/ch103/stacks-cohomology_full.tex.
% Do not edit this generated file; edit the bound locale source instead.

% OK, start here.
%

\setcounter{chapter}{102}
\chapter{대수적 스택의 코호몰로지}


\phantomsection
\label{stacks-cohomology-section-phantom}


\section{서론}
\label{stacks-cohomology-section-introduction}

\noindent
이 장에서는 대수적 스택의 코호몰로지를 다룬다.
특히 이는 준연접 층의 코호몰로지를 뜻한다. 즉,
『스킴의 코호몰로지』와 『대수공간의 코호몰로지』라는 장에 나오는
결과들의 유사한 명제들을 증명한다.
이 장의 결과들이 \cite{LM-B}의 결과들과 다른 주된 이유는
우리가 일관되게 ``큰 사이트''를 사용하기 때문이다.
이 장을 읽기 전에 『대수적 스택 위의 층』 장을 잠깐 살펴서
그곳에서 도입한 용어에 익숙해지기를 바란다. 『스택 위의 층』의 절
\ref{stacks-sheaves-section-introduction}을 보라.



\section{규약과 용어의 관용적 사용}
\label{stacks-cohomology-section-conventions}

\noindent
『스택의 성질』의 절
\ref{stacks-properties-section-conventions}에서 도입한 규약과
용어의 관용적 사용을 계속 따른다.








\section{표기법}
\label{stacks-cohomology-section-notation}

\noindent
서로 다른 위상들. 대수적 스택을
$\mathcal{X}, \mathcal{Y}, \mathcal{Z}$와 같은 서체 문자로 나타내면,
표기
$\mathcal{X}_{Zar}, \mathcal{X}_\etale, \mathcal{X}_{smooth},
\mathcal{X}_{syntomic}, \mathcal{X}_{fppf}$는
『스택 위의 층』의 정의
\ref{stacks-sheaves-definition-inherited-topologies}에서 도입한
사이트를 뜻한다. (``큰 사이트''라고 생각하라.) 이에 따라
$\mathcal{X}$의 구조층은 $\mathcal{X}_{fppf}$ 위의 층이다.
한편 대수공간과 스킴은 보통 $X, Y, Z$와 같은 로마체 대문자로 나타내며,
이 경우 $X_\etale$은 $X$의 작은 에탈 사이트를 뜻한다(『위상들』의 정의
\ref{topologies-definition-big-small-etale} 또는
『대수공간의 성질』의 정의
\ref{spaces-properties-definition-etale-site}에서 정의한 것이다).
이 구별은 충분히 명확할 것이다.

\medskip\noindent
기본 위상은 fppf 위상이다. 따라서 $\mathcal{X}_{fppf}$ 위의 층이나
$\textit{Mod}(\mathcal{X}_{fppf}, \mathcal{O}_\mathcal{X})$의 대상을
뜻하면서 때때로 ``$\mathcal{X}$ 위의 층'' 또는
``$\mathcal{O}_\mathcal{X}$-가군층''이라고 말할 것이다.

\medskip\noindent
$f : \mathcal{X} \to \mathcal{Y}$가 대수적 스택의 사상이면,
준층에 대해 정의한 함자 $f_*$와 $f^{-1}$은 위에서 언급한 어느 위상에
대해서도 층을 보존한다. 특히 층의 직상이나 당김을 논할 때에는
어느 위상을 사용하는지 언급할 필요가 없다. 코호몰로지 군이나
고차 직상을 계산할 때에는 그렇지 않다. 이 경우에는 항상
어느 위상을 사용하는지 명시할 것이다.

\medskip\noindent
$f : X \to \mathcal{Y}$가 대수공간 $X$에서 대수적 스택
$\mathcal{Y}$로 가는 사상이라고 하자. 어떤 위상 $\tau$에 대해
$\mathcal{G}$를 $\mathcal{Y}_\tau$ 위의 층이라 하자. 이 경우
$f^{-1}\mathcal{G}$는 $\mathcal{S}_X$ 위의 $\tau$ 위상에 대한 층이다
($X$에 대응하는 대수적 스택). 실제로 (우리의 규약에 따라) $f$는
$1$-사상 $f : \mathcal{S}_X \to \mathcal{Y}$이기 때문이다.
$\tau = \etale$이거나 그보다 강한 위상이면,
$f^{-1}\mathcal{G}|_{X_\etale}$로 $X$의 에탈 사이트에 대한 제한을
나타낸다. 『스택 위의 층』의 절
\ref{stacks-sheaves-section-compare}을 보라.
$\mathcal{G}$가 $\mathcal{O}_\mathcal{X}$-가군이면 때때로 그 대신
$f^*\mathcal{G}$와 $f^*\mathcal{G}|_{X_\etale}$라고 쓴다.

\section{준연접 가군의 당김}
\label{stacks-cohomology-section-pullback}

\noindent
$f : \mathcal{X} \to \mathcal{Y}$를 대수적 스택의 사상이라 하자.
환 달린 토포스 위의 준연접 가군이 당김과 양립한다는 것은 매우 일반적인 사실이다.
특히 당김 $f^*$는 준연접 가군을 보존하므로 함자
$$
f^* :
\QCoh(\mathcal{O}_\mathcal{Y})
\longrightarrow
\QCoh(\mathcal{O}_\mathcal{X}),
$$
를 얻는다. 『스택 위의 층』의 보조정리
\ref{stacks-sheaves-lemma-pullback-quasi-coherent}를 보라.
일반적으로 이 함자는 완전하지 않지만, $f$가 평탄하면 완전하다.

\begin{lemma}
\label{stacks-cohomology-lemma-flat-pullback-quasi-coherent}
$f : \mathcal{X} \to \mathcal{Y}$가 대수적 스택의 평탄 사상이면
$f^* : \QCoh(\mathcal{O}_\mathcal{Y}) \to
\QCoh(\mathcal{O}_\mathcal{X})$는 완전 함자이다.
\end{lemma}

\begin{proof}
스킴 $V$와 매끄러운 전사 사상 $V \to \mathcal{Y}$를 택한다.
스킴 $U$와 매끄러운 전사 사상
$U \to V \times_\mathcal{Y} \mathcal{X}$를 택한다.
그러면 $U \to \mathcal{X}$는 그러한 두 사상의 합성이므로 여전히 매끄럽고 전사이다.
가환 도식
$$
\xymatrix{
U \ar[d] \ar[r]_{f'} & V \ar[d] \\
\mathcal{X} \ar[r]^f & \mathcal{Y}
}
$$
으로부터 아벨 범주들의 가환 도식
$$
\xymatrix{
\QCoh(\mathcal{O}_U) & \QCoh(\mathcal{O}_V) \ar[l] \\
\QCoh(\mathcal{O}_\mathcal{X}) \ar[u] &
\QCoh(\mathcal{O}_\mathcal{Y}) \ar[l] \ar[u]
}
$$
을 얻는다. 이 도식의 아래 두 범주가 아벨 범주라는 우리의 증명은
수직 함자들이 충실한 완전 함자임을 보였다(『스택 위의 층』의 보조정리
\ref{stacks-sheaves-lemma-quasi-coherent-algebraic-stack}의 증명을 보라).
$f'$은 스킴의 평탄 사상이므로(대수적 스택의 평탄 사상에 대한 우리의 정의에 의해),
$(f')^*$는 $V$ 위의 준연접 층에 대한 완전 함자이다. 따라서 결론을 얻는다.
\end{proof}

\begin{lemma}
\label{stacks-cohomology-lemma-quasi-coherent-check-exact}
$\mathcal{X}$를 대수적 스택이라 하자. $I$를 집합이라 하고,
$i \in I$에 대해 $x_i : U_i \to \mathcal{X}$를
$\mathcal{X}$의 대상이라 하자. $x_i$가 평탄하고
$\coprod x_i : \coprod U_i \to \mathcal{X}$가 전사라고 가정하자.
$\varphi : \mathcal{F} \to \mathcal{G}$를
$\QCoh(\mathcal{O}_\mathcal{X})$의 화살이라 하자.
$(U_i)_\etale$에 대한 $\varphi$의 제한을 $\varphi_i$라 쓰자.
그러면 $\varphi$가 단사인 것, 각각 전사인 것, 각각 동형인 것은
모든 $\varphi_i$가 그러한 것과 동치이다.
\end{lemma}

\begin{proof}
스킴 $U$와 매끄러운 전사 사상
$x : U \to \mathcal{X}$를 택한다. $x$를 $\mathcal{X}$의 대상으로
생각해도 되고 그렇게 하겠다.
이로부터 공간에서의 어떤 준군 $(U, R, s, t, c)$에 대한 표시
$\mathcal{X} = [U/R]$를 얻고, 이에 대응하여 동치
$$
\QCoh(\mathcal{O}_\mathcal{X}) = \QCoh(U, R, s, t, c)
$$
를 얻는다. 『스택 위의 층』의 절
\ref{stacks-sheaves-section-quasi-coherent-algebraic-stacks}의 논의를 보라.
오른쪽의 아벨 범주 구조에 따르면 $\varphi$가 단사인 것, 각각 전사인 것,
각각 동형인 것은 제한 $\varphi|_{U_\etale}$가 그러한 것과 동치이다.
『공간에서의 준군』의 보조정리 \ref{spaces-groupoids-lemma-abelian}을 보라.

\medskip\noindent
각 $i$에 대해 스킴들에 의한 에탈 피복
$\{W_{i, j} \to V \times_\mathcal{X} U_i\}_{j \in J_i}$를 택한다.
자명한 화살들을 $g_{i, j} : W_{i, j} \to V$와
$h_{i, j} : W_{i, j} \to U_i$로 나타내자.
스킴 사상 $g_{i, j} : W_{i, j} \to U$ 각각은 평탄하고,
이들은 함께 전사이다. 마찬가지로 고정된 각 $i$에 대해 스킴 사상
$h_{i, j} : W_{i, j} \to U_i$들은 평탄하고 함께 전사이다.
『스택 위의 층』의 보조정리 \ref{stacks-sheaves-lemma-quasi-coherent}에 의해,
제한 $\varphi|_{U_\etale}$를 $(g_{i, j})_{small}$로 당긴 것은
제한 $\varphi|_{(W_{i, j})_\etale}$이고,
제한 $\varphi|_{(U_i)_\etale}$를 $(h_{i, j})_{small}$로 당긴 것도
제한 $\varphi|_{(W_{i, j})_\etale}$이다.
스킴의 평탄 사상에 의한 준연접 가군의 당김은 완전하고,
함께 전사인 평탄 스킴 사상들의 족에 의한 당김은 준연접 가군의
단사 사상, 각각 전사 사상, 각각 전단사 사상을 반영한다
(실제로 줄기에서 완전성을 확인할 수 있으므로 이는 모든 가군에 대해 성립한다).
따라서
$$
\varphi|_{U_\etale} \text{가 단사}
\Leftrightarrow
\varphi|_{(W_{i, j})_\etale} \text{가 모든 경우 단사: }i, j
\Leftrightarrow
\varphi|_{(U_i)_\etale} \text{가 모든 경우 단사: }i
$$
임을 알 수 있다. 이로써 증명이 끝난다.
\end{proof}

\section{여러 유형의 가군의 고차 직상}
\label{stacks-cohomology-section-key}

\noindent
다음 보조정리는 특정 유형의 가군층들의 고차 직상을 이해하는 기초가 된다.
두 가지 판본이 있는데, 하나는 에탈 위상에 대한 것이고 다른 하나는
fppf 위상에 대한 것이다.

\begin{lemma}
\label{stacks-cohomology-lemma-general-pushforward}
$\mathcal{M}$을 각 대수적 스택 $\mathcal{X}$에
$\textit{Mod}(\mathcal{X}_\etale, \mathcal{O}_\mathcal{X})$의 부분범주
$\mathcal{M}_\mathcal{X}$를 대응시키는 규칙이라 하자. 다음을 가정한다.
\begin{enumerate}
\item 모든 대수적 스택 $\mathcal{X}$에 대해
$\mathcal{M}_\mathcal{X}$는
$\textit{Mod}(\mathcal{X}_\etale, \mathcal{O}_\mathcal{X})$의
약한 Serre 부분범주이다(『호몰로지』의 정의
\ref{homology-definition-serre-subcategory}를 보라).
\item 대수적 스택의 매끄러운 사상
$f : \mathcal{Y} \to \mathcal{X}$에 대해 함자 $f^*$는
$\mathcal{M}_\mathcal{X}$를 $\mathcal{M}_\mathcal{Y}$ 안으로 보낸다.
\item $f_i : \mathcal{X}_i \to \mathcal{X}$가
$|\mathcal{X}| = \bigcup |f_i|(|\mathcal{X}_i|)$를 만족하는
대수적 스택의 매끄러운 사상들의 족이면,
$\textit{Mod}(\mathcal{X}_\etale, \mathcal{O}_\mathcal{X})$의 대상
$\mathcal{F}$가 $\mathcal{M}_\mathcal{X}$에 속하는 것과
모든 $i$에 대해 $f_i^*\mathcal{F}$가
$\mathcal{M}_{\mathcal{X}_i}$에 속하는 것은 동치이다.
\item $f : \mathcal{Y} \to \mathcal{X}$가 대수적 스택의 사상이고
$\mathcal{X}$와 $\mathcal{Y}$가 아핀 스킴으로 표현가능하면,
$R^if_*$는 $\mathcal{M}_\mathcal{Y}$를
$\mathcal{M}_\mathcal{X}$ 안으로 보낸다.
\end{enumerate}
그러면 대수적 스택의 임의의 준콤팩트 준분리 사상
$f : \mathcal{Y} \to \mathcal{X}$에 대해
$R^if_*$는 $\mathcal{M}_\mathcal{Y}$를
$\mathcal{M}_\mathcal{X}$ 안으로 보낸다. (고차 직상은 에탈 위상에서 계산한다.)
\end{lemma}

\begin{proof}
$f : \mathcal{Y} \to \mathcal{X}$를 대수적 스택의 준콤팩트 준분리 사상이라 하고,
$\mathcal{F}$를 $\mathcal{M}_\mathcal{Y}$의 대상이라 하자.
스킴으로 표현가능한 $\mathcal{U}$와 매끄러운 전사 사상
$\mathcal{U} \to \mathcal{X}$를 택한다.
『스택 위의 층』의 보조정리
\ref{stacks-sheaves-lemma-base-change-higher-direct-images}에 의해
고차 직상을 취하는 것은 밑변환과 가환한다.
조건 (2)에 따르면 $\mathcal{F}$를
$\mathcal{U} \times_\mathcal{X} \mathcal{Y}$로 당긴 것은
$\mathcal{M}_{\mathcal{U} \times_\mathcal{X} \mathcal{Y}}$에 속한다.
실제로 사영
$\mathcal{U} \times_\mathcal{X} \mathcal{Y} \to \mathcal{Y}$는
매끄러운 사상의 밑변환이므로 매끄럽다. 따라서 (3)에 의해
$\mathcal{Y} \to \mathcal{X}$를 사영
$\mathcal{U} \times_\mathcal{X} \mathcal{Y} \to \mathcal{U}$로
바꾸어도 된다. 달리 말해 $\mathcal{X}$가 스킴으로 표현가능하다고 가정해도 된다.
(3)을 한 번 더 사용하면 이 문제는 $\mathcal{X}$ 위에서 Zariski 국소적이므로,
$\mathcal{X}$가 아핀 스킴으로 표현가능하다고 가정해도 된다.
$f$가 준콤팩트이므로 $\mathcal{Y}$도 준콤팩트이다.
따라서 아핀 스킴으로 표현가능한 $\mathcal{V}$와 매끄러운 전사 사상
$g : \mathcal{V} \to \mathcal{Y}$를 택할 수 있다.

\medskip\noindent
이 상황에서는 『스택 위의 층』의 명제
\ref{stacks-sheaves-proposition-smooth-covering-compute-direct-image}의
스펙트럼 열
$$
E_2^{p, q} = R^q(f \circ g_p)_*g_p^*\mathcal{F}
\Rightarrow
R^{p + q}f_*\mathcal{F}
$$
이 있다. 이것은 제1사분면 스펙트럼 열이므로 『호몰로지』의 보조정리
\ref{homology-lemma-first-quadrant-ss}의 마지막 부분을 사용할 수 있음을 상기하라.
사상
$$
g_p : \mathcal{V}_p =
\mathcal{V} \times_\mathcal{Y} \ldots \times_\mathcal{Y} \mathcal{V}
\longrightarrow
\mathcal{Y}
$$
은 매끄러운 사상 $g$의 밑변환들을 합성한 것이므로 매끄럽다.
따라서 (2)에 의해 층 $g_p^*\mathcal{F}$는
$\mathcal{M}_{\mathcal{V}_p}$에 속한다. 그러므로 사상
$$
\mathcal{V}_p =
\mathcal{V} \times_\mathcal{Y} \ldots \times_\mathcal{Y} \mathcal{V}
\longrightarrow
\mathcal{X}
$$
에 의한 $\mathcal{M}_{\mathcal{V}_p}$의 대상들의 고차 직상이
$\mathcal{M}_\mathcal{X}$에 속함을 증명하면 충분하다.
대수적 스택 $\mathcal{V}_p$는 『스택의 사상』의 보조정리
\ref{stacks-morphisms-lemma-quasi-compact-quasi-separated-permanence}에 의해
준콤팩트이고 준분리이다.
물론 각 $\mathcal{V}_p$는 대수공간으로 표현가능하다
(대수적 스택 $\mathcal{Y}$의 대각사상은 대수공간으로 표현가능하다).
따라서 $\mathcal{Y}$가 대수공간으로 표현가능하고
$\mathcal{X}$가 아핀 스킴으로 표현가능한 경우로 환원된다.

\medskip\noindent
$\mathcal{Y}$가 대수공간으로 표현가능하고 $\mathcal{X}$가 아핀 스킴으로
표현가능한 상황에서, 아핀 스킴으로 표현가능한 $\mathcal{V}$와
매끄러운 전사 사상 $\mathcal{V} \to \mathcal{Y}$를 새로 택한다.
위 논증을 다시 밟으면 다시 사상 $\mathcal{V}_p \to \mathcal{X}$의 경우로
환원된다. 그러나 지금 상황에서는 대수적 스택 $\mathcal{V}_p$가
준콤팩트 준분리 스킴으로 표현가능하다
(대수공간의 대각사상은 스킴으로 표현가능하기 때문이다).

\medskip\noindent
따라서 $\mathcal{Y}$가 스킴으로 표현가능하고
$\mathcal{X}$가 아핀 스킴으로 표현가능하다고 가정해도 된다.
(다시) 아핀 스킴으로 표현가능한 $\mathcal{V}$와
매끄러운 전사 사상 $\mathcal{V} \to \mathcal{Y}$를 택한다.
이 경우 모든 대수적 스택 $\mathcal{V}_p$는 분리 스킴으로 표현가능하다
(스킴의 대각사상은 분리이기 때문이다).

\medskip\noindent
따라서 $\mathcal{Y}$가 분리 스킴으로 표현가능하고
$\mathcal{X}$가 아핀 스킴으로 표현가능하다고 가정해도 된다.
(또다시) 아핀 스킴으로 표현가능한 $\mathcal{V}$와
매끄러운 전사 사상 $\mathcal{V} \to \mathcal{Y}$를 택한다.
이 경우 모든 대수적 스택 $\mathcal{V}_p$는 아핀 스킴으로 표현가능하다
(분리 스킴의 대각사상은 닫힌 몰입이므로 아핀이기 때문이다).
이 경우는 조건 (4)로 처리된다. 이로써 증명이 끝난다.
\end{proof}

\noindent
다음은 fppf 위상에 대한 판본이다.

\begin{lemma}
\label{stacks-cohomology-lemma-general-pushforward-fppf}
$\mathcal{M}$을 각 대수적 스택 $\mathcal{X}$에
$\textit{Mod}(\mathcal{O}_\mathcal{X})$의 부분범주
$\mathcal{M}_\mathcal{X}$를 대응시키는 규칙이라 하자. 다음을 가정한다.
\begin{enumerate}
\item 모든 대수적 스택 $\mathcal{X}$에 대해
$\mathcal{O}_\mathcal{X}$는
$\textit{Mod}(\mathcal{O}_\mathcal{X})$의 약한 Serre 부분범주이다.
\item 대수적 스택의 매끄러운 사상
$f : \mathcal{Y} \to \mathcal{X}$에 대해 함자 $f^*$는
$\mathcal{M}_\mathcal{X}$를 $\mathcal{M}_\mathcal{Y}$ 안으로 보낸다.
\item $f_i : \mathcal{X}_i \to \mathcal{X}$가
$|\mathcal{X}| = \bigcup |f_i|(|\mathcal{X}_i|)$를 만족하는
대수적 스택의 매끄러운 사상들의 족이면,
$\textit{Mod}(\mathcal{O}_\mathcal{X})$의 대상 $\mathcal{F}$가
$\mathcal{M}_\mathcal{X}$에 속하는 것과 모든 $i$에 대해
$f_i^*\mathcal{F}$가 $\mathcal{M}_{\mathcal{X}_i}$에 속하는 것은 동치이다.
\item $f : \mathcal{Y} \to \mathcal{X}$가 대수적 스택의 사상이고
$\mathcal{X}$와 $\mathcal{Y}$가 아핀 스킴으로 표현가능하면,
$R^if_*$는 $\mathcal{M}_\mathcal{Y}$를
$\mathcal{M}_\mathcal{X}$ 안으로 보낸다.
\end{enumerate}
그러면 대수적 스택의 임의의 준콤팩트 준분리 사상
$f : \mathcal{Y} \to \mathcal{X}$에 대해
$R^if_*$는 $\mathcal{M}_\mathcal{Y}$를
$\mathcal{M}_\mathcal{X}$ 안으로 보낸다. (고차 직상은 fppf 위상에서 계산한다.)
\end{lemma}

\begin{proof}
보조정리 \ref{stacks-cohomology-lemma-general-pushforward}의 증명과 같다.
\end{proof}

\section{국소 준연접 가군}
\label{stacks-cohomology-section-locally-quasi-coherent}

\noindent
$\mathcal{X}$를 대수적 스택이라 하자. $\mathcal{F}$를
$\mathcal{O}_\mathcal{X}$-가군의 준층이라 하자. $\mathcal{F}$가
{\it 국소 준연접}인지 물을 수 있다. 『스택 위의 층』의 정의
\ref{stacks-sheaves-definition-locally-quasi-coherent}를 보라.
간단히 말해, 이는 $\mathcal{F}$가 에탈 위상에 대한
$\mathcal{O}_\mathcal{X}$-가군이고, 임의의 사상
$f : U \to \mathcal{X}$에 대해 제한
$f^*\mathcal{F}|_{U_\etale}$가 $U_\etale$ 위에서 준연접임을 뜻한다.
(실제 정의는 조금 다르지만 동치이다.) 유용한 사실은
$$
\textit{LQCoh}(\mathcal{O}_\mathcal{X}) \subset
\textit{Mod}(\mathcal{X}_\etale, \mathcal{O}_\mathcal{X})
$$
이 약한 Serre 부분범주라는 것이다. 『스택 위의 층』의 보조정리
\ref{stacks-sheaves-lemma-lqc-colimits}를 보라.

\begin{lemma}
\label{stacks-cohomology-lemma-check-lqc-on-etale-covering}
$\mathcal{X}$를 대수적 스택이라 하자.
$f_j : \mathcal{X}_j \to \mathcal{X}$를
$|\mathcal{X}| =\bigcup |f_j|(|\mathcal{X}_j|)$를 만족하는
대수적 스택의 매끄러운 사상들의 족이라 하자.
$\mathcal{F}$를 $\mathcal{X}_\etale$ 위의
$\mathcal{O}_\mathcal{X}$-가군층이라 하자.
각 $f_j^{-1}\mathcal{F}$가 국소 준연접이면 $\mathcal{F}$도 그러하다.
\end{lemma}

\begin{proof}
각 대수적 스택 $\mathcal{X}_j$를 스킴 $U_j$로 바꾸어도 된다
(모든 대수적 스택은 스킴에 의한 매끄러운 피복을 가지며 매끄러운 사상들의
합성이 매끄럽다는 사실을 사용한다. 『스택의 사상』의 보조정리
\ref{stacks-morphisms-lemma-composition-smooth}를 보라).
$\mathcal{F}$를 $(\Sch/U_j)_\etale$로 당긴 것은 여전히 국소 준연접이다.
『스택 위의 층』의 보조정리 \ref{stacks-sheaves-lemma-pullback-lqc}를 보라.
그러면 $f = \coprod f_j : U = \coprod U_j \to \mathcal{X}$는
매끄러운 전사 사상이다. $x$를 $\mathcal{X}$의 대상이라 하자.
『스택 위의 층』의 보조정리
\ref{stacks-sheaves-lemma-surjective-flat-locally-finite-presentation}에 의해
각 $x_i$가 $(\Sch/U)_\etale$의 어떤 대상 $u_i$로 올려지는 에탈 피복
$\{x_i \to x\}_{i \in I}$가 존재한다. 이는 $x$, $x_i$가 스킴
$V$, $V_i$ 위에 놓이고, $\{V_i \to V\}$가 에탈 피복이며,
$x_i$가 어떤 사상 $u_i : V_i \to U$에서 온다는 뜻일 뿐이다.
제한 $x_i^*\mathcal{F}|_{V_{i, \etale}}$는
$f^*\mathcal{F}$의 $V_{i, \etale}$에 대한 제한과 같다.
『스택 위의 층』의 보조정리 \ref{stacks-sheaves-lemma-comparison}을 보라.
따라서 $x^*\mathcal{F}|_{V_\etale}$는 $V$의 작은 에탈 사이트 위의 층이고,
각 $i$에 대해 $V_{i, \etale}$로 제한하면 준연접이다.
그러므로 이는 준연접이다(원하는 바이다). 예를 들어 『대수공간의 성질』의 보조정리
\ref{spaces-properties-lemma-characterize-quasi-coherent}를 보라.
\end{proof}

\begin{lemma}
\label{stacks-cohomology-lemma-pushforward-locally-quasi-coherent}
$f : \mathcal{X} \to \mathcal{Y}$를 대수적 스택의 준콤팩트 준분리 사상이라 하자.
$\mathcal{F}$를 $\mathcal{X}_\etale$ 위의 국소 준연접
$\mathcal{O}_\mathcal{X}$-가군이라 하자.
그러면 $R^if_*\mathcal{F}$는(에탈 위상에서 계산한 것으로)
$\mathcal{Y}_\etale$ 위에서 국소 준연접이다.
\end{lemma}

\begin{proof}
이를 증명하기 위해 보조정리 \ref{stacks-cohomology-lemma-general-pushforward}를 사용한다.
그 조건 (1)--(4)를 확인하겠다. (1)과 (2)는 『스택 위의 층』의 보조정리
\ref{stacks-sheaves-lemma-lqc-colimits}에서 따른다.
(3)은 보조정리 \ref{stacks-cohomology-lemma-check-lqc-on-etale-covering}에서 따른다.
따라서 (4)를 보이면 충분하다.

\medskip\noindent
$f : \mathcal{X} \to \mathcal{Y}$가 대수적 스택의 사상이고
$\mathcal{X}$와 $\mathcal{Y}$가 아핀 스킴 $X$와 $Y$로 표현가능하다고 하자.
스킴 $V$ 위에 놓인 $\mathcal{Y}$의 임의의 대상 $y$를 택한다.
명료성을 위해 $V$에 대응하는 대수적 스택을
$\mathcal{V} = (\Sch/V)_{fppf}$라 쓰자. 카르테시안 도식
$$
\xymatrix{
\mathcal{Z} \ar[d] \ar[r]_g \ar[d]_{f'} & \mathcal{X} \ar[d]^f \\
\mathcal{V} \ar[r]^y & \mathcal{Y}
}
$$
을 생각하자. 그러면 $\mathcal{Z}$는 스킴 $Z = V \times_Y X$로 표현가능하고,
$f'$은 준콤팩트이고 분리이다(실제로는 아핀이다).
『스택 위의 층』의 보조정리
\ref{stacks-sheaves-lemma-compare-representable-morphism-cohomology}에 의해
$$
R^if_*\mathcal{F}|_{V_\etale} =
R^if'_{small, *}\big(g^*\mathcal{F}|_{Z_\etale}\big)
$$
이다. 『대수공간의 코호몰로지』의 보조정리
\ref{spaces-cohomology-lemma-higher-direct-image}에 의해 오른쪽은
$V_\etale$ 위의 준연접 층이다. 따라서 왼쪽도 준연접이며, 이것이 증명할 내용이다.
\end{proof}

\begin{lemma}
\label{stacks-cohomology-lemma-check-lqc-on-flat-covering}
$\mathcal{X}$를 대수적 스택이라 하자.
$f_j : \mathcal{X}_j \to \mathcal{X}$를
$|\mathcal{X}| =\bigcup |f_j|(|\mathcal{X}_j|)$를 만족하는
대수적 스택의 평탄하고 국소 유한 표시인 사상들의 족이라 하자.
$\mathcal{F}$를 $\mathcal{X}_{fppf}$ 위의
$\mathcal{O}_\mathcal{X}$-가군층이라 하자.
각 $f_j^{-1}\mathcal{F}$가 국소 준연접이면 $\mathcal{F}$도 그러하다.
\end{lemma}

\begin{proof}
먼저 전사이고 평탄하며 국소 유한 표시이고 준콤팩트이며 준분리인 사상
$a : \mathcal{U} \to \mathcal{X}$가 존재하고,
$a^*\mathcal{F}$가 국소 준연접이라고 하자. 그러면 완전열
$$
0 \to \mathcal{F} \to a_*a^*\mathcal{F} \to b_*b^*\mathcal{F}
$$
이 있다. 여기서 $b$는 사상
$b : \mathcal{U} \times_\mathcal{X} \mathcal{U} \to \mathcal{X}$이다.
『스택 위의 층』의 명제
\ref{stacks-sheaves-proposition-exactness-cech-complex}과 보조정리
\ref{stacks-sheaves-lemma-surjective-flat-locally-finite-presentation}를 보라.
또한 당김 $b^*\mathcal{F}$는 사영 사상들 중 하나에 의한
$a^*\mathcal{F}$의 당김이므로 국소 준연접이다
(『스택 위의 층』의 보조정리 \ref{stacks-sheaves-lemma-pullback-lqc}).
가군 $a_*a^*\mathcal{F}$와 $b_*b^*\mathcal{F}$는 보조정리
\ref{stacks-cohomology-lemma-pushforward-locally-quasi-coherent}에 의해 국소 준연접이다.
($a_*$와 $b_*$는 이를 계산하는 데 어느 위상을 사용하는지와 무관하다.)
따라서 $\mathcal{F}$는 국소 준연접이다. 『스택 위의 층』의 보조정리
\ref{stacks-sheaves-lemma-lqc-colimits}를 보라.

\medskip\noindent
이제 일반적인 경우의 증명을 첫 문단의 상황으로 환원하겠다.
$x$를 스킴 $U$ 위에 놓인 $\mathcal{X}$의 대상이라 하자.
$\mathcal{F}|_{U_\etale}$가 준연접 $\mathcal{O}_U$-가군임을 보여야 한다.
이를 $U$ 위에서 (Zariski) 국소적으로 보이면 충분하므로 $U$가 아핀이라고
가정해도 된다. 『스택의 사상』의 보조정리
\ref{stacks-morphisms-lemma-surjective-family-flat-locally-finite-presentation}에 의해,
각 $x \circ a_i$가 어떤 $f_j$를 통해 분해되는 fppf 피복
$\{a_i : U_i \to U\}$가 존재한다. 따라서 $a_i^*\mathcal{F}$는
$(\Sch/U_i)_{fppf}$ 위에서 국소 준연접이다.
피복을 세분하여 $\{U_i \to U\}_{i = 1, \ldots, n}$이 표준 fppf 피복이라고
가정해도 된다. 그러면 $x^*\mathcal{F}$는 $(\Sch/U)_{fppf}$ 위의
fppf 가군이고, 사상
$a : U_1 \amalg \ldots \amalg U_n \to U$에 의한 당김은 국소 준연접이다.
따라서 첫 문단에 의해 $x^*\mathcal{F}$는 국소 준연접이고,
이는 분명 $\mathcal{F}|_{U_\etale}$가 준연접임을 함의한다.
\end{proof}

\section{평탄 비교 사상}
\label{stacks-cohomology-section-flat-comparison}

\noindent
$\mathcal{X}$를 대수적 스택이라 하고 $\mathcal{F}$를
$\textit{Mod}(\mathcal{X}_\etale, \mathcal{O}_\mathcal{X})$의 대상이라 하자.
스킴 $U$ 위에 놓인 $\mathcal{X}$의 대상 $x$가 주어지면,
제한 $\mathcal{F}|_{U_\etale}$는 $x^{-1}\mathcal{F}$를 $U$의 작은 에탈
사이트에 제한한 것이다. 『스택 위의 층』의 정의
\ref{stacks-sheaves-definition-pullback}을 보라.
다음으로 $\varphi : x \to x'$를 스킴의 사상 $f : U \to U'$ 위에 놓인
$\mathcal{X}$의 사상이라 하자. 따라서 다음 $2$-가환 도식
$$
\xymatrix{
U \ar[rd]_x \ar[rr]_f & & U' \ar[ld]^{x'} \\
& \mathcal{X}
}
$$
$\varphi$에 대응하여 제한들 사이의 비교 사상
\begin{equation}
\label{stacks-cohomology-equation-comparison-modules}
c_\varphi :
f_{small}^*(\mathcal{F}|_{U'_\etale})
\longrightarrow
\mathcal{F}|_{U_\etale}
\end{equation}
을 얻는다. 『스택 위의 층』의 식
(\ref{stacks-sheaves-equation-comparison-modules})을 보라.
이 상황에서 $\mathcal{F}$의 다음 성질을 생각할 수 있다.

\begin{definition}
\label{stacks-cohomology-definition-flat-base-change}
$\mathcal{X}$를 대수적 스택이라 하고 $\mathcal{F}$를
$\textit{Mod}(\mathcal{X}_\etale, \mathcal{O}_\mathcal{X})$ 안에 둔다.
$f$가 평탄할 때마다 $c_\varphi$가 동형인 경우, 그리고 그 경우에만
$\mathcal{F}$가 {\it 평탄 밑변환 성질}\footnote{이는 표준적이지 않은
표기일 수 있다.}을 가진다고 한다.
\end{definition}

\noindent
다음 보조정리는 이 개념의 몇 가지 성질을 제시한다.

\begin{lemma}
\label{stacks-cohomology-lemma-check-flat-comparison-on-etale-covering}
$\mathcal{X}$를 대수적 스택이라 하자. $\mathcal{F}$를
$\mathcal{X}_\etale$ 위의 $\mathcal{O}_\mathcal{X}$-가군이라 하자.
\begin{enumerate}
\item $\mathcal{F}$가 평탄 밑변환 성질을 가지면, 대수적 스택의 임의의 사상
$g : \mathcal{Y} \to \mathcal{X}$에 대해 당김 $g^*\mathcal{F}$도
그 성질을 가진다.
\item 평탄 밑변환 성질을 가진 가군들로 이루어진
$\textit{Mod}(\mathcal{X}_\etale, \mathcal{O}_\mathcal{X})$의 충만한
부분범주는 약한 Serre 부분범주이다.
\item $f_i : \mathcal{X}_i \to \mathcal{X}$를
$|\mathcal{X}| = \bigcup_i |f_i|(|\mathcal{X}_i|)$를 만족하는
대수적 스택의 매끄러운 사상들의 족이라 하자. 각
$f_i^*\mathcal{F}$가 평탄 밑변환 성질을 가지면 $\mathcal{F}$도 그 성질을 가진다.
\item $\mathcal{X}_\etale$ 위에서 평탄 밑변환 성질을 가진
$\mathcal{O}_\mathcal{X}$-가군들의 범주는 쌍대극한을 가지며, 이들은
$\textit{Mod}(\mathcal{X}_\etale, \mathcal{O}_\mathcal{X})$에서의
쌍대극한과 일치한다.
\item $\textit{Mod}(\mathcal{X}_\etale, \mathcal{O}_\mathcal{X})$의
$\mathcal{F}$와 $\mathcal{G}$가 평탄 밑변환 성질을 가지면 텐서곱
$\mathcal{F} \otimes_{\mathcal{O}_\mathcal{X}} \mathcal{G}$도
평탄 밑변환 성질을 가진다.
\item $\textit{Mod}(\mathcal{X}_\etale, \mathcal{O}_\mathcal{X})$의
$\mathcal{F}$와 $\mathcal{G}$가 주어지고 $\mathcal{F}$가 유한 표시이며
$\mathcal{G}$가 평탄 밑변환 성질을 가지면 층
$\SheafHom_{\mathcal{O}_\mathcal{X}}(\mathcal{F}, \mathcal{G})$도
평탄 밑변환 성질을 가진다.
\end{enumerate}
\end{lemma}

\begin{proof}
$g : \mathcal{Y} \to \mathcal{X}$가 (1)과 같다고 하자.
$y$를 스킴 $V$ 위에 놓인 $\mathcal{Y}$의 대상이라 하자.
『스택 위의 층』의 보조정리 \ref{stacks-sheaves-lemma-comparison}에 의해
$(g^*\mathcal{F})|_{V_\etale} = \mathcal{F}|_{V_\etale}$이다.
더욱이 $\mathcal{Y}$ 위의 층 $g^*\mathcal{F}$에 대한 비교 사상은
$\mathcal{X}$ 위의 층 $\mathcal{F}$에 대한 비교 사상의 특수한 경우이다.
『스택 위의 층』의 보조정리 \ref{stacks-sheaves-lemma-comparison}을 보라.
이로써 (1)은 명확하다.

\medskip\noindent
(2)의 증명. 『호몰로지』의 보조정리
\ref{homology-lemma-characterize-weak-serre-subcategory}에 나오는
약한 Serre 부분범주의 특성화를 사용한다.
평탄 밑변환 성질을 가진 층들 사이의 사상의 핵과 여핵도
평탄 밑변환 성질을 가진다. 스킴의 평탄 사상에 대해 $f_{small}^*$가
완전하고, 제한 함자 $(-)|_{U_\etale}$도 완전하므로 이는 명확하다
(우리는 에탈 위상에서 작업하고 있다). 마지막으로
$0 \to \mathcal{F}_1 \to \mathcal{F}_2 \to \mathcal{F}_3 \to 0$가
$\textit{Mod}(\mathcal{X}_\etale, \mathcal{O}_\mathcal{X})$의 짧은 완전열이고
바깥의 두 층이 평탄 밑변환 성질을 가지면 가운데 층도 그 성질을 가진다.
이 역시 $f_{small}^*$와 제한 함자의 완전성(그리고 5-보조정리)에서 따른다.

\medskip\noindent
(3)의 증명.
$f_i : \mathcal{X}_i \to \mathcal{X}$를 함께 전사인 대수적 스택의
매끄러운 사상들의 족이라 하고, 각 $f_i^*\mathcal{F}$가 평탄 밑변환
성질을 가진다고 가정하자. (1), 대수적 스택의 정의, 그리고 매끄러운
사상들의 합성이 매끄럽다는 사실(『스택의 사상』의 보조정리
\ref{stacks-morphisms-lemma-composition-smooth}를 보라)에 의해,
각 $\mathcal{X}_i$가 스킴으로 표현가능하다고 가정해도 된다.
$\varphi : x \to x'$를 스킴의 평탄 사상 $a : U \to U'$ 위에 놓인
$\mathcal{X}$의 사상이라 하자. 『스택 위의 층』의 보조정리
\ref{stacks-sheaves-lemma-surjective-flat-locally-finite-presentation}에 의해
$U'_i \to U' \to \mathcal{X}$가 $\mathcal{X}_i$를 통해 분해되는,
함께 전사인 에탈 사상들의 족 $U'_i \to U'$가 존재한다.
따라서 가환 도식들
$$
\xymatrix{
U_i = U \times_{U'} U_i' \ar[r]_-{a_i} \ar[d] &
U_i' \ar[r]_{x_i'} \ar[d] & \mathcal{X}_i \ar[d]^{f_i} \\
U \ar[r]^a & U' \ar[r]^{x'} & \mathcal{X}
}
$$
을 얻는다. 각 $a_i$는 $a$의 밑변환이므로 스킴의 평탄 사상임에 유의하라.
$a_i$ 위에 놓이고 목표가 $x_i'$인 $\mathcal{X}_i$의 사상을
$\psi_i : x_i \to x'_i$라 쓰자. 가정에 따라 비교 사상
$c_{\psi_i} :
(a_i)_{small}^*\big(f_i^*\mathcal{F}|_{(U'_i)_\etale}\big)
\to f_i^*\mathcal{F}|_{(U_i)_\etale}$는 동형이다.
수직 화살 $U_i' \to U'$와 $U_i \to U$가 에탈이므로,
층 $f_i^*\mathcal{F}|_{(U_i')_\etale}$와
$f_i^*\mathcal{F}|_{(U_i)_\etale}$는 각각
$\mathcal{F}|_{U'_\etale}$와 $\mathcal{F}|_{U_\etale}$의 제한이고,
사상 $c_{\psi_i}$는 $c_\varphi$를 $(U_i)_\etale$에 제한한 것이다.
『스택 위의 층』의 보조정리 \ref{stacks-sheaves-lemma-comparison}을 보라.
$\{U_i \to U\}$가 에탈 피복이므로 비교 사상 $c_\varphi$는 동형이다.
이것이 증명할 내용이다.

\medskip\noindent
(4)의 증명. 도식
$\mathcal{I} \to
\textit{Mod}(\mathcal{X}_\etale, \mathcal{O}_\mathcal{X})$,
$i \mapsto \mathcal{F}_i$가 주어지고, 각 $\mathcal{F}_i$가
평탄 밑변환 성질을 가진다고 하자.
$\varphi : x \to x'$를 스킴의 평탄 사상 $f : U \to U'$ 위에 놓인
$\mathcal{X}$의 사상이라 하자. $\colim_i \mathcal{F}_i$는 준층 쌍대극한의
층화임을 상기하라. 에탈 위상을 사용하고 있으므로
$$
(\colim_i \mathcal{F}_i)|_{U_\etale} =
\colim_i {\mathcal{F}_i}|_{U_\etale}
$$
임이 명확하고, $U'_\etale$에 대한 제한도 마찬가지이다. 따라서
\begin{align*}
f_{small}^*((\colim_i \mathcal{F}_i)|_{U'_\etale})
& =
f_{small}^*(\colim_i {\mathcal{F}_i}|_{U'_\etale}) \\
& =
\colim_i f_{small}^*({\mathcal{F}_i}|_{U'_\etale}) \\
& \xrightarrow{\colim c_\varphi}
\colim_i \mathcal{F}_i|_{U_\etale} \\
& =
(\colim_i \mathcal{F}_i)|_{U_\etale}
\end{align*}
이다. 두 번째 등식에는 $f_{small}^*$가 (왼쪽 수반이므로) 쌍대극한과
가환한다는 사실을 사용했다. 각 $\mathcal{F}_i$가 평탄 밑변환 성질을
가지므로 이 화살은 동형이다. 따라서 쌍대극한도 평탄 밑변환 성질을 가지며
(4)가 성립한다.

\medskip\noindent
(5)는 텐서곱이 당김과 가환하기 때문에 성립한다.
『사이트 위의 가군』의 보조정리
\ref{sites-modules-lemma-tensor-product-pullback}을 보라. 세부사항은 생략한다.

\medskip\noindent
$\mathcal{F}$와 $\mathcal{G}$가 (6)과 같다고 하자.
$\mathcal{F}$는 준연접이므로 『스택 위의 층』의 보조정리
\ref{stacks-sheaves-lemma-quasi-coherent}에 의해 평탄 밑변환 성질을 가진다.
$\varphi : x \to x'$를 스킴의 평탄 사상 $f : U \to U'$ 위에 놓인
$\mathcal{X}$의 사상이라 하자. 에탈 위상을 사용하고 있으므로
$$
\SheafHom_{\mathcal{O}_\mathcal{X}}(\mathcal{F}, \mathcal{G})|_{U_\etale} =
\SheafHom_{\mathcal{O}_U}(\mathcal{F}|_{U_\etale}, \mathcal{G}|_{U_\etale})
$$
이고 $U'_\etale$에 대한 제한도 마찬가지이다(세부사항은 생략한다). 따라서
\begin{align*}
f_{small}^*(
\SheafHom_{\mathcal{O}_\mathcal{X}}(\mathcal{F}, \mathcal{G})|_{U'_\etale})
& =
f_{small}^*(
\SheafHom_{\mathcal{O}_{U'}}(
\mathcal{F}|_{U'_\etale}, \mathcal{G}|_{U'_\etale})) \\
& =
\SheafHom_{\mathcal{O}_{U'}}(
f_{small}^*(\mathcal{F}|_{U'_\etale}),
f_{small}^*(\mathcal{G}|_{U'_\etale})) \\
& \xrightarrow{c_\varphi}
\SheafHom_{\mathcal{O}_U}(\mathcal{F}|_{U_\etale}, \mathcal{G}|_{U_\etale}) \\
& =
\SheafHom_{\mathcal{O}_\mathcal{X}}(\mathcal{F}, \mathcal{G})|_{U_\etale}
\end{align*}
이다. 여기서 두 번째 등식은 『사이트 위의 가군』의 보조정리
\ref{sites-modules-lemma-pullback-internal-hom}이며, 이 보조정리는
$f : U \to U'$가 평탄하여 환 달린 사이트의 사상 $f_{small}$도
평탄하다는 사실을 사용한다. $\mathcal{F}$와 $\mathcal{G}$가 모두
평탄 밑변환 성질을 가지므로 이 화살은 동형이다. 따라서 원하는 대로
우리의 $\SheafHom$도 평탄 밑변환 성질을 가진다.
\end{proof}

\begin{lemma}
\label{stacks-cohomology-lemma-flat-comparison}
$f : \mathcal{X} \to \mathcal{Y}$를 대수적 스택의 준콤팩트 준분리 사상이라 하자.
$\mathcal{F}$를 국소 준연접이고 평탄 밑변환 성질을 가지는
$\textit{Mod}(\mathcal{X}_\etale, \mathcal{O}_\mathcal{X})$의 대상이라 하자.
그러면 각 $R^if_*\mathcal{F}$는(에탈 위상에서 계산한 것으로)
평탄 밑변환 성질을 가진다.
\end{lemma}

\begin{proof}
이를 증명하기 위해 보조정리 \ref{stacks-cohomology-lemma-general-pushforward}를 사용한다.
모든 대수적 스택 $\mathcal{X}$에 대해
$\textit{LQCoh}^{fbc}(\mathcal{O}_\mathcal{X})$를
국소 준연접이고 평탄 밑변환 성질을 가진 층들로 이루어진
$\textit{Mod}(\mathcal{X}_\etale, \mathcal{O}_\mathcal{X})$의
충만한 부분범주라 하자. 보조정리 \ref{stacks-cohomology-lemma-general-pushforward}의 조건
(1)--(4)를 확인하면 이 보조정리가 따른다. 성질 (1), (2), (3)은
『스택 위의 층』의 보조정리
\ref{stacks-sheaves-lemma-pullback-lqc}와
\ref{stacks-sheaves-lemma-lqc-colimits}, 그리고 보조정리
\ref{stacks-cohomology-lemma-check-lqc-on-etale-covering}와
\ref{stacks-cohomology-lemma-check-flat-comparison-on-etale-covering}에서 따른다.
따라서 (4)를 보이면 충분하다.

\medskip\noindent
$f : \mathcal{X} \to \mathcal{Y}$가 대수적 스택의 사상이고
$\mathcal{X}$와 $\mathcal{Y}$가 아핀 스킴 $X$와 $Y$로 표현가능하다고 하자.
이 경우 $\psi : y \to y'$를 스킴의 평탄 사상 $b : V \to V'$ 위에 놓인
$\mathcal{Y}$의 사상이라 하자. 명료성을 위해 대응하는 대수적 스택들을
$\mathcal{V} = (\Sch/V)_{fppf}$와 $\mathcal{V}' = (\Sch/V')_{fppf}$라 쓰자.
두 정사각형이 모두 카르테시안인 대수적 스택의 도식
$$
\xymatrix{
\mathcal{Z} \ar[d]_{f''} \ar[r]_a &
\mathcal{Z}' \ar[r]_{x'} \ar[d]_{f'} & \mathcal{X} \ar[d]^f \\
\mathcal{V} \ar[r]^b & \mathcal{V}' \ar[r]^{y'} & \mathcal{Y}
}
$$
을 생각하자. $f$는 스킴으로 표현가능하므로(그리고 준콤팩트이고 분리이며,
실제로 아핀이므로), $\mathcal{Z}$와 $\mathcal{Z}'$는 스킴 $Z$와 $Z'$로
표현가능하고 실제로 $Z = V \times_{V'} Z'$이다.
$\mathcal{F}$가 평탄 밑변환 성질을 가지므로
$$
a_{small}^*\big(\mathcal{F}|_{Z'_\etale}\big)
\longrightarrow
\mathcal{F}|_{Z_\etale}
$$
은 동형이다. 더욱이 『스택 위의 층』의 보조정리
\ref{stacks-sheaves-lemma-compare-representable-morphism-cohomology}에 의해
$$
R^if_*\mathcal{F}|_{V'_\etale} =
R^i(f')_{small, *}\big(\mathcal{F}|_{Z'_\etale}\big)
$$
이고
$$
R^if_*\mathcal{F}|_{V_\etale} =
R^i(f'')_{small, *}\big(\mathcal{F}|_{Z_\etale}\big)
$$
이다. 따라서 비교 사상
$$
c_\psi :
b_{small}^*(R^if_*\mathcal{F}|_{V'_\etale})
\longrightarrow
R^if_*\mathcal{F}|_{V_\etale}
$$
은 『대수공간의 코호몰로지』의 보조정리
\ref{spaces-cohomology-lemma-flat-base-change-cohomology}에 의해 동형이다.
그러므로 $R^if_*\mathcal{F}$는 평탄 밑변환 성질을 가진다.
보조정리 \ref{stacks-cohomology-lemma-pushforward-locally-quasi-coherent}에 의해
$R^if_*\mathcal{F}$는 국소 준연접이기도 하므로 결론을 얻는다.
\end{proof}

\section{평탄 밑변환 성질을 가진 국소 준연접 가군}
\label{stacks-cohomology-section-loc-qcoh-flat-base-change}

\noindent
$\mathcal{X}$를 대수적 스택이라 하자. 우리는\footnote{끔찍한 표기에 대해
양해를 구한다.}
$$
\textit{LQCoh}^{fbc}(\mathcal{O}_\mathcal{X})
\subset
\textit{Mod}(\mathcal{X}_\etale, \mathcal{O}_\mathcal{X})
$$
로, 에탈 $\mathcal{O}_\mathcal{X}$-가군 $\mathcal{F}$ 가운데
국소 준연접이고(절 \ref{stacks-cohomology-section-locally-quasi-coherent}) 동시에
평탄 밑변환 성질을 가지는(절 \ref{stacks-cohomology-section-flat-comparison}) 것들을
대상으로 하는 충만한 부분범주를 나타내겠다.
『스택 위의 층』의 보조정리 \ref{stacks-sheaves-lemma-quasi-coherent}에 의해
$$
\QCoh(\mathcal{O}_\mathcal{X}) \subset
\textit{LQCoh}^{fbc}(\mathcal{O}_\mathcal{X})
$$
이다.

\begin{proposition}
\label{stacks-cohomology-proposition-loc-qcoh-flat-base-change}
평탄 밑변환 성질을 가진 국소 준연접 가군에 관한 결과들을 요약하면 다음과 같다.
\begin{enumerate}
\item $\mathcal{X}$를 대수적 스택이라 하자.
$\mathcal{F}$가 $\textit{LQCoh}^{fbc}(\mathcal{O}_\mathcal{X})$에 속하면,
$\mathcal{F}$는 fppf 위상에 대한 층이다. 즉,
$\textit{Mod}(\mathcal{O}_\mathcal{X})$의 대상이다.
\item 범주 $\textit{LQCoh}^{fbc}(\mathcal{O}_\mathcal{X})$는
$\textit{Mod}(\mathcal{O}_\mathcal{X})$와
$\textit{Mod}(\mathcal{X}_\etale, \mathcal{O}_\mathcal{X})$ 양쪽의
약한 Serre 부분범주이다.
\item 대수적 스택의 임의의 사상 $f : \mathcal{X} \to \mathcal{Y}$에
따른 당김 $f^*$는 함자
$f^* : \textit{LQCoh}^{fbc}(\mathcal{O}_\mathcal{Y}) \to
\textit{LQCoh}^{fbc}(\mathcal{O}_\mathcal{X})$를 유도한다.
\item $f : \mathcal{X} \to \mathcal{Y}$가 대수적 스택의
준콤팩트 준분리 사상이고 $\mathcal{F}$가
$\textit{LQCoh}^{fbc}(\mathcal{O}_\mathcal{X})$의 대상이면 다음이 성립한다.
\begin{enumerate}
\item 전체 직상 $Rf_*\mathcal{F}$와 고차 직상
$R^if_*\mathcal{F}$는 에탈 위상과 fppf 위상 어느 쪽에서도 계산할 수 있고
그 결과는 같다.
\item 각 $R^if_*\mathcal{F}$는
$\textit{LQCoh}^{fbc}(\mathcal{O}_\mathcal{Y})$의 대상이다.
\end{enumerate}
\item 범주 $\textit{LQCoh}^{fbc}(\mathcal{O}_\mathcal{X})$는 쌍대극한을
가지며, 이들은 $\textit{Mod}(\mathcal{X}_\etale, \mathcal{O}_\mathcal{X})$와
$\textit{Mod}(\mathcal{O}_\mathcal{X})$ 양쪽에서의 쌍대극한과 일치한다.
\item $\textit{LQCoh}^{fbc}(\mathcal{O}_\mathcal{X})$의
$\mathcal{F}$와 $\mathcal{G}$가 주어지면 텐서곱
$\mathcal{F} \otimes_{\mathcal{O}_\mathcal{X}} \mathcal{G}$는
$\textit{LQCoh}^{fbc}(\mathcal{O}_\mathcal{X})$에 속한다.
\item $\mathcal{F}$가 유한 표시이고 $\mathcal{G}$가
$\textit{LQCoh}^{fbc}(\mathcal{O}_\mathcal{X})$에 속하면
$\SheafHom_{\mathcal{O}_\mathcal{X}}(\mathcal{F}, \mathcal{G})$는
$\textit{LQCoh}^{fbc}(\mathcal{O}_\mathcal{X})$에 속한다.
\end{enumerate}
\end{proposition}

\begin{proof}
(1)은 『스택 위의 층』의 보조정리
\ref{stacks-sheaves-lemma-lqc-flat-base-change-fppf-sheaf}이다.

\medskip\noindent
매장
$\textit{LQCoh}^{fbc}(\mathcal{O}_\mathcal{X}) \subset
\textit{Mod}(\mathcal{X}_\etale, \mathcal{O}_\mathcal{X})$에 대한 (2)는
보조정리 \ref{stacks-cohomology-lemma-flat-comparison}의 증명에서 보았다.
이제 매장
$\textit{LQCoh}^{fbc}(\mathcal{O}_\mathcal{X}) \subset
\textit{Mod}(\mathcal{O}_\mathcal{X})$에 대해 (2)를 증명하자.
$\varphi : \mathcal{F} \to \mathcal{G}$를
$\textit{LQCoh}^{fbc}(\mathcal{O}_\mathcal{X})$의 대상들 사이의 사상이라 하자.
$\Ker(\varphi)$는 에탈 위상에서 계산하든 fppf 위상에서 계산하든 같으므로,
에탈 경우에 의해 $\Ker(\varphi)$는
$\textit{LQCoh}^{fbc}(\mathcal{O}_\mathcal{X})$에 속한다.
한편 fppf 위상에서 계산한 여핵은 에탈 위상에서 계산한 여핵의 fppf 층화이다.
그러나 이 에탈 여핵은
$\textit{LQCoh}^{fbc}(\mathcal{O}_\mathcal{X})$에 속하므로 (1)에 의해
fppf 층이고, 따라서 여핵은
$\textit{LQCoh}^{fbc}(\mathcal{O}_\mathcal{X})$에 속한다.
마지막으로
$$
0 \to \mathcal{F}_1 \to \mathcal{F}_2 \to \mathcal{F}_3 \to 0
$$
이 $\textit{Mod}(\mathcal{O}_\mathcal{X})$에서의 완전열이고
(즉, fppf 위상을 사용하여), $\mathcal{F}_1$, $\mathcal{F}_2$가
$\textit{LQCoh}^{fbc}(\mathcal{O}_\mathcal{X})$에 속한다고 하자.
$\mathcal{F}_2$가 $\textit{LQCoh}^{fbc}(\mathcal{O}_\mathcal{X})$의 대상임을
보이려면 이 열이 에탈 위상에서도 완전함을 보이면 충분하다.
이를 위해서는 (임의의 $\mathcal{X}$의 대상 $x$에 대해)
$H^1_{fppf}(x, \mathcal{F}_1)$의 임의의 원소가 $x$의 어떤 에탈 피복의
각 원소 위에서 영이 됨을 보이면 충분하다. 『스택 위의 층』의 보조정리
\ref{stacks-sheaves-lemma-compare-fppf-etale}에 의해
$H^1_{fppf}(x, \mathcal{F}_1) = H^1_\etale(x, \mathcal{F}_1)$이고,
코호몰로지의 국소성 때문에 이는 참이다. 『사이트 위의 코호몰로지』의 보조정리
\ref{sites-cohomology-lemma-kill-cohomology-class-on-covering}을 보라.
이로써 (2)가 증명되었다.

\medskip\noindent
(3)은 보조정리 \ref{stacks-cohomology-lemma-check-flat-comparison-on-etale-covering}와
『스택 위의 층』의 보조정리 \ref{stacks-sheaves-lemma-pullback-lqc}에서 따른다.

\medskip\noindent
에탈 코호몰로지에서 계산한 $R^if_*\mathcal{F}$에 대한 (4)(b)는
보조정리 \ref{stacks-cohomology-lemma-flat-comparison}에서 따른다.
그러면 (4)(a)는 『스택 위의 층』의 보조정리
\ref{stacks-sheaves-lemma-compare-fppf-etale}와 위의 (1)을 결합하면 따른다.

\medskip\noindent
에탈 위상에 대한 (5)는 『스택 위의 층』의 보조정리
\ref{stacks-sheaves-lemma-lqc-colimits}와 보조정리
\ref{stacks-cohomology-lemma-check-flat-comparison-on-etale-covering}에서 따른다.
그러면 에탈 위상에서의 쌍대극한이 (1)에 의해 이미 fppf 층이므로
fppf 판본도 따른다.

\medskip\noindent
(6)과 (7)은 보조정리 \ref{stacks-cohomology-lemma-check-flat-comparison-on-etale-covering}의
대응하는 부분과 『스택 위의 층』의 보조정리
\ref{stacks-sheaves-lemma-lqc-colimits}에서 따른다.
\end{proof}

\begin{lemma}
\label{stacks-cohomology-lemma-check-lqc-fbc-on-covering}
$\mathcal{X}$를 대수적 스택이라 하자.
\begin{enumerate}
\item $f_j : \mathcal{X}_j \to \mathcal{X}$를
$|\mathcal{X}| =\bigcup |f_j|(|\mathcal{X}_j|)$를 만족하는
대수적 스택의 매끄러운 사상들의 족이라 하자.
$\mathcal{F}$를 $\mathcal{X}_\etale$ 위의
$\mathcal{O}_\mathcal{X}$-가군층이라 하자. 각 $f_j^{-1}\mathcal{F}$가
$\textit{LQCoh}^{fpc}(\mathcal{O}_{\mathcal{X}_i})$에 속하면,
$\mathcal{F}$는 $\textit{LQCoh}^{fbc}(\mathcal{O}_\mathcal{X})$에 속한다.
\item $f_j : \mathcal{X}_j \to \mathcal{X}$를
$|\mathcal{X}| =\bigcup |f_j|(|\mathcal{X}_j|)$를 만족하는
대수적 스택의 평탄하고 국소 유한 표시인 사상들의 족이라 하자.
$\mathcal{F}$를 $\mathcal{X}_{fppf}$ 위의
$\mathcal{O}_\mathcal{X}$-가군층이라 하자. 각 $f_j^{-1}\mathcal{F}$가
$\textit{LQCoh}^{fbc}(\mathcal{O}_{\mathcal{X}_i})$에 속하면,
$\mathcal{F}$는 $\textit{LQCoh}^{fbc}(\mathcal{O}_\mathcal{X})$에 속한다.
\end{enumerate}
\end{lemma}

\begin{proof}
(1)은 보조정리 \ref{stacks-cohomology-lemma-check-lqc-on-etale-covering}와
\ref{stacks-cohomology-lemma-check-flat-comparison-on-etale-covering}를 결합하면 따른다.
(2)의 증명은 보조정리 \ref{stacks-cohomology-lemma-check-lqc-on-flat-covering}의 증명과 유사하다.
$\mathcal{F}$를 $\mathcal{X}_{fppf}$ 위의
$\mathcal{O}_\mathcal{X}$-가군층인 것으로 하자.

\medskip\noindent
먼저 전사이고 평탄하며 국소 유한 표시이고 준콤팩트이며 준분리인 사상
$a : \mathcal{U} \to \mathcal{X}$가 존재하고,
$a^*\mathcal{F}$가 국소 준연접이며 평탄 밑변환 성질을 가진다고 하자.
그러면 완전열
$$
0 \to \mathcal{F} \to a_*a^*\mathcal{F} \to b_*b^*\mathcal{F}
$$
이 있다. 여기서 $b$는 사상
$b : \mathcal{U} \times_\mathcal{X} \mathcal{U} \to \mathcal{X}$이다.
『스택 위의 층』의 명제
\ref{stacks-sheaves-proposition-exactness-cech-complex}과 보조정리
\ref{stacks-sheaves-lemma-surjective-flat-locally-finite-presentation}를 보라.
더욱이 당김 $b^*\mathcal{F}$는 사영 사상들 중 하나에 의한
$a^*\mathcal{F}$의 당김이므로 국소 준연접이고 평탄 밑변환 성질을 가진다.
명제 \ref{stacks-cohomology-proposition-loc-qcoh-flat-base-change}를 보라.
가군 $a_*a^*\mathcal{F}$와 $b_*b^*\mathcal{F}$는 명제
\ref{stacks-cohomology-proposition-loc-qcoh-flat-base-change}에 의해 국소 준연접이고
평탄 밑변환 성질을 가진다. 명제
\ref{stacks-cohomology-proposition-loc-qcoh-flat-base-change}에 의해 $\mathcal{F}$도
국소 준연접이고 평탄 밑변환 성질을 가진다는 결론을 얻는다.

\medskip\noindent
스킴 $U$와 매끄러운 전사 사상 $x : U \to \mathcal{X}$를 택한다.
(1)에 의해 $x^*\mathcal{F}$가 국소 준연접이고 평탄 밑변환 성질을 가짐을
보이면 충분하다. 다시 (1)에 의해 이를 $U$ 위에서 (Zariski) 국소적으로
보이면 충분하므로 $U$가 아핀이라고 가정해도 된다. 『스택의 사상』의 보조정리
\ref{stacks-morphisms-lemma-surjective-family-flat-locally-finite-presentation}에 의해,
각 $x \circ a_i$가 어떤 $f_j$를 통해 분해되는 fppf 피복
$\{a_i : U_i \to U\}$가 존재한다. 따라서 $(\Sch/U_i)_{fppf}$ 위의 가군
$a_i^*\mathcal{F}$는 국소 준연접이고 평탄 밑변환 성질을 가진다.
피복을 세분하여 $\{U_i \to U\}_{i = 1, \ldots, n}$이 표준 fppf 피복이라고
가정해도 된다. 그러면 $x^*\mathcal{F}$는 $(\Sch/U)_{fppf}$ 위의
fppf 가군이고, 사상 $a : U_1 \amalg \ldots \amalg U_n \to U$에 의한
당김은 국소 준연접이고 평탄 밑변환 성질을 가진다.
따라서 앞 문단에 의해 원하는 대로 $x^*\mathcal{F}$가 국소 준연접이고
평탄 밑변환 성질을 가짐을 알 수 있다.
\end{proof}

\begin{lemma}
\label{stacks-cohomology-lemma-loc-qcoh-fbc-compute-higher-direct-image}
$f : \mathcal{X} \to \mathcal{Y}$를 준콤팩트이고 준분리이며
대수공간으로 표현가능한 대수적 스택의 사상이라 하자.
$\mathcal{F}$가 $\textit{LQCoh}^{fbc}(\mathcal{O}_\mathcal{X})$에 속한다고 하자.
$\mathcal{Y}$의 대상 $y : V \to \mathcal{Y}$에 대해
$$
(R^if_*\mathcal{F})|_{V_\etale} = R^if'_{small, *}(\mathcal{F}|_{U_\etale})
$$
이다. 여기서 $f' : U = V \times_\mathcal{Y} \mathcal{X} \to V$는
$f$의 밑변환이다.
\end{lemma}

\begin{proof}
『스택 위의 층』의 보조정리
\ref{stacks-sheaves-lemma-base-change-higher-direct-images}에 의해
$\mathcal{X}$가 $U$로 표현되고 $\mathcal{Y}$가 $V$로 표현되는 경우로
환원할 수 있다. 물론 이는 명제 \ref{stacks-cohomology-proposition-loc-qcoh-flat-base-change}에 의해
$\mathcal{F}$를 $U$로 당긴 것이
$\textit{LQCoh}^{fbc}(\mathcal{O}_U)$에 속한다는 사실도 사용한다.
이제 결과는 『스택 위의 층』의 보조정리
\ref{stacks-sheaves-lemma-compare-morphism-cohomology}와, 명제
\ref{stacks-cohomology-proposition-loc-qcoh-flat-base-change}에 의해 $R^if_*$를 에탈 위상에서
계산할 수 있다는 사실에서 따른다.
\end{proof}

\begin{lemma}
\label{stacks-cohomology-lemma-loc-qcoh-fbc-affine-direct-image}
$f : \mathcal{X} \to \mathcal{Y}$를 대수적 스택의 아핀 사상이라 하자.
함자 $f_* : \textit{LQCoh}^{fbc}(\mathcal{O}_\mathcal{X}) \to
\textit{LQCoh}^{fbc}(\mathcal{O}_\mathcal{Y})$는 완전하고 직합과 가환한다.
$i > 0$에 대한 함자 $R^if_*$는
$\textit{LQCoh}^{fbc}(\mathcal{O}_\mathcal{X})$에서 영이다.
\end{lemma}

\begin{proof}
함자들은 명제 \ref{stacks-cohomology-proposition-loc-qcoh-flat-base-change}에 의해 존재한다.
보조정리 \ref{stacks-cohomology-lemma-loc-qcoh-fbc-compute-higher-direct-image}에 의해 이는
대수공간 위의 준연접 가군의 맥락에서 대수공간의 아핀 사상의 고차 직상을
취하는 경우로 환원된다. 『대수공간의 코호몰로지』의 절
\ref{spaces-cohomology-section-higher-direct-image}의 논의에 의해
스킴의 아핀 사상의 경우로 환원된다. 스킴의 아핀 사상에 대해서는
『스킴의 코호몰로지』의 보조정리
\ref{coherent-lemma-relative-affine-vanishing}에 의해 준연접 가군의 고차
직상이 소멸한다. $R^1f_*$의 소멸은 $f_*$의 완전성을 함의한다.
직합과의 가환은 예를 들어 『사상』의 보조정리
\ref{morphisms-lemma-affine-equivalence-modules}에서 따른다.
\end{proof}

\section{기생 가군}
\label{stacks-cohomology-section-parasitic}

\noindent
다음 정의는 『하강』의 정의 \ref{descent-definition-parasitic}와 양립한다.

\begin{definition}
\label{stacks-cohomology-definition-parasitic}
$\mathcal{X}$를 대수적 스택이라 하자.
$\mathcal{O}_\mathcal{X}$-가군의 준층 $\mathcal{F}$가 다음 조건을 만족하면
{\it 기생}이라고 한다. $\mathcal{X}$의 대상 $x$가 스킴 $U$ 위에 놓이고
이에 대응하는 사상 $x : U \to \mathcal{X}$가 평탄하면
$\mathcal{F}(x) = 0$이다.
\end{definition}

\noindent
다음 보조정리는 이 개념의 몇 가지 성질을 제시한다.

\begin{lemma}
\label{stacks-cohomology-lemma-parasitic}
$\mathcal{X}$를 대수적 스택이라 하자. $\mathcal{F}$를
$\mathcal{O}_\mathcal{X}$-가군의 준층이라 하자.
\begin{enumerate}
\item $\mathcal{F}$가 기생이고
$g : \mathcal{Y} \to \mathcal{X}$가 대수적 스택의 평탄 사상이면,
$g^*\mathcal{F}$는 기생이다.
\item $\tau \in \{Zariski, \etale, smooth, syntomic, fppf\}$에 대해 다음이 성립한다.
\begin{enumerate}
\item 기생 가군 준층의 $\tau$ 층화는 기생이다.
\item 기생 가군들로 이루어진
$\textit{Mod}(\mathcal{X}_\tau, \mathcal{O}_\mathcal{X})$의 충만한
부분범주는 Serre 부분범주이다.
\end{enumerate}
\item $\mathcal{F}$가 에탈 위상에 대한 층이라고 하자.
$f_i : \mathcal{X}_i \to \mathcal{X}$를
$|\mathcal{X}| = \bigcup_i |f_i|(|\mathcal{X}_i|)$를 만족하는
대수적 스택의 매끄러운 사상들의 족이라 하자.
각 $f_i^*\mathcal{F}$가 기생이면 $\mathcal{F}$도 기생이다.
\item $\mathcal{F}$가 fppf 위상에 대한 층이라고 하자.
$f_i : \mathcal{X}_i \to \mathcal{X}$를
$|\mathcal{X}| = \bigcup_i |f_i|(|\mathcal{X}_i|)$를 만족하는
대수적 스택의 평탄하고 국소 유한 표시인 사상들의 족이라 하자.
각 $f_i^*\mathcal{F}$가 기생이면 $\mathcal{F}$도 기생이다.
\end{enumerate}
\end{lemma}

\begin{proof}
(1)을 보기 위해 $y$를 스킴 $V$ 위에 놓인 $\mathcal{Y}$의 대상이라 하고,
이에 대응하는 사상 $y : V \to \mathcal{Y}$가 평탄하다고 하자.
그러면 $g(y) : V \to \mathcal{Y} \to \mathcal{X}$는 평탄 사상들의 합성이므로
평탄하고(『스택의 사상』의 보조정리
\ref{stacks-morphisms-lemma-composition-flat}을 보라), 가정에 의해
$\mathcal{F}(g(y))$는 영이다. 그런데
$g^*\mathcal{F} = g^{-1}\mathcal{F}(y) = \mathcal{F}(g(y))$이므로
$g^*\mathcal{F}$가 기생이라는 결론을 얻는다.

\medskip\noindent
(2)(a)를 보기 위해 $\{x_i \to x\}$가 $\mathcal{X}$의 $\tau$-피복이면
각 사상 $x_i \to x$가 스킴의 평탄 사상 위에 놓임에 유의하라.
따라서 $x$가 스킴 $U$ 위에 놓이고 $x : U \to \mathcal{X}$가 평탄하면,
모든 대상 $x_i$도 그러하다. 그러므로 준층 $\mathcal{F}$가 기생이면
준층 $\mathcal{F}^+$도 기생이다(『사이트』의 절
\ref{sites-section-sheafification}을 보라). $\mathcal{F}$의 층화가
$(\mathcal{F}^+)^+$이므로 이는 (2)(a)를 증명한다.

\medskip\noindent
$\mathcal{F}$를 기생 $\tau$-가군이라 하자. 정의에서 즉시
$\mathcal{F}$의 모든 부분가군이 기생임을 알 수 있다. 한편
$\mathcal{F}' \subset \mathcal{F}$가 부분가군이면 준층
$x \mapsto \mathcal{F}(x)/\mathcal{F}'(x)$가 기생임도 똑같이 명확하다.
따라서 (2)(a)에 의해 몫 $\mathcal{F}/\mathcal{F}'$은 기생 가군이다.
마지막으로 $\mathcal{F}_1$과 $\mathcal{F}_3$가 기생인 짧은 완전열
$0 \to \mathcal{F}_1 \to \mathcal{F}_2 \to \mathcal{F}_3 \to 0$가
주어지면 $\mathcal{F}_2$가 기생임을 보여야 한다.
$\mathcal{X}$ 위에서 평탄한 스킴 위에 놓인 $x$에서 값을 취하면 이는 즉시 따른다.
이로써 (2)(b)가 증명되었다. 『호몰로지』의 보조정리
\ref{homology-lemma-characterize-serre-subcategory}를 보라.

\medskip\noindent
$f_i : \mathcal{X}_i \to \mathcal{X}$를 함께 전사인 대수적 스택의
매끄러운 사상들의 족이라 하고, 각 $f_i^*\mathcal{F}$가 기생이라고 가정하자.
$x$를 스킴 $U$ 위에 놓인 $\mathcal{X}$의 대상이라 하고
$x : U \to \mathcal{X}$가 평탄하다고 하자.
매끄러운 전사 피복 $W_i \to U \times_{x, \mathcal{X}} \mathcal{X}_i$를
생각하자. 사영을 $y_i : W_i \to \mathcal{X}_i$라 쓰자.
그러면 $\{f_i(y_i) \to x\}$는 $\mathcal{X}$의 매끄러운 위상에 대한
피복이다. 평탄 사상들의 합성은 평탄하므로
$f_i^*\mathcal{F}(y_i) = 0$임을 알 수 있다. 한편 (1)의 증명에서 보았듯이
$f_i^*\mathcal{F}(y_i) = \mathcal{F}(f_i(y_i))$이다.
따라서 $\mathcal{X}$ 안의 어떤 매끄러운 피복
$\{x_i \to x\}_{i \in I}$에 대해 $\mathcal{F}(x_i) = 0$이다.
매끄러운 위상은 에탈 위상과 같으므로 이는 $\mathcal{F}(x) = 0$을 함의한다.
『사상 심화』의 보조정리
\ref{more-morphisms-lemma-etale-dominates-smooth}를 보라.
더 구체적으로, $\{x_i \to x\}_{i \in I}$는 스킴들의 매끄러운 피복
$\{U_i \to U\}_{i \in I}$ 위에 놓인다. 방금 인용한 보조정리에 의해
$\{U_i \to U\}_{i \in I}$를 세분하는 에탈 피복
$\{V_j \to U\}_{j \in J}$가 존재한다. $x'_j = x|_{V_j}$라 쓰자.
그러면 $\{x'_j \to x\}$는 $\{x_i \to x\}_{i \in I}$를 세분하는
$\mathcal{X}$ 안의 에탈 피복이다. 이는 사상
$\mathcal{F}(x) \to \prod_{j \in J} \mathcal{F}(x'_j)$가
$\mathcal{F}(x) \to \prod_{i \in I} \mathcal{F}(x_i)$를 통해
분해된다는 뜻이다. 전자는 $\mathcal{F}$가 에탈 위상의 층이므로 단사이고,
후자는 영이다. 따라서 원하는 대로 $\mathcal{F}(x) = 0$이다.

\medskip\noindent
(4)의 증명은 생략한다. 힌트: (3)의 증명과 유사하지만 더 간단하다.
\end{proof}

\noindent
기생 가군은 정말로 어떠한 직상에 의해서도 보존된다.

\begin{lemma}
\label{stacks-cohomology-lemma-pushforward-parasitic}
$\tau \in \{\etale, fppf\}$라 하자.
$\mathcal{X}$를 대수적 스택이라 하자. $\mathcal{F}$를
$\textit{Mod}(\mathcal{X}_\tau, \mathcal{O}_\mathcal{X})$의 기생 대상이라 하자.
\begin{enumerate}
\item 모든 $i$에 대해 $H^i_\tau(\mathcal{X}, \mathcal{F}) = 0$이다.
\item $f : \mathcal{X} \to \mathcal{Y}$를 대수적 스택의 사상이라 하자.
그러면 $R^if_*\mathcal{F}$는($\tau$-위상에서 계산한 것으로)
$\textit{Mod}(\mathcal{Y}_\tau, \mathcal{O}_\mathcal{Y})$의 기생 대상이다.
\end{enumerate}
\end{lemma}

\begin{proof}
먼저 (2)를 (1)로 환원한다.
『스택 위의 층』의 보조정리 \ref{stacks-sheaves-lemma-pushforward-restriction}에 의해
$R^if_*\mathcal{F}$는 준층
$$
y \longmapsto
H^i_\tau\Big(V \times_{y, \mathcal{Y}} \mathcal{X},
\ \text{pr}^{-1}\mathcal{F}\Big)
$$
에 연관된 층이다. 여기서 $y$는 스킴 $V$ 위에 놓인 $\mathcal{Y}$의
일반적인 대상이다. 보조정리 \ref{stacks-cohomology-lemma-parasitic}에 의해,
$y : V \to \mathcal{Y}$가 평탄할 때 이 코호몰로지 군들이 영임을 보이면 충분하다.
$\text{pr} : V \times_{y, \mathcal{Y}} \mathcal{X} \to \mathcal{X}$는
$y$의 밑변환이므로 평탄하다. 따라서 보조정리 \ref{stacks-cohomology-lemma-parasitic}에 의해
$\text{pr}^{-1}\mathcal{F}$가 기생임을 알 수 있다. 그러므로 (1)을 증명하면 충분하다.

\medskip\noindent
(1)을 보기 위해 『스택 위의 층』의 명제
\ref{stacks-sheaves-proposition-smooth-covering-compute-cohomology}의
스펙트럼 열을 사용하여 $\mathcal{X}$가 대수공간으로 표현가능한 대수적 스택인
경우로 환원할 수 있다. 스펙트럼 열에서 각
$f_p^{-1}\mathcal{F} = f_p^*\mathcal{F}$는 보조정리
\ref{stacks-cohomology-lemma-parasitic}에 의해 기생 가군임에 유의하라. 실제로 사상
$f_p : \mathcal{U}_p =
\mathcal{U} \times_\mathcal{X} \ldots
\times_\mathcal{X} \mathcal{U} \to \mathcal{X}$가 평탄하다.
이 스펙트럼 열을 한 번 더 사용하면(보조정리
\ref{stacks-cohomology-lemma-general-pushforward}의 증명에서처럼) 대수적 스택
$\mathcal{X}$가 스킴 $X$로 표현가능한 경우로 환원된다.
그러면 $H^i_\tau(\mathcal{X}, \mathcal{F}) = H^i((\Sch/X)_\tau, \mathcal{F})$이다.
이 경우 소멸은 {\v C}ech 피복에 관한 논증에서 쉽게 따른다.
『하강』의 보조정리 \ref{descent-lemma-cohomology-parasitic}를 보라.
\end{proof}

\noindent
다음 보조정리는 우리가 기생 가군에 관심을 갖는 주된 이유 가운데 하나이다.
명제를 이해하려면 함자
$\QCoh(\mathcal{O}_\mathcal{X}) \to
\textit{Mod}(\mathcal{X}_\etale, \mathcal{O}_\mathcal{X})$와
$\QCoh(\mathcal{O}_\mathcal{X}) \to
\textit{Mod}(\mathcal{O}_\mathcal{X})$가 일반적으로 완전하지 않음을 상기하라.

\begin{lemma}
\label{stacks-cohomology-lemma-exact-sequence-quasi-coherent-parasitic-cohomology}
$\mathcal{X}$를 대수적 스택이라 하자.
$\alpha : \mathcal{F} \to \mathcal{G}$와
$\beta : \mathcal{G} \to \mathcal{H}$를
$\beta \circ \alpha = 0$을 만족하는
$\QCoh(\mathcal{O}_\mathcal{X})$의 사상이라 하자. 다음은 동치이다.
\begin{enumerate}
\item 아벨 범주 $\QCoh(\mathcal{O}_\mathcal{X})$에서 복합체
$\mathcal{F} \to \mathcal{G} \to \mathcal{H}$는
$\mathcal{G}$에서 완전하다.
\item $\textit{Mod}(\mathcal{X}_\etale, \mathcal{O}_\mathcal{X})$ 또는
$\textit{Mod}(\mathcal{X}_{fppf}, \mathcal{O}_\mathcal{X})$ 어느 쪽에서
계산한 $\Ker(\beta)/\Im(\alpha)$도 기생이다.
\end{enumerate}
\end{lemma}

\begin{proof}
절 \ref{stacks-cohomology-section-loc-qcoh-flat-base-change}에 의해
$\QCoh(\mathcal{O}_\mathcal{X}) \subset
\textit{LQCoh}^{fbc}(\mathcal{O}_\mathcal{X})$이다.
따라서 명제 \ref{stacks-cohomology-proposition-loc-qcoh-flat-base-change}에 의해
$\textit{Mod}(\mathcal{X}_\etale, \mathcal{O}_\mathcal{X})$에서 계산한
$\Ker(\beta)/\Im(\alpha)$와
$\textit{Mod}(\mathcal{X}_{fppf}, \mathcal{O}_\mathcal{X})$에서 계산한 것은
일치한다. 이제부터 $\mathcal{X}$ 위의 에탈 위상을 사용하겠다.

\medskip\noindent
$\mathcal{E}$를 아벨 범주 $\QCoh(\mathcal{O}_\mathcal{X})$에서 계산한
$\mathcal{F} \to \mathcal{G} \to \mathcal{H}$의 코호몰로지라 하자.
$x : U \to \mathcal{X}$를 평탄 사상이라 하고 $U$를 스킴이라 하자.
에탈 위상을 사용하고 있으므로 제한 함자
$\textit{Mod}(\mathcal{X}_\etale, \mathcal{O}_\mathcal{X})
\to \textit{Mod}(U_\etale, \mathcal{O}_U)$는 완전하다.
한편 보조정리 \ref{stacks-cohomology-lemma-flat-pullback-quasi-coherent}와
『스택 위의 층』의 보조정리
\ref{stacks-sheaves-lemma-compare-quasi-coherent}에 의해 제한 함자
$$
\QCoh(\mathcal{O}_\mathcal{X}) \xrightarrow{x^*}
\QCoh((\Sch/U)_\etale, \mathcal{O}) \xrightarrow{{-}|_{U_\etale}}
\QCoh(U_\etale, \mathcal{O}_U)
$$
도 완전하다. 따라서
$\mathcal{E}|_{U_\etale} = (\Ker(\beta)/\Im(\alpha))|_{U_\etale}$이다.

\medskip\noindent
(1)이 성립하면 $\mathcal{E} = 0$이므로
$\mathcal{X}$ 위에서 평탄한 모든 $U$에 대해
$\Ker(\beta)/\Im(\alpha)$를 $U_\etale$에 제한한 것은 영이다.
이것이 기생 가군의 정의이다. (2)가 성립하면
$\mathcal{X}$ 위에서 평탄한 모든 $U$에 대해
$\Ker(\beta)/\Im(\alpha)$를 $U_\etale$에 제한한 것이 영이므로,
$\mathcal{X}$ 위에서 평탄한 모든 $U$에 대해
$\mathcal{E}$를 $U_\etale$에 제한한 것도 영이다.
이는 분명 준연접 가군 $\mathcal{E}$가 영임을 함의한다. 예를 들어
사상 $0 \to \mathcal{E}$에 보조정리
\ref{stacks-cohomology-lemma-quasi-coherent-check-exact}를 적용하라.
\end{proof}

\section{준연접 가군}
\label{stacks-cohomology-section-quasi-coherent}

\noindent
대수적 스택 위의 준연접 가군의 범주가 그 표시 위의 준연접 가군의 범주와
동치임을 보았다. 『스택 위의 층』의 절
\ref{stacks-sheaves-section-quasi-coherent-algebraic-stacks}을 보라.
이 사실이 다음 결과의 기초이다.

\begin{lemma}
\label{stacks-cohomology-lemma-adjoint}
$\mathcal{X}$를 대수적 스택이라 하자.
$\textit{LQCoh}^{fbc}(\mathcal{O}_\mathcal{X})$를 평탄 밑변환 성질을 가진
국소 준연접 가군들의 범주라 하자. 절
\ref{stacks-cohomology-section-loc-qcoh-flat-base-change}을 보라. 포함 함자
$i : \QCoh(\mathcal{O}_\mathcal{X}) \to
\textit{LQCoh}^{fbc}(\mathcal{O}_\mathcal{X})$는 우수반
$$
Q : \textit{LQCoh}^{fbc}(\mathcal{O}_\mathcal{X}) \to
\QCoh(\mathcal{O}_\mathcal{X})
$$
을 가지며, $Q \circ i$는 항등함자이다.
\end{lemma}

\begin{proof}
스킴 $U$와 매끄러운 전사 사상 $f : U \to \mathcal{X}$를 택한다.
$R = U \times_\mathcal{X} U$라 놓으면 대수공간에서의 매끄러운 준군
$(U, R, s, t, c)$를 얻으며 $\mathcal{X} = [U/R]$이다.
『대수적 스택』의 보조정리 \ref{algebraic-lemma-stack-presentation}을 보라.
$\mathcal{X}$를 $[U/R]$로 바꾸어도 되고 그렇게 하겠다.
『스택 위의 층』의 명제 \ref{stacks-sheaves-proposition-quasi-coherent}에 의해
동치
$$
q_1 :
\QCoh(U, R, s, t, c)
\longrightarrow
\QCoh(\mathcal{O}_\mathcal{X})
$$
가 있다. 이제 함자
$$
q_2 :
\textit{LQCoh}^{fbc}(\mathcal{O}_\mathcal{X})
\longrightarrow
\QCoh(U, R, s, t, c)
$$
를 다음 규칙으로 구성하자. $\mathcal{F}$가
$\textit{LQCoh}^{fbc}(\mathcal{O}_\mathcal{X})$의 대상이면
$$
q_2(\mathcal{F}) = (f^*\mathcal{F}|_{U_\etale}, \alpha)
$$
로 놓는다. 여기서 $\alpha$는 동형
$$
t_{small}^*(f^*\mathcal{F}|_{U_\etale})
\to
t^*f^*\mathcal{F}|_{R_\etale} \to
s^*f^*\mathcal{F}|_{R_\etale} \to
s_{small}^*(f^*\mathcal{F}|_{U_\etale})
$$
이고, 바깥의 두 사상은 비교 사상이다. $\mathcal{F}$가 국소 준연접이기
때문에 정확히 $q_2(\mathcal{F})$가 준연접이며, 하강 데이터 $\alpha$를
구성할 때 평탄 밑변환 성질을 사용했고 또 필요했음에 유의하라.
코사이클 조건(『공간에서의 준군』의 정의
\ref{spaces-groupoids-definition-groupoid-module}을 보라)이 성립하는지에 대한
확인은 생략한다. 『스택 위의 층』의 명제
\ref{stacks-sheaves-proposition-quasi-coherent}의 증명을 살펴보면
$q_2 \circ i$가 $q_1$의 준역임을 알 수 있다.
$Q = q_1 \circ q_2$로 정의하자.
$\mathcal{F}$를 $\textit{LQCoh}^{fbc}(\mathcal{O}_\mathcal{X})$의 대상이라 하고,
$\mathcal{G}$를 $\QCoh(\mathcal{O}_\mathcal{X})$의 대상이라 하자. 그러면
\begin{align*}
\Mor_{\textit{LQCoh}^{fbc}(\mathcal{O}_\mathcal{X})}
(i(\mathcal{G}), \mathcal{F})
& =
\Mor_{\QCoh(U, R, s, t, c)}(q_2(i(\mathcal{G})), q_2(\mathcal{F})) \\
& =
\Mor_{\QCoh(\mathcal{O}_\mathcal{X})}(\mathcal{G}, Q(\mathcal{F}))
\end{align*}
이다. 첫 번째 등식은 『스택 위의 층』의 보조정리
\ref{stacks-sheaves-lemma-map-from-quasi-coherent}이고,
두 번째 등식은 $q_1 \circ i$와 $q_2$가 범주들의 서로 준역인 동치이기 때문에
성립한다. 주장 $Q \circ i \cong \text{id}$는 $i$가 완전충실하다는 사실의
형식적 귀결이다.
\end{proof}

\begin{lemma}
\label{stacks-cohomology-lemma-adjoint-kernel-parasitic}
$\mathcal{X}$를 대수적 스택이라 하자. $Q :
\textit{LQCoh}^{fbc}(\mathcal{O}_\mathcal{X}) \to
\QCoh(\mathcal{O}_\mathcal{X})$를 보조정리 \ref{stacks-cohomology-lemma-adjoint}에서
구성한 함자라 하자.
\begin{enumerate}
\item $Q$의 핵은 정확히
$\textit{LQCoh}^{fbc}(\mathcal{O}_\mathcal{X})$의 기생 대상들의 모임이다.
\item $\textit{LQCoh}^{fbc}(\mathcal{O}_\mathcal{X})$의 임의의 대상
$\mathcal{F}$에 대해 수반 사상 $Q(\mathcal{F}) \to \mathcal{F}$의
핵과 여핵은 모두 기생이다.
\item 함자 $Q$는 완전하고 모든 극한 및 쌍대극한과 가환한다.
\end{enumerate}
\end{lemma}

\begin{proof}
보조정리 \ref{stacks-cohomology-lemma-adjoint}의 증명에서처럼 $\mathcal{X} = [U/R]$라 쓰자.
$\mathcal{F}$를 $\textit{LQCoh}^{fbc}(\mathcal{O}_\mathcal{X})$의 대상이라 하자.
보조정리 \ref{stacks-cohomology-lemma-adjoint}의 증명에서 $\mathcal{F}$가 $Q$의 핵에 속하는
것과 $\mathcal{F}|_{U_\etale} = 0$은 동치임이 명확하다.
특히 $\mathcal{F}$가 기생이면 $\mathcal{F}$는 그 핵에 속한다.
다음으로 $x : V \to \mathcal{X}$를 평탄 사상이라 하되 $V$는 스킴이라 하자.
$W = V \times_\mathcal{X} U$라 놓고 도식
$$
\xymatrix{
W \ar[d]_p \ar[r]_q & V \ar[d] \\
U \ar[r] & \mathcal{X}
}
$$
을 생각하자. 사영 $p : W \to U$는 평탄하고 사영 $q : W \to V$는
매끄럽고 전사이다. 따라서 $q_{small}^*$는 준연접 가군에서 충실한 함자이다.
가정에 따라 $\mathcal{F}$는 평탄 밑변환 성질을 가지므로
$p_{small}^*\mathcal{F}|_{U_\etale} \cong
q_{small}^*\mathcal{F}|_{V_\etale}$를 얻는다.
그러므로 $\mathcal{F}$가 $Q$의 핵에 속하면
$\mathcal{F}|_{V_\etale} = 0$이다. 이로써 (1)의 증명이 끝난다.

\medskip\noindent
(2)는 위 논의와, 사상 $Q(\mathcal{F}) \to \mathcal{F}$를
$U_\etale$에 제한하면 동형이 된다는 사실에서 따른다.

\medskip\noindent
(3)을 보기 위해 $Q$가 우수반이므로 왼쪽 완전함에 유의하라.
$0 \to \mathcal{F} \to \mathcal{G} \to \mathcal{H} \to 0$을
$\textit{LQCoh}^{fbc}(\mathcal{O}_\mathcal{X})$의 짧은 완전열이라 하자.
다음 가환 도식을 생각하자.
$$
\xymatrix{
0 \ar[r] &
Q(\mathcal{F}) \ar[r] \ar[d]_a &
Q(\mathcal{G}) \ar[r] \ar[d]_b &
Q(\mathcal{H}) \ar[r] \ar[d]_c & 0 \\
0 \ar[r] &
\mathcal{F} \ar[r] &
\mathcal{G} \ar[r] &
\mathcal{H} \ar[r] & 0
}
$$
(2)에 의해 $a$, $b$, $c$의 핵과 여핵은 기생이고 아래 행은 짧은
완전열이므로, 위 행을 $\mathcal{O}_\mathcal{X}$-가군의 복합체로 보았을 때
그 코호몰로지 층들이 기생임을 알 수 있다(세부사항은 생략한다. 이는 기생
가군들의 범주가 모든 가군의 범주의 Serre 부분범주라는 사실을 사용한다).
$Q$의 왼쪽 완전성에 의해 $Q(\mathcal{H})$에서의 완전성만 문제이다.
그러나 $Q(\mathcal{G}) \to Q(\mathcal{H}))$의 여핵 $\mathcal{Q}$는
$\textit{Mod}(\mathcal{O}_\mathcal{X})$에서 계산하든
$\QCoh(\mathcal{O}_\mathcal{X})$에서 계산하든 결과가 같다.
실제로 포함 함자 $\QCoh(\mathcal{O}_\mathcal{X}) \to
\textit{LQCoh}^{fbc}(\mathcal{O}_\mathcal{X})$는 왼쪽 수반이므로 오른쪽
완전하다. 따라서 $\mathcal{Q} = Q(\mathcal{Q})$는 준연접이면서 기생이고,
(1)에 의해 원하는 대로 $0$이다.

\medskip\noindent
우수반인 $Q$는 모든 극한과 가환한다. $Q$가 완전하므로, $Q$가 모든
쌍대극한과 가환함을 보이려면 $Q$가 직합과 가환함을 보이면 충분하다.
『범주』의 보조정리
\ref{categories-lemma-colimits-coproducts-coequalizers}를 보라.
$\mathcal{F}_i$, $i \in I$를
$\textit{LQCoh}^{fbc}(\mathcal{O}_\mathcal{X})$의 대상들의 족이라 하자.
$Q(\bigoplus \mathcal{F}_i)$가 $\bigoplus Q(\mathcal{F}_i)$와 같음을 보기 위해
보조정리 \ref{stacks-cohomology-lemma-adjoint}의 증명에 나오는 $Q$의 구성을 살펴보자.
여기서는 $U$가 스킴인 표시 $\mathcal{X} = [U/R]$를 사용한다.
그러면 먼저 $\QCoh(U, R, s, t, c)$에서 쌍
$(\mathcal{F}|_{U_\etale}, \alpha)$를 취하고, 이어서 동치
$\QCoh(U, R, s, t, c) \cong \QCoh(\mathcal{O}_\mathcal{X})$를 사용하여
$Q(\mathcal{F})$를 계산한다. 제한 함자
$\textit{Mod}(\mathcal{O}_\mathcal{X}) \to
\textit{Mod}(\mathcal{O}_{U_\etale})$,
$\mathcal{F} \mapsto \mathcal{F}|_{U_\etale}$가 직합과 가환하므로
원하는 등식은 명확하다.
\end{proof}

\begin{lemma}
\label{stacks-cohomology-lemma-flat-morphism-and-Q}
$f : \mathcal{X} \to \mathcal{Y}$를 대수적 스택의 평탄 사상이라 하자.
그러면 $Q_\mathcal{X} \circ f^* = f^* \circ Q_\mathcal{Y}$이다.
여기서 $Q_\mathcal{X}$와 $Q_\mathcal{Y}$는 보조정리
\ref{stacks-cohomology-lemma-adjoint}에서와 같다.
\end{lemma}

\begin{proof}
$f^*$는 $\QCoh$와 $\textit{LQCoh}^{fbc}$를 모두 보존한다.
『스택 위의 층』의 보조정리
\ref{stacks-sheaves-lemma-pullback-quasi-coherent}와 명제
\ref{stacks-cohomology-proposition-loc-qcoh-flat-base-change}를 보라.
$\mathcal{F}$가 $\textit{LQCoh}^{fbc}(\mathcal{O}_\mathcal{Y})$에 속하면,
보조정리 \ref{stacks-cohomology-lemma-adjoint-kernel-parasitic}에 의해
$Q_\mathcal{Y}(\mathcal{F}) \to \mathcal{F}$의 핵과 여핵은 기생이다.
$f$가 평탄하므로 보조정리 \ref{stacks-cohomology-lemma-parasitic}에 의해
$f^*Q_\mathcal{Y}(\mathcal{F}) \to f^*\mathcal{F}$의 핵과 여핵은 기생이다.
따라서 유도된 사상
$f^*Q_\mathcal{Y}(\mathcal{F}) \to Q_\mathcal{X}(f^*\mathcal{F})$의
핵과 여핵은 기생이고, 예를 들어 보조정리
\ref{stacks-cohomology-lemma-exact-sequence-quasi-coherent-parasitic-cohomology}에 의해 동형이다.
\end{proof}

\begin{lemma}
\label{stacks-cohomology-lemma-flat-object-and-Q}
$\mathcal{X}$를 대수적 스택이라 하자. $x$를 스킴 $U$ 위에 놓인
$\mathcal{X}$의 대상이라 하고 $x : U \to \mathcal{X}$가 평탄하다고 하자.
그러면 $\QCoh^{fbc}(\mathcal{O}_\mathcal{X})$의 $\mathcal{F}$에 대해
$Q(\mathcal{F})|_{U_\etale} = \mathcal{F}|_{U_\etale}$이다.
\end{lemma}

\begin{proof}
$Q(\mathcal{F}) \to \mathcal{F}$의 핵과 여핵이 기생이기 때문에 참이다.
보조정리 \ref{stacks-cohomology-lemma-adjoint-kernel-parasitic}를 보라.
\end{proof}

\begin{remark}
\label{stacks-cohomology-remark-QCoh-abelian}
$\mathcal{X}$를 대수적 스택이라 하자. 범주
$\QCoh(\mathcal{O}_\mathcal{X})$는 아벨 범주이고, 포함 함자
$\QCoh(\mathcal{O}_\mathcal{X}) \to \textit{Mod}(\mathcal{O}_\mathcal{X})$는
오른쪽 완전하지만 일반적으로 완전하지 않다. 『스택 위의 층』의 보조정리
\ref{stacks-sheaves-lemma-quasi-coherent-algebraic-stack}을 보라.
보조정리 \ref{stacks-cohomology-lemma-adjoint}와 \ref{stacks-cohomology-lemma-adjoint-kernel-parasitic}의 함자
$Q$를 사용하여 이를 이해할 수 있다. 구체적으로
$\varphi : \mathcal{F} \to \mathcal{G}$를 준연접
$\mathcal{O}_\mathcal{X}$-가군들의 사상이라 하자. 그러면 다음이 성립한다.
\begin{enumerate}
\item $\textit{Mod}(\mathcal{O}_\mathcal{X})$에서 계산한 여핵
$\Coker(\varphi)$는 준연접이고, $\QCoh(\mathcal{O}_\mathcal{X})$에서의
$\varphi$의 여핵이다.
\item $\textit{Mod}(\mathcal{O}_\mathcal{X})$에서 계산한 상
$\Im(\varphi)$는 준연접이고, $\QCoh(\mathcal{O}_\mathcal{X})$에서의
$\varphi$의 상이다.
\item $\textit{Mod}(\mathcal{O}_\mathcal{X})$에서 계산한 핵
$\Ker(\varphi)$는 명제 \ref{stacks-cohomology-proposition-loc-qcoh-flat-base-change}에 의해
$\textit{LQCoh}^{fbc}(\mathcal{O}_\mathcal{X})$에 속하고,
$Q(\Ker(\varphi))$는 $\QCoh(\mathcal{O}_\mathcal{X})$에서의 핵이다.
\end{enumerate}
이는 앞서 제시한 참고문헌에서 따른다.
\end{remark}

\begin{remark}
\label{stacks-cohomology-remark-QCoh-tensor}
$\mathcal{X}$를 대수적 스택이라 하자. 두 준연접
$\mathcal{O}_\mathcal{X}$-가군 $\mathcal{F}$와 $\mathcal{G}$가 주어지면
텐서곱 가군 $\mathcal{F} \otimes_{\mathcal{O}_\mathcal{X}} \mathcal{G}$는
준연접이다. 『스택 위의 층』의 보조정리
\ref{stacks-sheaves-lemma-quasi-coherent-algebraic-stack}의 (5)를 보라.
마찬가지로 평탄 밑변환 성질을 가진 두 국소 준연접 가군이 주어지면
그 텐서곱도 같은 성질을 가진다. 명제
\ref{stacks-cohomology-proposition-loc-qcoh-flat-base-change}를 보라. 따라서 포함 함자
$$
\QCoh(\mathcal{O}_\mathcal{X}) \to
\textit{LQCoh}^{fbc}(\mathcal{O}_\mathcal{X}) \to
\textit{Mod}(\mathcal{O}_\mathcal{X})
$$
는 대칭 모노이드 범주들의 함자이다. 더 흥미로운 사실은 함자
$$
Q :
\textit{LQCoh}^{fbc}(\mathcal{O}_\mathcal{X})
\longrightarrow
\QCoh(\mathcal{O}_\mathcal{X})
$$
도 대칭 모노이드 범주들의 함자라는 것이다. 구체적으로
$\textit{LQCoh}^{fbc}(\mathcal{O}_\mathcal{X})$의
$\mathcal{F}$와 $\mathcal{G}$가 주어지면
$$
\xymatrix{
Q(\mathcal{F})
\otimes_{\mathcal{O}_\mathcal{X}}
Q(\mathcal{G}) \ar[rr] \ar[rd] & &
\mathcal{F}
\otimes_{\mathcal{O}_\mathcal{X}}
\mathcal{G} \\
&
Q(\mathcal{F} \otimes_{\mathcal{O}_\mathcal{X}} \mathcal{G}) \ar[ru]
}
$$
를 얻는다. 여기서 남서쪽 화살은 북서쪽 화살의 보편 성질에서 온다
(그리고 왼쪽 위 모서리의 대상이 준연접이라는 앞서 언급한 사실을 사용한다).
이 도식을 평탄한 $U \to \mathcal{X}$에 대해 $U_\etale$로 제한하면
세 화살 모두 동형이 된다(보조정리 \ref{stacks-cohomology-lemma-adjoint}와
\ref{stacks-cohomology-lemma-adjoint-kernel-parasitic}, 그리고 정의
\ref{stacks-cohomology-definition-parasitic}을 보라). 따라서
$Q(\mathcal{F}) \otimes_{\mathcal{O}_\mathcal{X}}
Q(\mathcal{G}) \to
Q(\mathcal{F} \otimes_{\mathcal{O}_\mathcal{X}} \mathcal{G})$는 동형이다.
예를 들어 보조정리 \ref{stacks-cohomology-lemma-quasi-coherent-check-exact}를 보라.
\end{remark}

\begin{remark}
\label{stacks-cohomology-remark-bousfield-colocalization}
$\mathcal{X}$를 대수적 스택이라 하자.
$\textit{Parasitic}(\mathcal{O}_\mathcal{X}) \subset
\textit{Mod}(\mathcal{O}_\mathcal{X})$를 기생 가군들로 이루어진
충만한 부분범주라 하자. 보조정리 \ref{stacks-cohomology-lemma-adjoint}와
\ref{stacks-cohomology-lemma-adjoint-kernel-parasitic}의 결과는
$$
\QCoh(\mathcal{O}_\mathcal{X}) =
\textit{LQCoh}^{fbc}(\mathcal{O}_\mathcal{X}) /
\textit{Parasitic}(\mathcal{O}_\mathcal{X})
\cap \textit{LQCoh}^{fbc}(\mathcal{O}_\mathcal{X})
$$
임을 함의한다. 말로 하면, 준연접 가군의 범주는 평탄 밑변환 성질을 가진
국소 준연접 가군의 범주를 기생 대상들로 이루어진 Serre 부분범주로 나눈 것이다.
『호몰로지』의 보조정리 \ref{homology-lemma-serre-subcategory-is-kernel}을 보라.
몫 함자의 왼쪽 수반인 포함 함자
$i : \QCoh(\mathcal{O}_\mathcal{X}) \to
\textit{LQCoh}^{fbc}(\mathcal{O}_\mathcal{X})$가 존재한다는 것이 이 상황의
핵심 특징이다. 『스택의 유도 범주』의 절
\ref{stacks-perfect-section-derived}, 특히 보조정리
\ref{stacks-perfect-lemma-bousfield-colocalization}에서는
유도 범주의 수준에서도 유사한 결과가 성립함을 증명한다.
\end{remark}

\begin{lemma}
\label{stacks-cohomology-lemma-internal-hom-fp-into-qcoh}
$\mathcal{X}$를 대수적 스택이라 하자. $\mathcal{F}$를 유한 표시인
$\mathcal{O}_\mathcal{X}$-가군이라 하고, $\mathcal{G}$를 준연접
$\mathcal{O}_\mathcal{X}$-가군이라 하자.
$\textit{Mod}(\mathcal{X}_\etale, \mathcal{O}_\mathcal{X})$와
$\textit{Mod}(\mathcal{O}_\mathcal{X})$에서 계산한 내부 Hom
$\SheafHom_{\mathcal{O}_\mathcal{X}}(\mathcal{F}, \mathcal{G})$들은
일치하며, 그 공통값은
$\textit{LQCoh}^{fbc}(\mathcal{O}_\mathcal{X})$의 대상이다.
준연접 가군
$
hom(\mathcal{F}, \mathcal{G}) =
Q(\SheafHom_{\mathcal{O}_\mathcal{X}}(\mathcal{F}, \mathcal{G}))
$
은 다음 보편 성질을 가진다.
$$
\Hom_\mathcal{X}(\mathcal{H}, hom(\mathcal{F}, \mathcal{G})) =
\Hom_\mathcal{X}(\mathcal{H} \otimes_{\mathcal{O}_\mathcal{X}} \mathcal{F},
\mathcal{G})
$$
여기서 $\mathcal{H}$는 $\QCoh(\mathcal{O}_\mathcal{X})$에 속한다.
\end{lemma}

\begin{proof}
『사이트 위의 가군』의 절 \ref{sites-modules-section-internal-hom}에 나오는
$\SheafHom_{\mathcal{O}_\mathcal{X}}(\mathcal{F}, \mathcal{G})$의 구성은
가군 준층으로서의 $\mathcal{F}$와 $\mathcal{G}$에만 의존한다.
$\mathcal{F}$와 $\mathcal{G}$가 fppf 위상의 층이라고 가정했으므로,
그 출력 $\SheafHom$은 fppf 위상의 층이다. 『사이트 위의 가군』의 보조정리
\ref{sites-modules-lemma-internal-hom}을 보라.
『스택 위의 층』의 보조정리 \ref{stacks-sheaves-lemma-lqc-colimits}에 의해
$\SheafHom_{\mathcal{O}_\mathcal{X}}(\mathcal{F}, \mathcal{G})$는
국소 준연접이다. 보조정리
\ref{stacks-cohomology-lemma-check-flat-comparison-on-etale-covering}에 의해
$\SheafHom_{\mathcal{O}_\mathcal{X}}(\mathcal{F}, \mathcal{G})$는
평탄 밑변환 성질을 가진다. 따라서
$\SheafHom_{\mathcal{O}_\mathcal{X}}(\mathcal{F}, \mathcal{G})$는
$\textit{LQCoh}^{fbc}(\mathcal{O}_\mathcal{X})$의 대상이고,
여기에 보조정리 \ref{stacks-cohomology-lemma-adjoint}의 함자 $Q$를 적용할 수 있다.
$Q$의 보편 성질에 의해, 준연접인 $\mathcal{H}$에 대해
$$
\Hom_\mathcal{X}(\mathcal{H},
Q(\SheafHom_{\mathcal{O}_\mathcal{X}}(\mathcal{F}, \mathcal{G}))) =
\Hom_\mathcal{X}(\mathcal{H},
\SheafHom_{\mathcal{O}_\mathcal{X}}(\mathcal{F}, \mathcal{G}))
$$
이다. 그러므로 보조정리의 표시된 식은 『사이트 위의 가군』의 보조정리
\ref{sites-modules-lemma-internal-hom-adjoint-tensor}에서 따른다.
\end{proof}

\begin{lemma}
\label{stacks-cohomology-lemma-flat-morphism-and-hom}
$f : \mathcal{X} \to \mathcal{Y}$를 대수적 스택의 평탄 사상이라 하자.
$\mathcal{F}$를 유한 표시인 $\mathcal{O}_\mathcal{Y}$-가군이라 하고,
$\mathcal{G}$를 준연접 $\mathcal{O}_\mathcal{Y}$-가군이라 하자.
그러면 보조정리 \ref{stacks-cohomology-lemma-internal-hom-fp-into-qcoh}의 표기법으로
$f^*hom(\mathcal{F}, \mathcal{G}) = hom(f^*\mathcal{F}, f^*\mathcal{G})$이다.
\end{lemma}

\begin{proof}
『사이트 위의 가군』의 보조정리
\ref{sites-modules-lemma-pullback-internal-hom}에 의해
$f^*\SheafHom_{\mathcal{O}_\mathcal{Y}}(\mathcal{F}, \mathcal{G}) =
\SheafHom_{\mathcal{O}_\mathcal{X}}(f^*\mathcal{F}, f^*\mathcal{G})$이다.
(이 단계에서 $f$의 평탄성을 사용하는 것은 아님에 유의하라. $f$에 대응하는
환 달린 토포스의 사상은 항상 평탄하기 때문이다. 『스택 위의 층』의 주
\ref{stacks-sheaves-remark-flat}을 보라.)
이제 보조정리 \ref{stacks-cohomology-lemma-flat-morphism-and-Q}를 적용한다
(그리고 여기서는 $f$의 평탄성을 사용한다).
\end{proof}

\section{준연접 가군의 직상}
\label{stacks-cohomology-section-pushforward-quasi-coherent}

\noindent
$f : \mathcal{X} \to \mathcal{Y}$를 대수적 스택의 사상이라 하자.
직상
$$
f_* :
\textit{Mod}(\mathcal{O}_\mathcal{X})
\longrightarrow
\textit{Mod}(\mathcal{O}_\mathcal{Y})
$$
을 생각하자. 이 함자는 준연접 층들의 부분범주를 거의 절대로 보존하지
않는다는 사실이 드러난다. 예를 들어 스킴의 사상
$$
j : X = \mathbf{A}^2_k \setminus \{0\} \longrightarrow \mathbf{A}^2_k = Y.
$$
를 생각하자. 이에 대응하는 대수적 스택의 사상
$$
f = j_{big} : \mathcal{X} = (\Sch/X)_{fppf} \to
(\Sch/Y)_{fppf} = \mathcal{Y}
$$
이 있다. 구조층의 직상 $f_*\mathcal{O}_\mathcal{X}$의 대역 단면은
$k[x, y]$이다. 따라서 $f_*\mathcal{O}_\mathcal{X}$가 $\mathcal{Y}$ 위에서
준연접이라면 $f_*\mathcal{O}_\mathcal{X} = \mathcal{O}_\mathcal{Y}$이어야 한다.
그러나 $0$으로 가는 $T = \Spec(k) \to \mathbf{A}^2_k = Y$를 생각하자.
$X \times_Y T = \emptyset$이므로
$\Gamma(T, f_*\mathcal{O}_\mathcal{X}) = 0$인 반면,
$\Gamma(T, \mathcal{O}_\mathcal{Y}) = k$이다.
긍정적인 면으로, 임의의 평탄 사상 $T \to Y$에 대해
$\Gamma(T, f_*\mathcal{O}_\mathcal{X}) = \Gamma(T, \mathcal{O}_\mathcal{Y})$이다.
이는 $j$가 준콤팩트이고 준분리라는 사실과 『스킴의 코호몰로지』의 보조정리
\ref{coherent-lemma-flat-base-change-cohomology}에서 따른다.

\medskip\noindent
$f : \mathcal{X} \to \mathcal{Y}$를 대수적 스택의 준콤팩트 준분리 사상이라 하자.
위에서 언급한 문제는 다음 세 가지 관찰을 사용하여 우회한다.
\begin{enumerate}
\item $f_*$는 국소 준연접 가군을 보존한다
(보조정리 \ref{stacks-cohomology-lemma-pushforward-locally-quasi-coherent}).
\item $f_*$는 준연접 층을 그 평탄 비교 사상들이 동형인 국소 준연접 층으로
보낸다(보조정리 \ref{stacks-cohomology-lemma-flat-comparison}).
\item 평탄 밑변환 성질을 가진 국소 준연접
$\mathcal{O}_\mathcal{Y}$-가군은 $\mathcal{Y}$의 표시 위의 준연접 가군을,
따라서 $\mathcal{Y}$ 위의 준연접 가군을 준다. 『스택 위의 층』의 절
\ref{stacks-sheaves-section-quasi-coherent-algebraic-stacks}을 보라.
\end{enumerate}
따라서 함자
$$
f_{\QCoh, *} :
\QCoh(\mathcal{O}_\mathcal{X})
\longrightarrow
\QCoh(\mathcal{O}_\mathcal{Y})
$$
를 얻는다. 이는
$f^* : \QCoh(\mathcal{O}_\mathcal{Y}) \to
\QCoh(\mathcal{O}_\mathcal{X})$의 우수반이며, 더 나아가 연관된 $1$-사상
$y : V \to \mathcal{Y}$가 평탄한 임의의 $y \in \Ob(\mathcal{Y})$에 대해
$$
\Gamma(y, f_*\mathcal{F}) = \Gamma(y, f_{\QCoh, *}\mathcal{F})
$$
이다. 보조정리 \ref{stacks-cohomology-lemma-direct-image-quasi-coherent-over-flat}을 보라.
더욱이 유사한 구성으로 함자 $R^if_{\QCoh, *}$를 만들 수 있다.
그러나 이 결과들은 (준연접 코호몰로지 층을 가진 복합체의) 전체 직상 함자를
만들기에는 충분하지 않다.

\begin{proposition}
\label{stacks-cohomology-proposition-direct-image-quasi-coherent}
$f : \mathcal{X} \to \mathcal{Y}$를 대수적 스택의 준콤팩트 준분리 사상이라 하자.
함자 $f^* : \QCoh(\mathcal{O}_\mathcal{Y}) \to
\QCoh(\mathcal{O}_\mathcal{X})$는 우수반
$$
f_{\QCoh, *} :
\QCoh(\mathcal{O}_\mathcal{X})
\longrightarrow
\QCoh(\mathcal{O}_\mathcal{Y})
$$
을 가진다. 이는 합성
$$
\QCoh(\mathcal{O}_\mathcal{X}) \to
\textit{LQCoh}^{fbc}(\mathcal{O}_\mathcal{X})
\xrightarrow{f_*} \textit{LQCoh}^{fbc}(\mathcal{O}_\mathcal{Y})
\xrightarrow{Q} \QCoh(\mathcal{O}_\mathcal{Y})
$$
으로 정의할 수 있다. 여기서 함자 $f_*$와 $Q$는 각각 명제
\ref{stacks-cohomology-proposition-loc-qcoh-flat-base-change}와 보조정리
\ref{stacks-cohomology-lemma-adjoint}에서와 같다. 더욱이 $R^if_{\QCoh, *}$를 합성
$$
\QCoh(\mathcal{O}_\mathcal{X}) \to
\textit{LQCoh}^{fbc}(\mathcal{O}_\mathcal{X})
\xrightarrow{R^if_*} \textit{LQCoh}^{fbc}(\mathcal{O}_\mathcal{Y})
\xrightarrow{Q} \QCoh(\mathcal{O}_\mathcal{Y})
$$
으로 정의하면, 함자들의 열 $\{R^if_{\QCoh, *}\}_{i \geq 0}$은
코호몰로지적 $\delta$-함자를 이룬다.
\end{proposition}

\begin{proof}
이는 명제에서 언급한 결과들을 결합한 것이다.
수반성은 다음과 같이 보일 수 있다. $\mathcal{F}$를 준연접
$\mathcal{O}_\mathcal{X}$-가군이라 하고, $\mathcal{G}$를 준연접
$\mathcal{O}_\mathcal{Y}$-가군이라 하자. 그러면
\begin{align*}
\Mor_{\QCoh(\mathcal{O}_\mathcal{X})}(f^*\mathcal{G}, \mathcal{F})
& =
\Mor_{\textit{LQCoh}^{fbc}(\mathcal{O}_\mathcal{Y})}
(\mathcal{G}, f_*\mathcal{F}) \\
& =
\Mor_{\QCoh(\mathcal{O}_\mathcal{Y})}(\mathcal{G}, Q(f_*\mathcal{F}))
\\
& =
\Mor_{\QCoh(\mathcal{O}_\mathcal{Y})}(\mathcal{G},
f_{\QCoh, *}\mathcal{F})
\end{align*}
이다. 첫 번째 등식은 (임의의 가군층에 대한) $f_*$와 $f^*$의 수반성에서
따른다. 명제 \ref{stacks-cohomology-proposition-loc-qcoh-flat-base-change}에 의해
$f_*\mathcal{F}$는
$\textit{LQCoh}^{fbc}(\mathcal{O}_\mathcal{Y})$의 대상이고
(fppf 위상과 에탈 위상 어느 쪽에서도 계산할 수 있다), 보조정리
\ref{stacks-cohomology-lemma-adjoint}에서 두 번째 등식을 얻는다. 세 번째 등식은
$f_{\QCoh, *}$의 정의이다.

\medskip\noindent
$\{R^if_{\QCoh, *}\}_{i \geq 0}$가 『호몰로지』의 정의
\ref{homology-definition-cohomological-delta-functor}에서 정의한
코호몰로지적 $\delta$-함자임을 보이기 위해
$$
0 \to \mathcal{F}_1 \to \mathcal{F}_2 \to \mathcal{F}_3 \to 0
$$
을 $\QCoh(\mathcal{O}_\mathcal{X})$의 짧은 완전열이라 하자.
이 열은 $\textit{Mod}(\mathcal{O}_\mathcal{X})$에서 완전열이 아닐 수 있지만,
기생 가군들을 무시하면 완전임을 알고 있다. 보조정리
\ref{stacks-cohomology-lemma-exact-sequence-quasi-coherent-parasitic-cohomology}를 보라.
따라서 이 열을 $\mathcal{P}_i$가 기생인
$\textit{Mod}(\mathcal{O}_\mathcal{X})$의 짧은 완전열들
$$
\begin{matrix}
0 \to \mathcal{P}_1 \to \mathcal{F}_1 \to \mathcal{I}_2 \to 0 \\
0 \to \mathcal{I}_2 \to \mathcal{F}_2 \to \mathcal{Q}_2 \to 0 \\
0 \to \mathcal{P}_2 \to \mathcal{Q}_2 \to \mathcal{I}_3 \to 0 \\
0 \to \mathcal{I}_3 \to \mathcal{F}_3 \to \mathcal{P}_3 \to 0
\end{matrix}
$$
로 나눌 수 있다. 각 층
$\mathcal{P}_j$, $\mathcal{I}_j$, $\mathcal{Q}_j$는 명제
\ref{stacks-cohomology-proposition-loc-qcoh-flat-base-change}에 의해
$\textit{LQCoh}^{fbc}(\mathcal{O}_\mathcal{X})$의 대상임에 유의하라.
$R^if_*$를 적용하면 긴 완전열들
$$
\begin{matrix}
0 \to f_*\mathcal{P}_1 \to f_*\mathcal{F}_1 \to f_*\mathcal{I}_2 \to
R^1f_*\mathcal{P}_1 \to \ldots \\
0 \to f_*\mathcal{I}_2 \to f_*\mathcal{F}_2 \to f_*\mathcal{Q}_2 \to
R^1f_*\mathcal{I}_2 \to \ldots \\
0 \to f_*\mathcal{P}_2 \to f_*\mathcal{Q}_2 \to f_*\mathcal{I}_3 \to
R^1f_*\mathcal{P}_2 \to \ldots \\
0 \to f_*\mathcal{I}_3 \to f_*\mathcal{F}_3 \to f_*\mathcal{P}_3 \to
R^1f_*\mathcal{I}_3 \to \ldots
\end{matrix}
$$
을 얻는데, 항들이 명제 \ref{stacks-cohomology-proposition-loc-qcoh-flat-base-change}에 의해
$\textit{LQCoh}^{fbc}(\mathcal{O}_\mathcal{Y})$의 대상이다.
보조정리 \ref{stacks-cohomology-lemma-pushforward-parasitic}에 의해 층
$R^if_*\mathcal{P}_j$는 기생이므로, 보조정리
\ref{stacks-cohomology-lemma-adjoint-kernel-parasitic}에 의해 함자 $Q$를 적용하면 소멸한다.
$Q$는 완전하므로 사상
$$
Q(R^if_*\mathcal{F}_3)
\cong
Q(R^if_*\mathcal{I}_3)
\cong
Q(R^if_*\mathcal{Q}_2)
\rightarrow
Q(R^{i + 1}f_*\mathcal{I}_2)
\cong
Q(R^{i + 1}f_*\mathcal{F}_1)
$$
을 연결 사상으로 삼아 함자족
$\{R^if_{\QCoh, *}\}_{i \geq 0}$을 코호몰로지적 $\delta$-함자로 만들 수 있다.
\end{proof}

\begin{lemma}
\label{stacks-cohomology-lemma-direct-image-quasi-coherent-over-flat}
$f : \mathcal{X} \to \mathcal{Y}$를 대수적 스택의 준콤팩트 준분리 사상이라 하자.
$y : V \to \mathcal{Y}$가 $\Ob(\mathcal{Y})$에 속하고 $y$가 평탄 사상이라 하자.
$\mathcal{F}$가 $\QCoh(\mathcal{O}_\mathcal{X})$에 속한다고 하자.
그러면 모든 $i \in \mathbf{Z}$에 대해
$(f_*\mathcal{F})(y) = (f_{\QCoh, *}\mathcal{F})(y)$이고
$(R^if_*\mathcal{F})(y) = (R^if_{\QCoh, *}\mathcal{F})(y)$이다.
\end{lemma}

\begin{proof}
이는 명제 \ref{stacks-cohomology-proposition-direct-image-quasi-coherent}의 함자
$R^if_{\QCoh, *}$의 구성, 정의 \ref{stacks-cohomology-definition-parasitic}의 기생 가군의 정의,
그리고 보조정리 \ref{stacks-cohomology-lemma-adjoint-kernel-parasitic}의 (2)에서 따른다.
\end{proof}

\begin{remark}
\label{stacks-cohomology-remark-direct-image-quasi-coherent-tensor}
$f : \mathcal{X} \to \mathcal{Y}$를 대수적 스택의 준콤팩트 준분리 사상이라 하자.
$\mathcal{F}$와 $\mathcal{G}$가
$\QCoh(\mathcal{O}_\mathcal{X})$에 속한다고 하자. 그러면 표준 가환 도식
$$
\xymatrix{
f_{\QCoh, *}\mathcal{F}
\otimes_{\mathcal{O}_\mathcal{Y}}
f_{\QCoh, *}\mathcal{G} \ar[r] \ar[d] &
f_*\mathcal{F}
\otimes_{\mathcal{O}_\mathcal{Y}}
f_*\mathcal{G} \ar[d]^c \\
f_{\QCoh, *}(\mathcal{F}
\otimes_{\mathcal{O}_\mathcal{X}}
\mathcal{G}) \ar[r] &
f_*(\mathcal{F}
\otimes_{\mathcal{O}_\mathcal{X}}
\mathcal{G})
}
$$
이 있다. 오른쪽의 수직 화살 $c$는 ($0$차에서의) 소박한 상대적 컵곱이다.
『사이트 위의 코호몰로지』의 절
\ref{sites-cohomology-section-cup-product}을 보라.
$c$의 원천과 목표는 명제 \ref{stacks-cohomology-proposition-loc-qcoh-flat-base-change}에 의해
$\textit{LQCoh}^{fbc}(\mathcal{O}_\mathcal{X})$에 속한다.
$Q$가 텐서곱과 가환하므로 $c$에 $Q$를 적용하여 왼쪽의 수직 화살을 얻는다.
주 \ref{stacks-cohomology-remark-QCoh-tensor}를 보라. 이 구성은 $\mathcal{F}$와
$\mathcal{G}$에 대해 함자적이다.
\end{remark}

\begin{lemma}
\label{stacks-cohomology-lemma-leray}
$f : \mathcal{X} \to \mathcal{Y}$를 대수적 스택의 준콤팩트 준분리 사상이라 하자.
$\mathcal{F}$를 $\mathcal{X}$ 위의 준연접 층이라 하자. 그러면
$E_2$-쪽이
$$
E_2^{p, q} = H^p(\mathcal{Y}, R^qf_{\QCoh, *}\mathcal{F})
$$
이고 $H^{p + q}(\mathcal{X}, \mathcal{F})$로 수렴하는 스펙트럼 열이 존재한다.
\end{lemma}

\begin{proof}
『사이트 위의 코호몰로지』의 보조정리
\ref{sites-cohomology-lemma-Leray}에 의해 다음
$$
E_2^{p, q} = H^p(\mathcal{Y}, R^qf_*\mathcal{F})
$$
인 Leray 스펙트럼 열은 $H^{p + q}(\mathcal{X}, \mathcal{F})$로 수렴한다.
수반 사상
$$
R^qf_{\QCoh, *}\mathcal{F} \longrightarrow R^qf_*\mathcal{F}
$$
의 핵과 여핵은 $\mathcal{Y}$ 위의 기생 가군이고
(보조정리 \ref{stacks-cohomology-lemma-adjoint-kernel-parasitic}), 따라서 그 코호몰로지는
소멸한다(보조정리 \ref{stacks-cohomology-lemma-pushforward-parasitic}). 형식적으로
$H^p(\mathcal{Y}, R^qf_{\QCoh, *}\mathcal{F}) =
H^p(\mathcal{Y}, R^qf_*\mathcal{F})$가 따르며 결론을 얻는다.
\end{proof}

\begin{lemma}
\label{stacks-cohomology-lemma-relative-leray}
$f : \mathcal{X} \to \mathcal{Y}$와 $g : \mathcal{Y} \to \mathcal{Z}$를
대수적 스택의 준콤팩트 준분리 사상이라 하자.
$\mathcal{F}$를 $\mathcal{X}$ 위의 준연접 층이라 하자. 그러면
$E_2$-쪽이
$$
E_2^{p, q} = R^pg_{\QCoh, *}(R^qf_{\QCoh, *}\mathcal{F})
$$
이고 $R^{p + q}(g \circ f)_{\QCoh, *}\mathcal{F}$로 수렴하는
스펙트럼 열이 존재한다.
\end{lemma}

\begin{proof}
『사이트 위의 코호몰로지』의 보조정리
\ref{sites-cohomology-lemma-relative-Leray}에 의해 다음
$$
E_2^{p, q} = R^pg_*(R^qf_*\mathcal{F})
$$
인 Leray 스펙트럼 열은 $R^{p + q}(g \circ f)_*\mathcal{F}$로 수렴한다.
명제 \ref{stacks-cohomology-proposition-loc-qcoh-flat-base-change}의 결과들에 의해 이 스펙트럼
열의 모든 항은 $\textit{LQCoh}^{fbc}(\mathcal{O}_\mathcal{Z})$의 대상이다.
완전 함자 $Q_\mathcal{Z} :
\textit{LQCoh}^{fbc}(\mathcal{O}_\mathcal{Z}) \to
\QCoh(\mathcal{O}_\mathcal{Z})$를 적용하면
$R^{p + q}(g \circ f)_{\QCoh, *}\mathcal{F}$를 덮는
$\QCoh(\mathcal{O}_\mathcal{Z})$ 안의 스펙트럼 열을 얻는다. 따라서
$$
Q_\mathcal{Z}(R^pg_*(R^qf_*\mathcal{F})) =
Q_\mathcal{Z}(R^pg_*(Q_\mathcal{X}(R^qf_*\mathcal{F}))
$$
임을 보이면 결과가 따른다. 이는 사상
$$
Q_\mathcal{X}(R^qf_*\mathcal{F}) \longrightarrow R^qf_*\mathcal{F}
$$
의 핵과 여핵이 기생이라는 사실(보조정리
\ref{stacks-cohomology-lemma-adjoint-kernel-parasitic})과, $R^pg_*$가 기생 가군을 기생 가군으로
보낸다는 사실(보조정리 \ref{stacks-cohomology-lemma-pushforward-parasitic})에서 따른다.
\end{proof}

\noindent
이 절을 마치기 위해 스킴에 의한 매끄러운 피복에 대응하는 스펙트럼 열들을
명시하겠다. 『스택 위의 층』의 절
\ref{stacks-sheaves-section-cohomology}와
\ref{stacks-sheaves-section-higher-direct-images}를 비교하라.

\begin{proposition}
\label{stacks-cohomology-proposition-smooth-covering-compute-cohomology}
$f : \mathcal{U} \to \mathcal{X}$를 대수적 스택의 사상이라 하자.
$f$가 대수공간으로 표현가능하고, 전사이고, 평탄하며, 국소 유한 표시라고
가정하자. $\mathcal{F}$를 준연접 $\mathcal{O}_\mathcal{X}$-가군이라 하자.
그러면 스펙트럼 열
$$
E_2^{p, q} = H^q(\mathcal{U}_p, f_p^*\mathcal{F})
\Rightarrow
H^{p + q}(\mathcal{X}, \mathcal{F})
$$
이 있다. 여기서 $f_p$는 사상
$\mathcal{U} \times_\mathcal{X} \ldots \times_\mathcal{X} \mathcal{U} \to
\mathcal{X}$이다(인자는 $p + 1$개이다).
\end{proposition}

\begin{proof}
이는 『스택 위의 층』의 명제
\ref{stacks-sheaves-proposition-smooth-covering-compute-cohomology}의 특수한 경우이다.
\end{proof}

\begin{proposition}
\label{stacks-cohomology-proposition-smooth-covering-compute-direct-image}
$f : \mathcal{U} \to \mathcal{X}$와 $g : \mathcal{X} \to \mathcal{Y}$를
합성 가능한 대수적 스택의 사상이라 하자. 다음을 가정하자.
\begin{enumerate}
\item $f$는 대수공간으로 표현가능하고, 전사이고, 평탄하며,
국소 유한 표시이고, 준콤팩트이며, 준분리이다.
\item $g$는 준콤팩트이고 준분리이다.
\end{enumerate}
$\mathcal{F}$가 $\QCoh(\mathcal{O}_\mathcal{X})$에 속하면
$\QCoh(\mathcal{O}_\mathcal{Y})$ 안의 스펙트럼 열
$$
E_2^{p, q} = R^q(g \circ f_p)_{\QCoh, *}f_p^*\mathcal{F}
\Rightarrow
R^{p + q}g_{\QCoh, *}\mathcal{F}
$$
이 있다.
\end{proposition}

\begin{proof}
각 사상
$f_p : \mathcal{U} \times_\mathcal{X} \ldots \times_\mathcal{X} \mathcal{U} \to
\mathcal{X}$는 준콤팩트이고 준분리이며, 따라서 $g \circ f_p$도
준콤팩트이고 준분리이다. 그러므로 주장은 의미가 있다
(즉, 함자 $R^q(g \circ f_p)_{\QCoh, *}$가 정의된다).
『스택 위의 층』의 명제
\ref{stacks-sheaves-proposition-smooth-covering-compute-direct-image}에 의해
스펙트럼 열
$$
E_2^{p, q} = R^q(g \circ f_p)_*f_p^{-1}\mathcal{F}
\Rightarrow
R^{p + q}g_*\mathcal{F}
$$
이 있다. 완전 함자
$Q_\mathcal{Y} : \textit{LQCoh}^{fbc}(\mathcal{O}_\mathcal{Y}) \to
\QCoh(\mathcal{O}_\mathcal{Y})$를 적용하면 원하는
$\QCoh(\mathcal{O}_\mathcal{Y})$ 안의 스펙트럼 열을 얻는다.
\end{proof}

\section{준연접 가군에 관한 추가 주}
\label{stacks-cohomology-section-further-remarks}

\noindent
이 절에서는 대수적 스택 위의 준연접 가군을 사용하는 방법을 이해하는 데
도우려는 몇 가지 결과를 모은다.

\medskip\noindent
$f : \mathcal{U} \to \mathcal{X}$를 대수적 스택의 사상이라 하자.
$\mathcal{U}$가 대수공간 $U$로 표현된다고 가정하자. 당김
(『스택 위의 층』의 절 \ref{stacks-sheaves-section-modules})에 이어 제한
(『스택 위의 층』의 절
\ref{stacks-sheaves-section-restriction-algebraic-spaces})으로 주어지는 함자
$$
a :
\textit{Mod}(\mathcal{X}_\etale, \mathcal{O}_\mathcal{X})
\longrightarrow
\textit{Mod}(U_\etale, \mathcal{O}_U),\quad
\mathcal{F}
\longmapsto
f^*\mathcal{F}|_{U_\etale}
$$
를 생각하자. 이 함자를 국소 준연접 가군에 적용하면 함자
$$
b : \textit{LQCoh}(\mathcal{O}_\mathcal{X})
\longrightarrow
\QCoh(U_\etale, \mathcal{O}_U)
$$
를 얻는다. 『스택 위의 층』의 보조정리
\ref{stacks-sheaves-lemma-pullback-lqc}와
\ref{stacks-sheaves-lemma-compare-locally-quasi-coherent}를 보라.
함자를 더 작은 부분범주들로 한정하여
$$
c :
\textit{LQCoh}^{fbc}(\mathcal{O}_\mathcal{X})
\longrightarrow
\QCoh(U_\etale, \mathcal{O}_U)
$$
와
$$
d :
\QCoh(\mathcal{O}_\mathcal{X})
\longrightarrow
\QCoh(U_\etale, \mathcal{O}_U)
$$
를 얻을 수도 있다. 이 함자들에 관해서는 다음을 말할 수 있다.\footnote{제시한
참고문헌을 따라가기보다 냅킨 위에서 이 명제들이 왜 참인지 직접 풀어 보기를 권한다.}
\begin{enumerate}
\item 함자 $a$는 완전하다. 실제로 당김 $f^* = f^{-1}$은 완전하고
(『스택 위의 층』의 절 \ref{stacks-sheaves-section-modules}),
$U_\etale$에 대한 제한도 완전하다. 『스택 위의 층』의 식
(\ref{stacks-sheaves-equation-restrict})을 보라.
\item 함자 $b$는 완전하다. 실제로 『스택 위의 층』의 보조정리
\ref{stacks-sheaves-lemma-lqc-colimits}에 의해 포함
$\textit{LQCoh}(\mathcal{O}_\mathcal{X}) \to
\textit{Mod}(\mathcal{X}_\etale, \mathcal{O}_\mathcal{X})$는 완전하다.
\item 함자 $c$는 완전하다. 실제로 명제
\ref{stacks-cohomology-proposition-loc-qcoh-flat-base-change}에 의해 포함 함자
$\textit{LQCoh}^{fbc}(\mathcal{O}_\mathcal{X}) \to
\textit{Mod}(\mathcal{X}_\etale, \mathcal{O}_\mathcal{X})$는 완전하다.
\item 함자 $d$는 오른쪽 완전하지만 일반적으로 완전하지 않다.
실제로 『스택 위의 층』의 보조정리
\ref{stacks-sheaves-lemma-qc-colimits}에 의해 포함 함자
$\QCoh(\mathcal{O}_\mathcal{X}) \to
\textit{Mod}(\mathcal{X}_\etale, \mathcal{O}_\mathcal{X})$는 오른쪽 완전하다.
완전하지 않음을 보이는 예는 생략한다.
\item $f$가 평탄하면 $d$는 완전하다. 이는 보조정리
\ref{stacks-cohomology-lemma-flat-pullback-quasi-coherent}와 『스택 위의 층』의 보조정리
\ref{stacks-sheaves-lemma-compare-quasi-coherent}를 결합하면 따른다.
\item $f$가 평탄하면 $c$는 기생 대상을 영으로 보낸다.
실제로 보조정리 \ref{stacks-cohomology-lemma-parasitic}에 의해 $f^*$는 기생 대상을 보존한다.
그러면 $U$ 위에서 에탈이고 따라서 $\mathcal{X}$ 위에서 평탄한 임의의 스킴
$V$에 대해, 에탈 국소화와 제한의 양립성에 의해
$0 = f^*\mathcal{F}|_{V_\etale} = c(\mathcal{F})|_{V_\etale}$이다.
『스택 위의 층』의 주 \ref{stacks-sheaves-remark-compare}를 보라.
따라서 명백히 $c(\mathcal{F}) = 0$이다.
\item $f$가 평탄하면 $c = d \circ Q$이다. 실제로 보조정리
\ref{stacks-cohomology-lemma-adjoint-kernel-parasitic}에 의해
$Q(\mathcal{F}) \to \mathcal{F}$의 핵과 여핵은 기생이다.
따라서 $c$가 완전하고 (3) 기생 대상을 영으로 보내므로 (6),
$Q(\mathcal{F}) \to \mathcal{F}$에 $c$를 적용하면 동형이 된다.
\item 함자 $a, b, c, d$는 쌍대극한 및 임의의 직합과 가환한다.
$f^*$와 제한은 왼쪽 수반이므로 이에 대해 참하고, 따라서 $a$에 대해 참이다.
그러면 위에서 제시한 참고문헌들에 의해 $b$, $c$, $d$에 대해서도 따른다.
\item 함자 $a, b, c, d$는 텐서곱과 가환한다.
\item $f$가 평탄하고 전사이며,
$\mathcal{F}$가 $\textit{LQCoh}^{fbc}(\mathcal{O}_\mathcal{X})$에 속하고,
$c(\mathcal{F}) = 0$이면 $\mathcal{F}$는 기생이다. 실제로 (7)에 의해
$d(Q(\mathcal{F})) = 0$이다. 에탈 국소화와 제한의 양립성에 의해
$U$가 스킴이라고 가정해도 된다(위 참고문헌을 보라).
이제 $0 \to Q(\mathcal{F})$와 사상 $f : U \to \mathcal{X}$에
보조정리 \ref{stacks-cohomology-lemma-quasi-coherent-check-exact}를 적용하면
$Q(\mathcal{F}) = 0$임을 알 수 있다. 따라서 보조정리
\ref{stacks-cohomology-lemma-adjoint-kernel-parasitic}에 의해 $\mathcal{F}$는 기생이다.
\item $f$가 평탄하고 전사이면 함자 $d$는 완전성을 반영한다.
더 정확히, $\mathcal{F}^\bullet$을
$\QCoh(\mathcal{O}_\mathcal{X})$ 안의 복합체라 하자. 그러면
$\mathcal{F}^\bullet$이 $\QCoh(\mathcal{O}_\mathcal{X})$에서 완전인 것과
$d(\mathcal{F}^\bullet)$이 완전인 것은 동치이다. 실제로 한 방향은 (5)에서
보았다. 다른 방향을 위해 $H^i(d(\mathcal{F}^\bullet)) = 0$이라고 하자.
그러면 $\mathcal{G} = H^i(\mathcal{F}^\bullet)$은
$d(\mathcal{G}) = 0$을 만족하는
$\QCoh(\mathcal{O}_\mathcal{X})$의 대상이다. 따라서 (10)에 의해
$\mathcal{G}$는 준연접이면서 기생이고, 예를 들어 주
\ref{stacks-cohomology-remark-bousfield-colocalization}에 의해 $0$이다.
\item
\label{stacks-cohomology-item-hom-restriction}
$f$가 평탄이고,
$\mathcal{F}, \mathcal{G} \in \Ob(\QCoh(\mathcal{O}_\mathcal{X}))$이며,
$\mathcal{F}$가 유한 표시이고 두자. 그러면
$$
d(hom(\mathcal{F}, \mathcal{G})) =
\SheafHom_{\mathcal{O}_U}(d(\mathcal{F}), d(\mathcal{G}))
$$
이다. 표기법은 보조정리 \ref{stacks-cohomology-lemma-internal-hom-fp-into-qcoh}에서와 같다.
이를 보는 가장 쉬운 방법은 아마 다음과 같을 것이다.
\begin{align*}
d(hom(\mathcal{F}, \mathcal{G}))
& =
d(Q(\SheafHom_{\mathcal{O}_\mathcal{X}}(\mathcal{F}, \mathcal{G}))) \\
& =
c(\SheafHom_{\mathcal{O}_\mathcal{X}}(\mathcal{F}, \mathcal{G})) \\
& =
f^*\SheafHom_{\mathcal{O}_\mathcal{X}}(\mathcal{F}, \mathcal{G})|_{U_\etale} \\
& =
\SheafHom_{\mathcal{O}_\mathcal{U}}(f^*\mathcal{F},
f^*\mathcal{G})|_{U_\etale} \\
& =
\SheafHom_{\mathcal{O}_U}(f^*\mathcal{F}|_{U_\etale},
f^*\mathcal{G}|_{U_\etale})
\end{align*}
첫 번째 등식은 $hom$의 구성에서 따른다. 두 번째 등식은 (7)에서 따른다.
세 번째 등식은 $c$의 정의에서 따른다. 네 번째 등식은 『사이트 위의 가군』의
보조정리 \ref{sites-modules-lemma-pullback-internal-hom}에서 따른다.
마지막 등식은 『스택 위의 층』의 보조정리
\ref{stacks-sheaves-lemma-compare}의 환 달린 토포스의 평탄 사상
$i_U (U_\etale, \mathcal{O}_U) \to
(\mathcal{U}_\etale, \mathcal{O}_\mathcal{U})$에 같은 참고문헌을 적용하면 따른다.
\item 여기에 더 추가하라.
\end{enumerate}

\section{쌍대극한과 코호몰로지}
\label{stacks-cohomology-section-colimits}

\noindent
특히 다음 보조정리는 준연접 층의 도식에 적용된다.

\begin{lemma}
\label{stacks-cohomology-lemma-colimits}
$\mathcal{X}$를 준콤팩트 준분리 대수적 스택이라 하자. 그러면
$\mathcal{X}$ 위의 아벨 층들의 모든 여과 도식에 대해
$$
\colim_i H^p(\mathcal{X}, \mathcal{F}_i)
\longrightarrow
H^p(\mathcal{X}, \colim_i \mathcal{F}_i)
$$
는 동형이다. $\mathcal{X}_\etale$ 위의 아벨 층에 대해 에탈 위상에서
코호몰로지를 취하는 경우에도 마찬가지이다.
\end{lemma}

\begin{proof}
$\tau = fppf$, 각각 $\tau = \etale$라 하자. 이 보조정리는 사이트
$\mathcal{X}_\tau$에 『사이트 위의 코호몰로지』의 보조정리
\ref{sites-cohomology-lemma-colim-global}을 적용하면 따른다.
그 가정을 확인하기 위해 『사이트 위의 코호몰로지』의 주
\ref{sites-cohomology-remark-colim-global}를 사용한다.
구체적으로 $\mathcal{B} \subset \Ob(\mathcal{X}_\tau)$를 아핀 스킴 위에
놓인 대상들의 집합이라 하자. 달리 말해 $\mathcal{B}$의 원소는
$U$가 아핀인 사상 $x : U \to \mathcal{X}$이다.
그 주의 조건 (1)--(4)를 차례로 확인하자.
\begin{enumerate}
\item $\mathcal{X}$가 준콤팩트이므로 $U$가 아핀인 매끄러운 전사 사상
$x : U \to \mathcal{X}$가 존재한다(『스택의 성질』의 보조정리
\ref{stacks-properties-lemma-quasi-compact-stack}).
그러면 $h_x^\# \to *$는 $\mathcal{X}_\tau$ 위의 층의 전사 사상이다.
\item $\mathcal{X}_\tau$의 피복들은 fppf, 각각 에탈 피복이므로,
$U \in \mathcal{B}$의 모든 피복은 유한 아핀 fppf 피복에 의해 세분된다.
『위상들』의 보조정리 \ref{topologies-lemma-fppf-affine}, 각각 보조정리
\ref{topologies-lemma-etale-affine}을 보라.
\item $x : U \to \mathcal{X}$와 $x' : U' \to \mathcal{X}$가
$\mathcal{B}$에 속한다고 하자.
$\Sh(\mathcal{X}_\tau)$에서의 곱 $h_x^\# \times h_{x'}^\#$은
$\mathcal{X}$ 위의 대수공간
$W = U \times_{x, \mathcal{X}, x'} U'$로 결정되는
$\mathcal{X}_\tau$ 위의 층과 같다. 즉, $\mathcal{X}_\tau$의 대상
$y : V \to \mathcal{X}$에 대해
$(h_x^\# \times h_{x'}^\#)(y) = \{f : V \to W \mid y =
x \circ \text{pr}_1 \circ f = x' \circ \text{pr}_2 \circ f\}$이다.
$\mathcal{X}$가 준분리이므로 대수공간 $W$는 준콤팩트이다.
예를 들어 『스택의 사상』의 보조정리
\ref{stacks-morphisms-lemma-quasi-compact-quasi-separated-permanence}를 보라.
따라서 아핀 스킴 $U''$과 에탈 전사 사상 $U'' \to W$를 택할 수 있다.
$U'' \to W$와 $W \to \mathcal{X}$의 합성을
$x'' : U'' \to \mathcal{X}$라 쓰자. 그러면 원하는 대로
$h_{x''}^\# \to h_x^\# \times h_{x'}^\#$는 전사이다.
\item $x : U \to \mathcal{X}$와 $x' : U' \to \mathcal{X}$가
$\mathcal{B}$에 속한다고 하자. $a, b : U \to U'$를
$\mathcal{X}$ 위의 사상이라 하자. 즉, $a, b : x \to x'$는
$\mathcal{X}_\tau$의 사상이다. 그러면 $h_a$와 $h_b$의 등화자는
$a, b : U \to U'$의 등화자로 표현되며, 이는 $\mathcal{X}$ 위의 아핀
스킴이므로 $\mathcal{B}$에 속한다.
\end{enumerate}
이로써 증명을 끝냈다.
\end{proof}

\begin{lemma}
\label{stacks-cohomology-lemma-colimit-cohomology}
$f : \mathcal{X} \to \mathcal{Y}$를 대수적 스택의 준콤팩트 준분리 사상이라 하자.
$\mathcal{F} = \colim \mathcal{F}_i$를 $\mathcal{X}$ 위의 아벨 층들의 여과
쌍대극한이라 하자. 그러면 임의의 $p \geq 0$에 대해 다음이 성립한다.
$$
R^pf_*\mathcal{F} = \colim R^pf_*\mathcal{F}_i.
$$
$\mathcal{X}_\etale$ 위의 아벨 층에 대해 에탈 위상에서 고차
직상을 취하는 경우에도 마찬가지이다.
\end{lemma}

\begin{proof}
fppf 위상에 대해 증명하겠다. 에탈 위상에 대한 증명도 같다.
$R^if_*\mathcal{F}$는 준층
$$
(y : V \to \mathcal{Y})
\longmapsto
H^i(V \times_{y, \mathcal{Y}} \mathcal{X}, \text{pr}^{-1}\mathcal{F})
$$
에 연관된 $\mathcal{Y}_{fppf}$ 위의 층임을 상기하라.
『스택 위의 층』의 보조정리
\ref{stacks-sheaves-lemma-pushforward-restriction}을 보라.
쌍대극한은 준층 쌍대극한에 연관된 층임도 상기하라.
$V$가 아핀이면 섬유곱 $V \times_\mathcal{Y} \mathcal{X}$는
준콤팩트이고 준분리이다. 따라서 $V$가 아핀일 때
$H^p(V \times_\mathcal{Y} \mathcal{X}, -)$에 보조정리
\ref{stacks-cohomology-lemma-colimits}를 적용할 수 있다. 모든 $V$는 아핀 대상들에 의한
fppf 피복을 가지므로 보조정리가 증명된다. 일부 세부사항은 생략한다.
\end{proof}

\begin{lemma}
\label{stacks-cohomology-lemma-quasi-coherent-pushforward-direct-sums}
$f : \mathcal{X} \to \mathcal{Y}$를 대수적 스택의 준콤팩트 준분리 사상이라 하자.
함자 $f_{\QCoh, *}$와 함자 $R^if_{\QCoh, *}$는 직합 및 여과 쌍대극한과 가환한다.
\end{lemma}

\begin{proof}
보조정리 \ref{stacks-cohomology-lemma-colimit-cohomology}에 의해 함자 $f_*$와 $R^if_*$는
모든 가군에서 직합 및 여과 쌍대극한과 가환한다.
$f_{\QCoh, *} = Q \circ f_*$이고
$R^if_{\QCoh, *} = Q \circ R^if_*$이며, 보조정리
\ref{stacks-cohomology-lemma-adjoint-kernel-parasitic}에 의해 $Q$가 모든 쌍대극한과 가환하므로
결과가 따른다.
\end{proof}

\begin{lemma}
\label{stacks-cohomology-lemma-quasi-coherent-pushforward-affine}
$f : \mathcal{X} \to \mathcal{Y}$를 대수적 스택의 아핀 사상이라 하자.
함자 $R^if_{\QCoh, *}$, $i > 0$은 소멸하고, 함자 $f_{\QCoh, *}$는
완전하며 직합 및 모든 쌍대극한과 가환한다.
\end{lemma}

\begin{proof}
$R^if_{\QCoh, *} = Q \circ R^if_*$이므로 보조정리
\ref{stacks-cohomology-lemma-loc-qcoh-fbc-affine-direct-image}에서 소멸을 얻는다.
$\{R^if_{\QCoh, *}\}_{i \geq 0}$가 $\delta$-함자를 이루므로 그 소멸은
$f_{\QCoh, *}$가 완전임을 함의한다. 명제
\ref{stacks-cohomology-proposition-direct-image-quasi-coherent}를 보라.
그러면 예를 들어 보조정리 \ref{stacks-cohomology-lemma-quasi-coherent-pushforward-direct-sums}에
의해 $f_{\QCoh, *}$는 직합과 가환한다. 직합과 가환하는 완전 함자는
모든 쌍대극한과 가환한다.
\end{proof}

\noindent
다음 보조정리는 준콤팩트 준분리 대수적 스택에서 유한 표시 가군이
예상대로 행동함을 알려 준다.

\begin{lemma}
\label{stacks-cohomology-lemma-finite-presentation-quasi-compact-colimit}
$\mathcal{X}$를 준콤팩트 준분리 대수적 스택이라 하자. $I$를 유향 집합이라 하고,
$(\mathcal{F}_i, \varphi_{ii'})$를 $I$ 위의
$\mathcal{O}_\mathcal{X}$-가군들의 계라 하자. $\mathcal{G}$를 유한 표시인
$\mathcal{O}_\mathcal{X}$-가군이라 하자. 그러면 다음이 성립한다.
$$
\colim_i \Hom_\mathcal{X}(\mathcal{G}, \mathcal{F}_i)
=
\Hom_\mathcal{X}(\mathcal{G}, \colim_i \mathcal{F}_i).
$$
특히 $\Hom_\mathcal{X}(\mathcal{G}, -)$는
$\QCoh(\mathcal{O}_\mathcal{X})$에서 여과 쌍대극한과 가환한다.
\end{lemma}

\begin{proof}
표시된 등식은 『사이트 위의 가군』의 보조정리
\ref{sites-modules-lemma-finite-presentation-quasi-compact-colimit}의
특수한 경우이다. 이를 적용하려면 사이트 $\mathcal{X}_{fppf}$에 대해
『사이트』의 보조정리 \ref{sites-lemma-directed-colimits-global-sections}의
(4)의 가정을 확인해야 한다. 이를 위해 『사이트』의 주
\ref{sites-remark-stronger-conditions}의 가정 (2)(a), (2)(b), (2)(c)를
확인하겠다. 구체적으로 $\mathcal{B} \subset \Ob(\mathcal{X}_{fppf})$를
아핀 스킴 위에 놓인 대상들의 집합이라 하자. 달리 말해 $\mathcal{B}$의
원소는 $U$가 아핀인 사상 $x : U \to \mathcal{X}$이다.
그 주의 조건 (2)(a), (2)(b), (2)(c)를 차례로 확인하자.
\begin{enumerate}
\item $\mathcal{X}$가 준콤팩트이므로 $U$가 아핀인 매끄러운 전사 사상
$x : U \to \mathcal{X}$가 존재한다(『스택의 성질』의 보조정리
\ref{stacks-properties-lemma-quasi-compact-stack}).
그러면 $h_x^\# \to *$는 $\mathcal{X}_{fppf}$ 위의 층의 전사 사상이다.
\item $\mathcal{X}_{fppf}$의 피복들은 fppf 피복이므로,
$U \in \mathcal{B}$의 모든 피복은 유한 아핀 fppf 피복에 의해 세분된다.
『위상들』의 보조정리 \ref{topologies-lemma-fppf-affine}을 보라.
\item $x : U \to \mathcal{X}$와 $x' : U' \to \mathcal{X}$가
$\mathcal{B}$에 속한다고 하자.
$\Sh(\mathcal{X}_{fppf})$에서의 곱 $h_x^\# \times h_{x'}^\#$은
$\mathcal{X}$ 위의 대수공간
$W = U \times_{x, \mathcal{X}, x'} U'$로 결정되는
$\mathcal{X}_{fppf}$ 위의 층과 같다. 즉, $\mathcal{X}_{fppf}$의 대상
$y : V \to \mathcal{X}$에 대해
$(h_x^\# \times h_{x'}^\#)(y) = \{f : V \to W \mid y =
x \circ \text{pr}_1 \circ f = x' \circ \text{pr}_2 \circ f\}$이다.
$\mathcal{X}$가 준분리이므로 대수공간 $W$는 준콤팩트이다.
예를 들어 『스택의 사상』의 보조정리
\ref{stacks-morphisms-lemma-quasi-compact-quasi-separated-permanence}를 보라.
따라서 아핀 스킴 $U''$과 에탈 전사 사상 $U'' \to W$를 택할 수 있다.
$U'' \to W$와 $W \to \mathcal{X}$의 합성을
$x'' : U'' \to \mathcal{X}$라 쓰자. 그러면 원하는 대로
$h_{x''}^\# \to h_x^\# \times h_{x'}^\#$는 전사이다.
\end{enumerate}
마지막 명제를 위해 포함 함자
$\QCoh(\mathcal{O}_X) \to \textit{Mod}(\mathcal{O}_X)$가 쌍대극한과
가환하고 유한 표시 가군이 준연접임에 유의하라. 『스택 위의 층』의 보조정리
\ref{stacks-sheaves-lemma-quasi-coherent-algebraic-stack}을 보라.
\end{proof}

\section{리스-에탈 사이트와 평탄-fppf 사이트}
\label{stacks-cohomology-section-lisse-etale}

\noindent
책 \cite{LM-B}에서는 위의 결과들 가운데 많은 것을 대수적 스택의
리스-에탈 사이트를 사용하여 증명한다. 여기서 이 사이트를 정의한다.
『예』의 절 \ref{examples-section-lisse-etale-not-functorial}에서는
리스-에탈 사이트가 함자적이지 않음을 보인다.
또한 그 유사물인 평탄-fppf 사이트도 정의하는데, 이는 Stacks 프로젝트에서
제시하는 대수적 스택의 전개에 더 잘 맞는다(우리는 fppf 위상을 바탕 위상으로
사용하기 때문이다). 물론 평탄-fppf 사이트도 함자적이지 않다.

\begin{definition}
\label{stacks-cohomology-definition-lisse-etale}
$\mathcal{X}$를 대수적 스택이라 하자.
\begin{enumerate}
\item $\mathcal{X}$의 {\it 리스-에탈 사이트}는 $\mathcal{X}$의 충만한
부분범주 $\mathcal{X}_{lisse,\etale}$이다.\footnote{문헌에서는 이 사이트를
$\text{Lis-\'et}(\mathcal{X})$ 또는 $\text{Lis-Et}(\mathcal{X})$로 쓰고,
연관된 토포스를 $\mathcal{X}_{\text{lis-\'e}t}$ 또는
$\mathcal{X}_{\text{lis-et}}$로 쓴다. Stacks 프로젝트의 규약은 사이트에
이름을 붙이고 대응하는 토포스를 $\Sh(\mathcal{C})$로 나타내는 것이다.}
그 대상은 스킴 $U$ 위에 놓인 $x \in \Ob(\mathcal{X})$ 가운데
$x : U \to \mathcal{X}$가 매끄러운 것들이다.
$\mathcal{X}_{lisse,\etale}$의 피복은
$\mathcal{X}_\etale$의 피복을 이루는
$\mathcal{X}_{lisse,\etale}$의 사상들의 족
$\{x_i \to x\}_{i \in I}$이다.
\item $\mathcal{X}$의 {\it 평탄-fppf 사이트}는 $\mathcal{X}$의 충만한
부분범주 $\mathcal{X}_{flat,fppf}$이다. 그 대상은 스킴 $U$ 위에 놓인
$x \in \Ob(\mathcal{X})$ 가운데 $x : U \to \mathcal{X}$가 평탄한 것들이다.
$\mathcal{X}_{flat,fppf}$의 피복은 $\mathcal{X}_{fppf}$의 피복을 이루는
$\mathcal{X}_{flat,fppf}$의 사상들의 족
$\{x_i \to x\}_{i \in I}$이다.
\end{enumerate}
\end{definition}

\noindent
$\mathcal{O}_{\mathcal{X}_{lisse,\etale}}$로
$\mathcal{O}_\mathcal{X}$를 리스-에탈 사이트에 제한한 것을 나타내며,
$\mathcal{O}_{\mathcal{X}_{flat,fppf}}$도 마찬가지이다.
리스-에탈 사이트와 에탈 사이트의 관계는 다음과 같다
(이 보조정리에서는 주로 ``위상적'' 성질들만 다룬다).

\begin{lemma}
\label{stacks-cohomology-lemma-lisse-etale}
$\mathcal{X}$를 대수적 스택이라 하자.
\begin{enumerate}
\item 포함 함자
$\mathcal{X}_{lisse,\etale} \to \mathcal{X}_\etale$는 완전충실하고,
연속이며, 쌍연속이다. 따라서 다음이 성립한다.
\begin{enumerate}
\item 토포스의 사상
$$
g :
\Sh(\mathcal{X}_{lisse,\etale})
\longrightarrow
\Sh(\mathcal{X}_\etale)
$$
이 있고, $g^{-1}$은 제한으로 주어진다.
\item 집합의 층에서 함자 $g^{-1}$은 왼쪽 수반 $g_!^{Sh}$를 가진다.
\item 수반 사상 $g^{-1}g_* \to \text{id}$와
$\text{id} \to g^{-1}g_!^{Sh}$는 동형이다.
\item 아벨 층에서 함자 $g^{-1}$은 왼쪽 수반 $g_!$을 가진다.
\item 수반 사상 $\text{id} \to g^{-1}g_!$은 동형이다.
\item $g^{-1}\mathcal{O}_\mathcal{X} =
\mathcal{O}_{\mathcal{X}_{lisse,\etale}}$이다. 따라서 $g$는
$g^{-1} = g^*$인 환 달린 토포스의 평탄 사상을 유도한다.
\end{enumerate}
\item 포함 함자
$\mathcal{X}_{flat,fppf} \to \mathcal{X}_{fppf}$는 완전충실하고,
연속이며, 쌍연속이다. 따라서 다음이 성립한다.
\begin{enumerate}
\item 토포스의 사상
$$
g :
\Sh(\mathcal{X}_{flat,fppf})
\longrightarrow
\Sh(\mathcal{X}_{fppf})
$$
이 있고, $g^{-1}$은 제한으로 주어진다.
\item 집합의 층에서 함자 $g^{-1}$은 왼쪽 수반 $g_!^{Sh}$를 가진다.
\item 수반 사상 $g^{-1}g_* \to \text{id}$와
$\text{id} \to g^{-1}g_!^{Sh}$는 동형이다.
\item 아벨 층에서 함자 $g^{-1}$은 왼쪽 수반 $g_!$을 가진다.
\item 수반 사상 $\text{id} \to g^{-1}g_!$은 동형이다.
\item $g^{-1}\mathcal{O}_\mathcal{X} =
\mathcal{O}_{\mathcal{X}_{flat,fppf}}$이다. 따라서 $g$는
$g^{-1} = g^*$인 환 달린 토포스의 평탄 사상을 유도한다.
\end{enumerate}
\end{enumerate}
\end{lemma}

\begin{proof}
두 경우 모두 함자가 완전충실하고 연속이며 쌍연속임은 즉시 알 수 있다
(『사이트』의 정의 \ref{sites-definition-continuous}와
\ref{sites-definition-cocontinuous}를 보라).
따라서 성질 (a), (b), (c)는 『사이트』의 보조정리
\ref{sites-lemma-when-shriek}와 \ref{sites-lemma-back-and-forth}에서 따른다.
(d), (e)는 『사이트 위의 가군』의 보조정리
\ref{sites-modules-lemma-g-shriek-adjoint}와
\ref{sites-modules-lemma-back-and-forth}에서 따른다.
(f)는 즉시 성립한다.
\end{proof}

\begin{lemma}
\label{stacks-cohomology-lemma-lisse-etale-cohomology}
$\mathcal{X}$를 대수적 스택이라 하자. 표기법은 보조정리
\ref{stacks-cohomology-lemma-lisse-etale}에서와 같다.
\begin{enumerate}
\item $\mathcal{X}_\etale$ 위의 아벨 층 $\mathcal{F}$에 대해 다음이 성립한다.
\begin{enumerate}
\item $H^p(\mathcal{X}_\etale, \mathcal{F}) =
H^p(\mathcal{X}_{lisse,\etale}, g^{-1}\mathcal{F})$이다.
\item $\mathcal{X}_{lisse,\etale}$의 임의의 대상 $x$에 대해
$H^p(x, \mathcal{F}) =
H^p(\mathcal{X}_{lisse,\etale}/x, g^{-1}\mathcal{F})$이다.
\end{enumerate}
가군층에 대해서도 마찬가지이다.
\item $\mathcal{X}_{fppf}$ 위의 아벨 층 $\mathcal{F}$에 대해 다음이 성립한다.
\begin{enumerate}
\item $H^p(\mathcal{X}_{fppf}, \mathcal{F}) =
H^p(\mathcal{X}_{flat,fppf}, g^{-1}\mathcal{F})$이다.
\item $\mathcal{X}_{flat,fppf}$의 임의의 대상 $x$에 대해
$H^p(x, \mathcal{F}) =
H^p(\mathcal{X}_{flat,fppf}/x, g^{-1}\mathcal{F})$이다.
\end{enumerate}
가군층에 대해서도 마찬가지이다.
\end{enumerate}
\end{lemma}

\begin{proof}
(1)(a)는 포함 함자
$\mathcal{X}_{lisse,\etale} \to \mathcal{X}_\etale$에
『스택 위의 층』의 보조정리
\ref{stacks-sheaves-lemma-cohomology-on-subcategory}를 적용하면 따른다.
(1)(b)는 (1)(a)에서 따른다. 구체적으로 $x$가 스킴 $U$ 위에 놓이면
사이트 $\mathcal{X}_\etale/x$는 $(\Sch/U)_\etale$와 동치이고,
$\mathcal{X}_{lisse,\etale}$는 $U_{lisse,\etale}$와 동치이다.
(2)도 같은 방식으로 증명한다.
\end{proof}

\begin{lemma}
\label{stacks-cohomology-lemma-lisse-etale-modules}
$\mathcal{X}$를 대수적 스택이라 하자. 표기법은 보조정리
\ref{stacks-cohomology-lemma-lisse-etale}에서와 같다.
\begin{enumerate}
\item $g^*$의 왼쪽 수반인 함자
$$
g_! :
\textit{Mod}(\mathcal{X}_{lisse,\etale},
\mathcal{O}_{\mathcal{X}_{lisse,\etale}})
\longrightarrow
\textit{Mod}(\mathcal{X}_\etale, \mathcal{O}_{\mathcal{X}})
$$
이 존재한다. 더욱이 이는 아벨 층에서의 함자 $g_!$과 일치하고
$g^*g_! = \text{id}$이다.
\item $g^*$의 왼쪽 수반인 함자
$$
g_! :
\textit{Mod}(\mathcal{X}_{flat,fppf},
\mathcal{O}_{\mathcal{X}_{flat,fppf}})
\longrightarrow
\textit{Mod}(\mathcal{X}_{fppf}, \mathcal{O}_{\mathcal{X}})
$$
이 존재한다. 더욱이 이는 아벨 층에서의 함자 $g_!$과 일치하고
$g^*g_! = \text{id}$이다.
\end{enumerate}
\end{lemma}

\begin{proof}
두 경우 모두 함자 $g_!$의 존재는 『사이트 위의 가군』의 보조정리
\ref{sites-modules-lemma-lower-shriek-modules}에서 따른다.
$g_!$이 아벨 층에서의 함자와 일치함을 보기 위해 『사이트 위의 가군』의 식
(\ref{sites-modules-equation-compare-on-localizations})의 사상들이
동형임을 보이겠다.

\medskip\noindent
리스-에탈 경우. $x \in \Ob(\mathcal{X}_{lisse,\etale})$가 스킴 $U$ 위에
놓이고 $x : U \to \mathcal{X}$가 매끄럽다고 하자. 유도된 완전충실 함자
$$
g' :
\mathcal{X}_{lisse,\etale}/x
\longrightarrow
\mathcal{X}_\etale/x
$$
를 생각하자. 오른쪽은 $(\Sch/U)_\etale$와 동일시되고,
왼쪽은 합성 $U' \to U \to \mathcal{X}$가 매끄러운 스킴 $U'/U$들의
충만한 부분범주와 동일시된다. 따라서 『에탈 코호몰로지』의 보조정리
\ref{etale-cohomology-lemma-compare-structure-sheaves}를 적용할 수 있다.

\medskip\noindent
평탄-fppf 경우. $x \in \Ob(\mathcal{X}_{flat,fppf})$가 스킴 $U$ 위에
놓이고 $x : U \to \mathcal{X}$가 평탄하다고 하자. 유도된 완전충실 함자
$$
g' :
\mathcal{X}_{flat,fppf}/x
\longrightarrow
\mathcal{X}_{fppf}/x
$$
를 생각하자. 오른쪽은 $(\Sch/U)_{fppf}$와 동일시되고,
왼쪽은 합성 $U' \to U \to \mathcal{X}$가 평탄한 스킴 $U'/U$들의
충만한 부분범주와 동일시된다. 따라서 『에탈 코호몰로지』의 보조정리
\ref{etale-cohomology-lemma-compare-structure-sheaves}를 적용할 수 있다.

\medskip\noindent
두 경우 모두 등식 $g^*g_! = \text{id}$는 $g^* = g^{-1}$과
보조정리 \ref{stacks-cohomology-lemma-lisse-etale}의 아벨 층에 대한 등식에서 따른다.
\end{proof}

\begin{lemma}
\label{stacks-cohomology-lemma-lisse-etale-structure-sheaf}
$\mathcal{X}$를 대수적 스택이라 하자. 표기법은 보조정리
\ref{stacks-cohomology-lemma-lisse-etale}와 \ref{stacks-cohomology-lemma-lisse-etale-modules}에서와 같다.
\begin{enumerate}
\item $g_!\mathcal{O}_{\mathcal{X}_{lisse,\etale}} =
\mathcal{O}_\mathcal{X}$이다.
\item $g_!\mathcal{O}_{\mathcal{X}_{flat, fppf}} =
\mathcal{O}_\mathcal{X}$이다.
\end{enumerate}
\end{lemma}

\begin{proof}
이 증명에서는
$\mathcal{C} = \mathcal{X}_\etale$
(각각 $\mathcal{C} = \mathcal{X}_{fppf}$)라 쓰고,
$\mathcal{C}' = \mathcal{X}_{lisse,\etale}$
(각각 $\mathcal{C}' = \mathcal{X}_{flat, fppf}$)라 쓰겠다.
그러면 $\mathcal{C}'$은 $\mathcal{C}$의 충만한 부분범주이다.
이 증명에서는 $\mathcal{C}$의 대상 $V$를 $\mathcal{X}$ 위의 스킴으로,
$\mathcal{C}'$의 대상 $U$를 $\mathcal{X}$ 위에서 매끄러운
(각각 평탄한) 스킴으로 생각하겠다. 마지막으로
$\mathcal{O} = \mathcal{O}_\mathcal{X}$라 쓰고,
$\mathcal{O}' = \mathcal{O}_{\mathcal{X}_{lisse,\etale}}$
(각각 $\mathcal{O}' = \mathcal{O}_{\mathcal{X}_{flat,fppf}}$)라 쓰겠다.
위 표기법으로 $\mathcal{O}(V) = \Gamma(V, \mathcal{O}_V)$이고
$\mathcal{O}'(U) = \Gamma(U, \mathcal{O}_U)$이다.
동일시 $\mathcal{O}' = g^{-1}\mathcal{O}$의 수반인
$\mathcal{O}$-가군 준동형 $g_!\mathcal{O}' \to \mathcal{O}$를 생각하자.

\medskip\noindent
$g_!\mathcal{O}'$은 규칙
$$
V \longmapsto \colim_{V \to U} \mathcal{O}'(U)
$$
으로 주어지는 준층 $g_{p!}\mathcal{O}'$에 연관된 층임을 상기하라.
여기서 쌍대극한은 아벨 군의 범주에서 취한다(『사이트 위의 가군』의 정의
\ref{sites-modules-definition-g-shriek}). 아래에서는 사상
$$
V \to U \to U'
$$
이 있고 $f' \in \mathcal{O}'(U')$가 $f \in \mathcal{O}'(U)$로 제한되면,
$(V \to U, f)$와 $(V \to U', f')$가 쌍대극한의 같은 원소를 정의한다는 사실을
자주 사용하겠다. 또한 $g_!\mathcal{O}' \to \mathcal{O}$는 원소
$(V \to U, f)$를 단순히 $f$를 $V$로 당긴 것으로 보낸다.

\medskip\noindent
$g_!\mathcal{O}' \to \mathcal{O}$가 전사임을 증명하자.
$\mathcal{C}$의 어떤 대상 $V$에 대해 $h \in \mathcal{O}(V)$라 하자.
$h$가 국소적으로 상에 속함을 보이면 충분하다. 매끄러운 전사 사상
$U \to \mathcal{X}$에 대응하는 $\mathcal{C}'$의 대상 $U$를 택한다.
$U \times_\mathcal{X} V \to V$가 매끄러운 전사이므로, $V$를 $V$의 어떤
에탈 피복의 원소들로 바꾸면 사상 $V \to U$가 존재한다고 가정해도 된다.
『공간 위의 위상들』의 보조정리
\ref{spaces-topologies-lemma-etale-dominates-smooth}를 보라.
$h$를 사용하여 사상 $V \to U \times \mathbf{A}^1$을 얻는데,
$\mathbf{A}^1 = \Spec(\mathbf{Z}[t])$라 쓰면 원소
$t \in \mathcal{O}(U \times \mathbf{A}^1)$가 $h$로 당겨진다.
$U \times \mathbf{A}^1$은 $\mathcal{C}'$의 대상이므로
$(V \to U \times \mathbf{A}^1, t)$는 위 쌍대극한의 원소이고,
원하는 대로 $h \in \mathcal{O}(V)$로 간다.

\medskip\noindent
$s \in g_!\mathcal{O}'(V)$가 $\mathcal{O}(V)$의 영으로 가는 단면이라고 하자.
증명을 마치려면 $s$가 영임을 보여야 한다. $V$를 어떤 피복의 원소들로
바꾸어 $s$가 쌍대극한
$$
\colim_{V \to U} \mathcal{O}'(U)
$$
의 원소라고 가정해도 된다. $s = \sum (\varphi_i, s_i)$가 유한합이고,
$\varphi_i : V \to U_i$, $U_i$는 $\mathcal{X}$ 위에서 매끄러우며
(각각 평탄하며), $s_i \in \Gamma(U_i, \mathcal{O}_{U_i})$라고 하자.
대수공간 $U = U_1 \times_\mathcal{X} \ldots \times_\mathcal{X} U_n$ 위의
에탈 전사 스킴 $W$를 택한다. $W$는 여전히 $\mathcal{X}$ 위에서
매끄럽고(각각 평탄하고), 즉 $\mathcal{C}'$의 대상을 정의함에 유의하라.
섬유곱
$$
V' = V \times_{(\varphi_1, \ldots, \varphi_n), U} W
$$
은 $V$ 위에서 에탈 전사이므로, $s$가 $g_!\mathcal{O}'(V')$의 영으로
감을 보이면 충분하다. 제한 $\sum (\varphi_i, s_i)|_{V'}$는 함수 $s_i$들을
$W$로 당긴 것들의 합에 대응함에 유의하라. 달리 말해,
$\varphi : V \to U$가 $U$가 $\mathcal{C}'$에 속하는 사상이고
$s \in \mathcal{O}'(U)$가 $\mathcal{O}(V)$의 영으로 제한되는
$(\varphi, s)$의 경우로 환원했다. 가환 도식
$$
\xymatrix{
V \ar[rr]_-{(\varphi, 0)} \ar[rrd]_\varphi & & U \times \mathbf{A}^1 \\
& & U \ar[u]_{(\text{id}, 0)}
}
$$
에 의해
$((\varphi, 0) : V \to U \times \mathbf{A}^1, \text{pr}_2^*x)$는
위 쌍대극한의 영을 나타낸다. 따라서 $U$를 $U \times \mathbf{A}^1$로,
$\varphi$를 $(\varphi, 0)$으로, $s$를
$\text{pr}_1^*s + \text{pr}_2^*x$로 바꾸어도 된다. 그러므로 $U$ 안에서
$s$의 영점 자취 $Z : s = 0$이 $\mathcal{X}$ 위에서 매끄럽다고
(각각 평탄하다고) 가정해도 된다. 이제 $(V \to Z, 0)$과
$(\varphi, s)$가 쌍대극한에서 같은 값을 가짐을 알 수 있다. 즉, 원하는 대로
원소 $s$가 영임을 알 수 있다.
\end{proof}

\noindent
리스-에탈 사이트와 평탄-fppf 사이트를 사용하여 기생 가군을 다음과 같이
특성화할 수 있다.

\begin{lemma}
\label{stacks-cohomology-lemma-parasitic-in-terms-flat-fppf}
$\mathcal{X}$를 대수적 스택이라 하자.
\begin{enumerate}
\item $\mathcal{F}$를 $\mathcal{X}_\etale$ 위에서 평탄 밑변환 성질을 가진
$\mathcal{O}_\mathcal{X}$-가군이라 하자. 다음은 동치이다.
\begin{enumerate}
\item $\mathcal{F}$는 기생이다.
\item 보조정리 \ref{stacks-cohomology-lemma-lisse-etale}에서와 같은
$g : \Sh(\mathcal{X}_{lisse,\etale}) \to
\Sh(\mathcal{X}_\etale)$에 대해 $g^*\mathcal{F} = 0$이다.
\end{enumerate}
\item $\mathcal{F}$를 $\mathcal{X}_{fppf}$ 위의
$\mathcal{O}_\mathcal{X}$-가군이라 하자. 다음은 동치이다.
\begin{enumerate}
\item $\mathcal{F}$는 기생이다.
\item 보조정리 \ref{stacks-cohomology-lemma-lisse-etale}에서와 같은
$g :  \Sh(\mathcal{X}_{flat,fppf}) \to \Sh(\mathcal{X}_{fppf})$에 대해
$g^*\mathcal{F} = 0$이다.
\end{enumerate}
\end{enumerate}
\end{lemma}

\begin{proof}
(2)는 정의에서 즉시 따른다(이는 리스-에탈 사이트에 비해 평탄-fppf
사이트가 가지는 장점 가운데 하나이다). 함의 (1)(a) $\Rightarrow$ (1)(b)도
즉시 성립한다. (1)(b) $\Rightarrow$ (1)(a)를 보기 위해 $U$를 스킴이라 하고
$x : U \to \mathcal{X}$를 매끄러운 전사 사상이라 하자.
그러면 $x$는 $\mathcal{X}$의 리스-에탈 사이트의 대상이다. 따라서 (1)(b)는
$\mathcal{F}|_{U_\etale} = 0$을 함의한다. $V$가 스킴인 평탄 사상
$V \to \mathcal{X}$를 택하자. $W = U \times_\mathcal{X} V$라 놓고 도식
$$
\xymatrix{
W \ar[d]_p \ar[r]_q & V \ar[d] \\
U \ar[r] & \mathcal{X}
}
$$
을 생각하자. 사영 $p : W \to U$는 평탄하고 사영 $q : W \to V$는
매끄럽고 전사이다. 따라서 $q_{small}^*$는 준연접 가군에서 충실한 함자이다.
가정에 따라 $\mathcal{F}$는 평탄 밑변환 성질을 가지므로
$p_{small}^*\mathcal{F}|_{U_\etale} \cong
q_{small}^*\mathcal{F}|_{V_\etale}$를 얻는다.
그러므로 $\mathcal{F}$가 $g^*$의 핵에 속하면
$\mathcal{F}|_{V_\etale} = 0$이며, 이것이 원하는 바이다.
\end{proof}

\section{리스-에탈 사이트와 평탄-fppf 사이트의 함자성}
\label{stacks-cohomology-section-lisse-etale-functorial}

\noindent
리스-에탈 사이트는 대수적 스택의 매끄러운 사상에 대해 함자적이고,
평탄-fppf 사이트는 대수적 스택의 평탄 사상에 대해 함자적이다.
리스-에탈 및 평탄-fppf 토포스는 대수적 스택의 모든 사상에 대해 함자적인
것이 아님을 독자에게 경고한다. 『예』의 절
\ref{examples-section-lisse-etale-not-functorial}을 보라.

\begin{lemma}
\label{stacks-cohomology-lemma-lisse-etale-functorial}
$f : \mathcal{X} \to \mathcal{Y}$를 대수적 스택의 사상이라 하자.
\begin{enumerate}
\item $f$가 매끄러우면 $f$는 연속 쌍연속 함자
$\mathcal{X}_{lisse,\etale} \to \mathcal{Y}_{lisse,\etale}$로 제한되며,
다음 가환 도식에 들어가는 환 달린 토포스의 사상을 준다.
$$
\xymatrix{
\Sh(\mathcal{X}_{lisse,\etale}) \ar[r]_{g'} \ar[d]_{f'} &
\Sh(\mathcal{X}_\etale) \ar[d]^f \\
\Sh(\mathcal{Y}_{lisse,\etale}) \ar[r]^g &
\Sh(\mathcal{Y}_\etale)
}
$$
$f'_*(g')^{-1} = g^{-1}f_*$이고 $g'_!(f')^{-1} = f^{-1}g_!$이다.
\item $f$가 평탄하면 $f$는 연속 쌍연속 함자
$\mathcal{X}_{flat,fppf} \to \mathcal{Y}_{flat,fppf}$로 제한되며,
다음 가환 도식에 들어가는 환 달린 토포스의 사상을 준다.
$$
\xymatrix{
\Sh(\mathcal{X}_{flat,fppf}) \ar[r]_{g'} \ar[d]_{f'} &
\Sh(\mathcal{X}_{fppf}) \ar[d]^f \\
\Sh(\mathcal{Y}_{flat,fppf}) \ar[r]^g &
\Sh(\mathcal{Y}_{fppf})
}
$$
$f'_*(g')^{-1} = g^{-1}f_*$이고 $g'_!(f')^{-1} = f^{-1}g_!$이다.
\end{enumerate}
\end{lemma}

\begin{proof}
첫 명제는 다음 사실에서 온다. $x \in \Ob(\mathcal{X})$가 스킴 $U$ 위에
놓이고 $x : U \to \mathcal{X}$가 매끄러우며(각각 평탄하며), $f$가
매끄러우면(각각 평탄하면), $f(x) : U \to \mathcal{Y}$는 매끄럽다
(각각 평탄하다). 『스택의 사상』의 보조정리
\ref{stacks-morphisms-lemma-composition-smooth}와
\ref{stacks-morphisms-lemma-composition-flat}을 보라.
유도된 함자
$\mathcal{X}_{lisse,\etale} \to \mathcal{Y}_{lisse,\etale}$
(각각 $\mathcal{X}_{flat,fppf} \to \mathcal{Y}_{flat,fppf}$)는 이 범주들에서의
피복에 대한 우리의 정의에 의해 연속이고 쌍연속이다.
마지막으로 수평 사상들이 포함 함자들로 주어진다는 사실
(보조정리 \ref{stacks-cohomology-lemma-lisse-etale}를 보라)과 『사이트』의 보조정리
\ref{sites-lemma-composition-cocontinuous}에 의해 도식은 가환한다.

\medskip\noindent
$f'_*(g')^{-1} = g^{-1}f_*$임을 보이기 위해 $\mathcal{F}$를
$\mathcal{X}_\etale$ 위의 층(각각 $\mathcal{X}_{fppf}$ 위의 층)이라 하자.
표준 당김 사상
$$
g^{-1}f_*\mathcal{F} \longrightarrow f'_*(g')^{-1}\mathcal{F}
$$
이 있다. 『사이트』의 절 \ref{sites-section-pullback}을 보라.
이 사상이 동형임을 주장한다. 이를 증명하기 위해
$\mathcal{Y}_{lisse,\etale}$의 대상(각각 $\mathcal{Y}_{flat,fppf}$의 대상)
$y$를 택하자. $y$가 스킴 $V$ 위에 놓이고
$y : V \to \mathcal{Y}$가 매끄럽다고(각각 평탄하다고) 하자.
$g^{-1}$은 제한이므로 『스택 위의 층』의 식
(\ref{stacks-sheaves-equation-pushforward})에 의해
$$
\left(g^{-1}f_*\mathcal{F}\right)(y) =
\Gamma(V \times_{y, \mathcal{Y}} \mathcal{X},\ \text{pr}^{-1}\mathcal{F})
$$
이다. $(V \times_{y, \mathcal{Y}} \mathcal{X})' \subset
V \times_{y, \mathcal{Y}} \mathcal{X}$를, 유도된 사상
$W \to \mathcal{X}$가 매끄러운(각각 평탄한) 대상
$z : W \to V \times_{y, \mathcal{Y}} \mathcal{X}$들로 이루어진 충만한
부분범주라 하자. 위 식에 사용된 함자 $\text{pr}$의 제한을
$$
\text{pr}' :
(V \times_{y, \mathcal{Y}} \mathcal{X})'
\longrightarrow
\mathcal{X}_{lisse,\etale}
\ (\text{각각 }\mathcal{X}_{flat,fppf})
$$
로 쓰자. 『스택 위의 층』의 식
(\ref{stacks-sheaves-equation-pushforward})을 증명하는 것과 정확히 같은
논증에 의해, $\mathcal{X}_{lisse,\etale}$ 위의 임의의 층 $\mathcal{H}$
(각각 $\mathcal{X}_{flat,fppf}$ 위의 임의의 층)에 대해
\begin{equation}
\label{stacks-cohomology-equation-pushforward-lisse-etale}
f'_*\mathcal{H}(y) =
\Gamma((V \times_{y, \mathcal{Y}} \mathcal{X})',
\ (\text{pr}')^{-1}\mathcal{H})
\end{equation}
이다. $(g')^{-1}$은 제한이므로
$$
\left(f'_*(g')^{-1}\mathcal{F}\right)(y) =
\Gamma((V \times_{y, \mathcal{Y}} \mathcal{X})',
\ \text{pr}^{-1}\mathcal{F}|_{(V \times_{y, \mathcal{Y}} \mathcal{X})'})
$$
이다. 『스택 위의 층』의 보조정리
\ref{stacks-sheaves-lemma-cohomology-on-subcategory}에 의해
$$
\Gamma((V \times_{y, \mathcal{Y}} \mathcal{X})',
\ \text{pr}^{-1}\mathcal{F}|_{(V \times_{y, \mathcal{Y}} \mathcal{X})'})
=
\Gamma(V \times_{y, \mathcal{Y}} \mathcal{X},\ \text{pr}^{-1}\mathcal{F})
$$
로, 원하는 대로 같다. 보조정리의 가정에 대한 확인은 생략하지만,
$V \to \mathcal{Y}$가 매끄럽다는(각각 평탄하다는) 사실이 두 번째 조건을
확인하는 데 사용됨을 지적한다.

\medskip\noindent
마지막으로 등식 $g'_!(f')^{-1} = f^{-1}g_!$은
$f^{-1}$과 $f_*$의 수반성, $g_!$과 $g^{-1}$의 수반성 및 각각의
``프라임 붙은'' 판본에 의해 등식 $f'_*(g')^{-1} = g^{-1}f_*$에서
형식적으로 따른다.
\end{proof}

\begin{lemma}
\label{stacks-cohomology-lemma-lisse-etale-functorial-pushforward}
보조정리 \ref{stacks-cohomology-lemma-lisse-etale-functorial}과 같은 가정과 표기법을 사용하자.
$\mathcal{H}$를 $\mathcal{X}_{lisse,\etale}$ 위의 아벨 층
(각각 $\mathcal{X}_{flat,fppf}$ 위의 아벨 층)이라 하자. 그러면
\begin{equation}
\label{stacks-cohomology-equation-higher-direct-image-lisse-etale}
R^pf'_*\mathcal{H} =
\text{다음 준층에 연관된 층 }y \longmapsto
H^p((V \times_{y, \mathcal{Y}} \mathcal{X})', (\text{pr}')^{-1}\mathcal{H})
\end{equation}
이다. 여기서 $y$는 스킴 $V$ 위에 놓인
$\mathcal{Y}_{lisse,\etale}$의 대상(각각
$\mathcal{Y}_{flat,fppf}$의 대상)이고,
표기법 $(V \times_{y, \mathcal{Y}} \mathcal{X})'$와 $\text{pr}'$은
증명에서 설명한다.
\end{lemma}

\begin{proof}
보조정리 \ref{stacks-cohomology-lemma-lisse-etale-functorial}의 증명에서처럼
$(V \times_{y, \mathcal{Y}} \mathcal{X})' \subset
V \times_{y, \mathcal{Y}} \mathcal{X}$를 다음 대상 $(x, \varphi)$들로
이루어진 충만한 부분범주라 하자. $x$는
$\mathcal{X}_{lisse,\etale}$의 대상(각각
$\mathcal{X}_{flat,fppf}$의 대상)이고,
$\varphi : f(x) \to y$는 $\mathcal{Y}$의 사상이다.
식 (\ref{stacks-cohomology-equation-pushforward-lisse-etale})에 의해
$$
f'_*\mathcal{H}(y) =
\Gamma((V \times_{y, \mathcal{Y}} \mathcal{X})',
\ (\text{pr}')^{-1}\mathcal{H})
$$
이다. 여기서 $\text{pr}'$은 사영이다.
$(V \times_{y, \mathcal{Y}} \mathcal{X})'$의 대상 $(x, \varphi)$에 대해
$\varphi$를 $x$ 위의 $(f')^{-1}h_y$의 단면으로 생각할 수 있다.
따라서 $(V \times_\mathcal{Y} \mathcal{X})'$은 집합의 층
$(f')^{-1}h_y$에서 사이트 $\mathcal{X}_{lisse,\etale}$
(각각 $\mathcal{X}_{flat,fppf}$)을 국소화한 것이다.
『사이트』의 보조정리 \ref{sites-lemma-localize-topos-site}를 보라.
사상
$$
\text{pr}' : (V \times_{y, \mathcal{Y}} \mathcal{X})'
\to \mathcal{X}_{lisse,\etale}
\ (\text{각각 }
\text{pr}' : (V \times_{y, \mathcal{Y}} \mathcal{X})'
\to \mathcal{X}_{flat,fppf})
$$
은 국소화 사상이다. 특히 당김 $(\text{pr}')^{-1}$은 단사 아벨 층을
보존한다. 『사이트 위의 코호몰로지』의 보조정리
\ref{sites-cohomology-lemma-cohomology-on-sheaf-sets}를 보라.

\medskip\noindent
$\mathcal{X}_{lisse,\etale}$ 위에서(각각
$\mathcal{X}_{flat,fppf}$ 위에서) 단사 분해
$\mathcal{H} \to \mathcal{I}^\bullet$을 택한다.
직상 공식에 의해 $R^if'_*\mathcal{H}$는 $y$에 복합체
$$
\begin{matrix}
\Gamma\Big((V \times_{y, \mathcal{Y}} \mathcal{X})',
(\text{pr}')^{-1}\mathcal{I}^{i - 1}\Big) \\
\downarrow \\
\Gamma\Big((V \times_{y, \mathcal{Y}} \mathcal{X})',
(\text{pr}')^{-1}\mathcal{I}^i\Big) \\
\downarrow \\
\Gamma\Big((V \times_{y, \mathcal{Y}} \mathcal{X})',
(\text{pr}')^{-1}\mathcal{I}^{i + 1}\Big)
\end{matrix}
$$
의 코호몰로지를 대응시키는 준층에 연관된 층이다.
$(\text{pr}')^{-1}$은 완전하고 단사 대상을 보존하므로
복합체 $(\text{pr}')^{-1}\mathcal{I}^\bullet$은
$(\text{pr}')^{-1}\mathcal{H}$의 단사 분해이다. 이로써 증명이 끝난다.
\end{proof}

\begin{lemma}
\label{stacks-cohomology-lemma-lisse-etale-functorial-cohomology}
보조정리 \ref{stacks-cohomology-lemma-lisse-etale-functorial}과 같은 가정과 표기법을 사용하자.
$\mathcal{X}_\etale$ 위의 임의의 아벨 층 $\mathcal{F}$
(각각 $\mathcal{X}_{fppf}$ 위의 임의의 아벨 층)에 대해 표준 (밑변환) 사상
$$
g^{-1}Rf_*\mathcal{F} \longrightarrow Rf'_*(g')^{-1}\mathcal{F}
$$
은 동형이다.
\end{lemma}

\begin{proof}
『스택 위의 층』의 보조정리
\ref{stacks-sheaves-lemma-pushforward-restriction}과 보조정리
\ref{stacks-cohomology-lemma-lisse-etale-functorial-pushforward}에서 주어진
$g^{-1}R^pf_*\mathcal{F}$와 $R^pf'_*(g')^{-1}\mathcal{F}$의 공식을
비교하면
$$
H^p((V \times_{y, \mathcal{Y}} \mathcal{X})',
\ \text{pr}^{-1}\mathcal{F}|_{(V \times_{y, \mathcal{Y}} \mathcal{X})'})
=
H^p_\tau(V \times_{y, \mathcal{Y}} \mathcal{X},\ \text{pr}^{-1}\mathcal{F})
$$
임을 보이면 충분하다. 여기서 $\tau = \etale$
(각각 $\tau = fppf$)이다. 또 $y$는 스킴 $V$ 위에 놓인
$\mathcal{Y}$의 대상이고 사상 $y : V \to \mathcal{Y}$는 매끄럽다
(각각 평탄하다). 이 등식은 『스택 위의 층』의 보조정리
\ref{stacks-sheaves-lemma-cohomology-on-subcategory}에서 따른다.
그 보조정리의 가정 확인은 생략하지만, $V \to \mathcal{Y}$가 매끄럽다는
(각각 평탄하다는) 사실이 두 번째 조건을 확인하는 데 사용됨을 지적한다.
\end{proof}

\section{준연접 가군과 리스-에탈 및 평탄-fppf 사이트}
\label{stacks-cohomology-section-quasi-coherent-modules-II}

\noindent
이 절에서는 대수적 스택 위의 준연접 가군을 그 리스-에탈 또는 평탄-fppf
사이트의 관점에서 생각하는 방법을 설명한다.

\begin{lemma}
\label{stacks-cohomology-lemma-check-qc-on-etale-covering}
$\mathcal{X}$를 대수적 스택이라 하자.
\begin{enumerate}
\item $f_j : \mathcal{X}_j \to \mathcal{X}$를
$|\mathcal{X}| =\bigcup |f_j|(|\mathcal{X}_j|)$를 만족하는
대수적 스택의 매끄러운 사상들의 족이라 하자.
$\mathcal{F}$를 $\mathcal{X}_\etale$ 위의
$\mathcal{O}_\mathcal{X}$-가군층이라 하자.
각 $f_j^{-1}\mathcal{F}$가 준연접이면 $\mathcal{F}$도 준연접이다.
\item $f_j : \mathcal{X}_j \to \mathcal{X}$를
$|\mathcal{X}| =\bigcup |f_j|(|\mathcal{X}_j|)$를 만족하는
대수적 스택의 평탄하고 국소 유한 표시인 사상들의 족이라 하자.
$\mathcal{F}$를 $\mathcal{X}_{fppf}$ 위의
$\mathcal{O}_\mathcal{X}$-가군층이라 하자.
각 $f_j^{-1}\mathcal{F}$가 준연접이면 $\mathcal{F}$도 준연접이다.
\end{enumerate}
\end{lemma}

\begin{proof}
(1)의 증명. 각 대수적 스택 $\mathcal{X}_j$를 스킴 $U_j$로 바꾸어도 된다
(모든 대수적 스택은 스킴에 의한 매끄러운 피복을 가지며 매끄러운 사상들의
합성이 매끄럽다는 사실을 사용한다. 『스택의 사상』의 보조정리
\ref{stacks-morphisms-lemma-composition-smooth}를 보라).
$\mathcal{F}$를 $(\Sch/U_j)_\etale$로 당긴 것은 여전히 준연접이다.
『사이트 위의 가군』의 보조정리
\ref{sites-modules-lemma-local-pullback}을 보라.
그러면 $f = \coprod f_j : U = \coprod U_j \to \mathcal{X}$는
매끄러운 전사 사상이다. $x : V \to \mathcal{X}$를
$\mathcal{X}$의 대상이라 하자. 『스택 위의 층』의 보조정리
\ref{stacks-sheaves-lemma-surjective-flat-locally-finite-presentation}에 의해
각 $x_i$가 $(\Sch/U)_\etale$의 대상 $u_i$로 올려지는 에탈 피복
$\{x_i \to x\}_{i \in I}$가 존재한다. 이는 $x_i$가 스킴 $V_i$ 위에
놓이고, $\{V_i \to V\}$가 에탈 피복이며, $x_i$가 사상
$u_i : V_i \to U$에서 온다는 뜻일 뿐이다. 그러면
$x_i^*\mathcal{F} = u_i^*f^*\mathcal{F}$는 준연접이다.
예를 들어 『사이트 위의 가군』의 보조정리
\ref{sites-modules-lemma-local-final-object}에 의해, 이는
$(\Sch/V)_\etale$ 위의 $x^*\mathcal{F}$가 준연접임을 함의한다.
『스택 위의 층』의 보조정리
\ref{stacks-sheaves-lemma-characterize-quasi-coherent-bis}에 의해
$x^*\mathcal{F}$가 fppf 층임을 알 수 있고, $x$는 임의였으므로
$\mathcal{F}$가 fppf 위상의 층임을 알 수 있다.
『스택 위의 층』의 보조정리
\ref{stacks-sheaves-lemma-characterize-quasi-coherent}를 적용하면
$\mathcal{F}$가 준연접임을 알 수 있다.

\medskip\noindent
(2)의 증명. 정확히 같은 논증으로 증명하지만 여기서 전부 적겠다.
각 대수적 스택 $\mathcal{X}_j$를 스킴 $U_j$로 바꾸어도 된다
(모든 대수적 스택은 스킴에 의한 매끄러운 피복을 가지며 평탄하고 국소
유한 표시인 사상들이 합성에 의해 보존된다는 사실을 사용한다.
『스택의 사상』의 보조정리
\ref{stacks-morphisms-lemma-composition-flat}과
\ref{stacks-morphisms-lemma-composition-finite-presentation}을 보라).
$\mathcal{F}$를 $(\Sch/U_j)_\etale$로 당긴 것은 여전히 국소 준연접이다.
『스택 위의 층』의 보조정리
\ref{stacks-sheaves-lemma-pullback-quasi-coherent}를 보라.
그러면 $f = \coprod f_j : U = \coprod U_j \to \mathcal{X}$는
전사이고 평탄하며 국소 유한 표시이다. $x : V \to \mathcal{X}$를
$\mathcal{X}$의 대상이라 하자. 『스택 위의 층』의 보조정리
\ref{stacks-sheaves-lemma-surjective-flat-locally-finite-presentation}에 의해
각 $x_i$가 $(\Sch/U)_\etale$의 대상 $u_i$로 올려지는 fppf 피복
$\{x_i \to x\}_{i \in I}$가 존재한다. 이는 $x_i$가 스킴 $V_i$ 위에
놓이고, $\{V_i \to V\}$가 fppf 피복이며, $x_i$가 사상
$u_i : V_i \to U$에서 온다는 뜻일 뿐이다. 그러면
$x_i^*\mathcal{F} = u_i^*f^*\mathcal{F}$는 준연접이다.
예를 들어 『사이트 위의 가군』의 보조정리
\ref{sites-modules-lemma-local-final-object}에 의해, 이는
$(\Sch/V)_\etale$ 위의 $x^*\mathcal{F}$가 준연접임을 함의한다.
『스택 위의 층』의 보조정리
\ref{stacks-sheaves-lemma-characterize-quasi-coherent}에 의해
$\mathcal{F}$가 준연접임을 알 수 있다.
\end{proof}

\noindent
『사이트 위의 가군』의 절 \ref{sites-modules-section-local}에서 임의의
환 달린 토포스 위의 준연접 가군이라는 개념을 정의했음을 상기하자.

\begin{lemma}
\label{stacks-cohomology-lemma-shriek-quasi-coherent}
$\mathcal{X}$를 대수적 스택이라 하자. 표기법은 보조정리
\ref{stacks-cohomology-lemma-lisse-etale}에서와 같다.
\begin{enumerate}
\item $\mathcal{H}$를 $\mathcal{X}$의 리스-에탈 사이트 위의 준연접
$\mathcal{O}_{\mathcal{X}_{lisse,\etale}}$-가군이라 하자.
그러면 $g_!\mathcal{H}$는 $\mathcal{X}$ 위의 준연접 가군이다.
\item $\mathcal{H}$를 $\mathcal{X}$의 평탄-fppf 사이트 위의 준연접
$\mathcal{O}_{\mathcal{X}_{flat,fppf}}$-가군이라 하자.
그러면 $g_!\mathcal{H}$는 $\mathcal{X}$ 위의 준연접 가군이다.
\end{enumerate}
\end{lemma}

\begin{proof}
스킴 $U$와 매끄러운 전사 사상 $x : U \to \mathcal{X}$를 택한다.
『사이트 위의 가군』의 정의 \ref{sites-modules-definition-site-local}에 의해,
각 당김 $f_i^{-1}\mathcal{H}$가 대역 표시를 가지는 에탈
(각각 fppf) 피복 $\{U_i \to U\}_{i \in I}$가 존재한다
(『사이트 위의 가군』의 정의 \ref{sites-modules-definition-global}을 보라).
여기서 $f_i : U_i \to \mathcal{X}$는 합성 $U_i \to U \to \mathcal{X}$이고,
이는 대수적 스택의 사상이다. (당김이 $\mathcal{X}/f_i$에 대한 제한과
``같다''는 것을 상기하라. 『스택 위의 층』의 정의
\ref{stacks-sheaves-definition-pullback}과 이어지는 논의를 보라.)
보조정리 \ref{stacks-cohomology-lemma-lisse-etale-functorial}에 의해 각 $f_i$가 매끄러우므로
(각각 평탄하므로)
$f_i^{-1}g_!\mathcal{H} = g_{i, !}(f'_i)^{-1}\mathcal{H}$이다.
보조정리 \ref{stacks-cohomology-lemma-check-qc-on-etale-covering}을 사용하면
$\mathcal{H}$가 대역 표시를 가지는 경우로 보조정리의 명제를 환원할 수 있다.
즉, $\mathcal{O}$-가군의 열
$$
\bigoplus\nolimits_{j \in J} \mathcal{O} \longrightarrow
\bigoplus\nolimits_{i \in I} \mathcal{O} \longrightarrow
\mathcal{H} \longrightarrow 0
$$
이 있다고 하자. 여기서
$\mathcal{O} = \mathcal{O}_{\mathcal{X}_{lisse,\etale}}$
(각각 $\mathcal{O} = \mathcal{O}_{\mathcal{X}_{flat,fppf}}$)이다.
$g_!$은 (왼쪽 수반 함자이므로) 임의의 쌍대극한과 가환한다
(보조정리 \ref{stacks-cohomology-lemma-lisse-etale-modules}와 『범주』의 보조정리
\ref{categories-lemma-adjoint-exact}을 보라). 따라서 완전열
$$
\bigoplus\nolimits_{j \in J} g_!\mathcal{O} \longrightarrow
\bigoplus\nolimits_{i \in I} g_!\mathcal{O} \longrightarrow
g_!\mathcal{H} \longrightarrow 0
$$
이 존재한다. 보조정리 \ref{stacks-cohomology-lemma-lisse-etale-structure-sheaf}에 의해
$g_!\mathcal{O} = \mathcal{O}_\mathcal{X}$이다. (2)의 경우에는 끝났다.
(1)의 경우에는 『스택 위의 층』의 보조정리
\ref{stacks-sheaves-lemma-characterize-quasi-coherent-bis}를 적용하여 결론을 얻는다.
\end{proof}

\begin{lemma}
\label{stacks-cohomology-lemma-quasi-coherent}
$\mathcal{X}$를 대수적 스택이라 하자.
\begin{enumerate}
\item 리스-에탈 사이트에 대해 보조정리 \ref{stacks-cohomology-lemma-lisse-etale}에서와 같은
$g$를 사용하면 다음이 성립한다.
\begin{enumerate}
\item 함자 $g^{-1}$과 $g_!$은 서로 역인 함자
$$
\xymatrix{
\QCoh(\mathcal{O}_\mathcal{X}) \ar@<1ex>[r]^-{g^{-1}} &
\QCoh(\mathcal{X}_{lisse,\etale},
\mathcal{O}_{\mathcal{X}_{lisse,\etale}}) \ar@<1ex>[l]^-{g_!}
}
$$
를 정의한다.
\item $\mathcal{F}$가
$\textit{LQCoh}^{fbc}(\mathcal{O}_\mathcal{X})$에 속하면
$g^{-1}\mathcal{F}$는
$\QCoh(\mathcal{O}_{\mathcal{X}_{lisse,\etale}})$에 속한다.
\item $Q$가 보조정리 \ref{stacks-cohomology-lemma-adjoint}에서와 같을 때
$Q(\mathcal{F}) = g_!g^{-1}\mathcal{F}$이다.
\end{enumerate}
\item 평탄-fppf 사이트에 대해 보조정리 \ref{stacks-cohomology-lemma-lisse-etale}에서와 같은
$g$를 사용하면 다음이 성립한다.
\begin{enumerate}
\item 함자 $g^{-1}$과 $g_!$은 서로 역인 함자
$$
\xymatrix{
\QCoh(\mathcal{O}_\mathcal{X}) \ar@<1ex>[r]^-{g^{-1}} &
\QCoh(\mathcal{X}_{flat,fppf},
\mathcal{O}_{\mathcal{X}_{flat,fppf}}) \ar@<1ex>[l]^-{g_!}
}
$$
를 정의한다.
\item $\mathcal{F}$가
$\textit{LQCoh}^{fbc}(\mathcal{O}_\mathcal{X})$에 속하면
$g^{-1}\mathcal{F}$는
$\QCoh(\mathcal{O}_{\mathcal{X}_{flat,fppf}})$에 속한다.
\item $Q$가 보조정리 \ref{stacks-cohomology-lemma-adjoint}에서와 같을 때
$Q(\mathcal{F}) = g_!g^{-1}\mathcal{F}$이다.
\end{enumerate}
\end{enumerate}
\end{lemma}

\begin{proof}
환 달린 토포스의 임의의 사상에 의한 당김은 준연접 가군들의 범주를
보존한다. 『사이트 위의 가군』의 보조정리
\ref{sites-modules-lemma-local-pullback}을 보라. 따라서 $g^{-1}$은
준연접 가군들의 범주를 보존한다. 여기서 『스택 위의 층』의 보조정리
\ref{stacks-sheaves-lemma-characterize-quasi-coherent-bis}에 의해
$\QCoh(\mathcal{O}_\mathcal{X}) =
\QCoh(\mathcal{X}_\etale, \mathcal{O}_\mathcal{X})$라는 사실을 사용한다.
보조정리 \ref{stacks-cohomology-lemma-shriek-quasi-coherent}에 의해 $g_!$에 대해서도 같다.
보조정리 \ref{stacks-cohomology-lemma-lisse-etale}에 의해
$\mathcal{H} \to g^{-1}g_!\mathcal{H}$는 동형이다.
반대로 $\mathcal{F}$가 $\QCoh(\mathcal{O}_\mathcal{X})$에 속하면
사상 $g_!g^{-1}\mathcal{F} \to \mathcal{F}$는 $\mathcal{X}$ 위의 준연접
가군들의 사상이고, $\mathcal{X}$ 위의 임의의 매끄러운 스킴으로 제한하면
동형이 된다. 그러면 『스택 위의 층』의 절
\ref{stacks-sheaves-section-quasi-coherent-presentation}과
\ref{stacks-sheaves-section-quasi-coherent-algebraic-stacks}의 논의
(표시 위의 준연접 가군과 비교한다)에 의해 이 사상은 동형이다.
이로써 (1)(a)와 (2)(a)가 증명되었다.

\medskip\noindent
$\mathcal{F}$를
$\textit{LQCoh}^{fbc}(\mathcal{O}_\mathcal{X})$의 대상이라 하자.
보조정리 \ref{stacks-cohomology-lemma-adjoint-kernel-parasitic}에 의해 사상
$Q(\mathcal{F}) \to \mathcal{F}$의 핵과 여핵은 기생이다.
따라서 보조정리 \ref{stacks-cohomology-lemma-parasitic-in-terms-flat-fppf}와
$g^* = g^{-1}$이 완전하다는 사실에 의해
$g^*Q(\mathcal{F}) \to g^*\mathcal{F}$가 동형이라는 결론을 얻는다.
그러므로 $g^*\mathcal{F}$는 준연접이다. 이로써 (1)(b)와 (2)(b)가 증명되었다.
마지막으로 위 논증에 의해
$g_!g^*Q(\mathcal{F}) \to Q(\mathcal{F})$가 동형이므로
(1)(c)와 (2)(c)가 따른다.
\end{proof}

\begin{lemma}
\label{stacks-cohomology-lemma-quasi-coherent-weak-serre}
$\mathcal{X}$를 대수적 스택이라 하자.
\begin{enumerate}
\item $\QCoh(\mathcal{O}_{\mathcal{X}_{lisse,\etale}})$는
$\textit{Mod}(\mathcal{O}_{\mathcal{X}_{lisse,\etale}})$의
약한 Serre 부분범주이다.
\item $\QCoh(\mathcal{O}_{\mathcal{X}_{flat,fppf}})$는
$\textit{Mod}(\mathcal{O}_{\mathcal{X}_{flat,fppf}})$의
약한 Serre 부분범주이다.
\end{enumerate}
\end{lemma}

\begin{proof}
『호몰로지』의 보조정리
\ref{homology-lemma-characterize-weak-serre-subcategory}의 조건
(1), (2), (3), (4)를 확인하겠다.

\medskip\noindent
$0$은 임의의 환 달린 사이트 위에서 준연접 가군이므로 (1)이 성립한다.

\medskip\noindent
정의에 따라 $\QCoh(\mathcal{O})$는
$\textit{Mod}(\mathcal{O})$의 엄밀 충만한 부분범주이므로 (2)가 성립한다.

\medskip\noindent
$\varphi : \mathcal{G} \to \mathcal{F}$를
$\mathcal{X}_{lisse,\etale}$ 또는 $\mathcal{X}_{flat,fppf}$ 위의
준연접 가군들의 사상이라 하자. 보조정리 \ref{stacks-cohomology-lemma-lisse-etale-modules}에 의해
$g^*g_!\mathcal{F} = \mathcal{F}$이고, $\mathcal{G}$와 $\varphi$에 대해서도
마찬가지이다. 보조정리 \ref{stacks-cohomology-lemma-shriek-quasi-coherent}에 의해
$g_!\mathcal{F}$와 $g_!\mathcal{G}$는 준연접
$\mathcal{O}_\mathcal{X}$-가군이다. 『스택 위의 층』의 보조정리
\ref{stacks-sheaves-lemma-quasi-coherent-algebraic-stack}에 의해
$\Coker(g_!\varphi)$는 $\mathcal{X}$ 위의 준연접 가군이며
($\mathcal{X}$ 위의 준연접 가군 범주에서의 여핵이다).
$g^*$가 완전하므로(보조정리 \ref{stacks-cohomology-lemma-lisse-etale}를 보라)
$g^*\Coker(g_!\varphi) = \Coker(g^*g_!\varphi) = \Coker(\varphi)$도
준연접이다(보조정리 \ref{stacks-cohomology-lemma-quasi-coherent}를 보라).
명제 \ref{stacks-cohomology-proposition-loc-qcoh-flat-base-change}에 의해
핵 $\Ker(g_!\varphi)$는
$\textit{LQCoh}^{fbc}(\mathcal{O}_\mathcal{X})$에 속한다.
$g^*$가 완전하므로
$g^*\Ker(g_!\varphi) = \Ker(g^*g_!\varphi) = \Ker(\varphi)$이다.
보조정리 \ref{stacks-cohomology-lemma-quasi-coherent}에 의해 $g^*$는
$\textit{LQCoh}^{fbc}(\mathcal{O}_\mathcal{X})$의 대상을 준연접 가군으로
보내므로, $\Ker(\varphi)$도 준연접이라는 결론을 얻는다.
이로써 (3)이 증명되었다.

\medskip\noindent
마지막으로
$$
0 \to \mathcal{F} \to \mathcal{E} \to \mathcal{G} \to 0
$$
이 $\mathcal{O}_{\mathcal{X}_{lisse,\etale}}$-가군의 확대
(각각 $\mathcal{O}_{\mathcal{X}_{flat,fppf}}$-가군의 확대)이고
$\mathcal{F}$와 $\mathcal{G}$가 준연접이라고 하자.
(4)를 증명하여 전체 증명을 마치려면 $\mathcal{E}$가
$\mathcal{X}_{lisse,\etale}$ 위에서(각각 $\mathcal{X}_{flat,fppf}$ 위에서)
준연접임을 보여야 한다. $U$를 $\mathcal{X}_{lisse,\etale}$의 대상
(각각 $\mathcal{X}_{flat,fppf}$의 대상이라 하자. 우리는 $U$를
$\mathcal{X}$ 위에서 매끄러운(각각 평탄한) 스킴으로 생각한다.
$\mathcal{E}$를 $U_{lisse,\etale}$에 제한한 것
(각각 $=U_{flat,fppf}$)이 준연접임을 보여야 한다.
따라서 $\mathcal{X} = U$가 스킴이라고 가정해도 된다.
$\mathcal{G}$가 $U_{lisse,\etale}$ 위에서(각각 $U_{flat,fppf}$ 위에서)
준연접이므로, $U$를 어떤 에탈(각각 fppf) 피복의 원소들로 바꾸어
$\mathcal{G}$가 $U_{lisse,\etale}$ 위에서(각각 $U_{flat,fppf}$ 위에서)
표시
$$
\bigoplus\nolimits_{j \in J} \mathcal{O} \longrightarrow
\bigoplus\nolimits_{i \in I} \mathcal{O} \longrightarrow
\mathcal{G} \longrightarrow 0
$$
를 가진다고 가정해도 된다. 여기서 $\mathcal{O}$는 그 사이트 위의 구조층이다.
$U$가 아핀이라고 가정해도 된다. $\mathcal{F}$가 준연접이므로
$$
H^1(U_{lisse,\etale}, \mathcal{F}) = 0,
\quad\text{각각}\quad
H^1(U_{flat,fppf}, \mathcal{F}) = 0
$$
이다. 구체적으로 $\mathcal{F}$는 $U$의 큰 사이트 위의 어떤 준연접 가군
$\mathcal{F}'$의 당김이고(보조정리 \ref{stacks-cohomology-lemma-quasi-coherent}),
$\mathcal{F}$와 $\mathcal{F}'$의 코호몰로지는 일치하며
(보조정리 \ref{stacks-cohomology-lemma-lisse-etale-cohomology}), 아핀 스킴 $U$의 큰 사이트
위에서 $\mathcal{F}'$의 코호몰로지가 영임을 알고 있다
(현재 상황에서 이를 얻으려면 『하강』의 명제
\ref{descent-proposition-equivalence-quasi-coherent}와
\ref{descent-proposition-same-cohomology-quasi-coherent}를
『스킴의 코호몰로지』의 보조정리
\ref{coherent-lemma-quasi-coherent-affine-cohomology-zero}와 결합해야 한다).
따라서 사상 $\bigoplus_{i \in I} \mathcal{O} \to \mathcal{G}$를
$\mathcal{E}$로 올릴 수 있다. 도식 추적을 하면 완전열
$$
\bigoplus\nolimits_{j \in J} \mathcal{O} \to
\mathcal{F} \oplus \bigoplus\nolimits_{i \in I} \mathcal{O}
\to
\mathcal{E} \to 0
$$
을 얻는다. 위에서 증명한 (3)에 의해 원하는 대로 $\mathcal{E}$는 준연접이다.
\end{proof}

\section{국소 뇌터 스택 위의 연접층}
\label{stacks-cohomology-section-coherent-sheaves}

\noindent
이 절은 『대수공간의 코호몰로지』의 절
\ref{spaces-cohomology-section-coherent}에 대응한다.
『사이트 위의 가군』의 절 \ref{sites-modules-section-local}에서 임의의 환 달린
토포스 위의 연접 가군이라는 개념을 정의했다. 그러나 임의의 대수적 스택
$\mathcal{X}$에 대해 연접 $\mathcal{O}_\mathcal{X}$-가군의 범주는 영이다.
이는 본질적으로 사이트 $\mathcal{X}$가 뇌터가 아닌 대상을 너무 많이
포함하기 때문이다($\mathcal{X}$ 자체가 국소 뇌터인 경우에도 그렇다).
대신 다음 보조정리를 사용하여 연접 가군을 정의하겠다.

\begin{lemma}
\label{stacks-cohomology-lemma-coherent-Noetherian}
$\mathcal{X}$를 국소 뇌터 대수적 스택이라 하자.
$\mathcal{F}$를 $\mathcal{O}_\mathcal{X}$-가군이라 하자. 다음은 동치이다.
\begin{enumerate}
\item $\mathcal{F}$는 유한형 준연접 $\mathcal{O}_\mathcal{X}$-가군이다.
\item $\mathcal{F}$는 유한 표시 $\mathcal{O}_\mathcal{X}$-가군이다.
\item $\mathcal{F}$는 준연접이고, $U$가 국소 뇌터 대수공간인 임의의
사상 $f : U \to \mathcal{X}$에 대해 당김
$f^*\mathcal{F}|_{U_\etale}$는 연접이다.
\item $\mathcal{F}$는 준연접이고, 국소 유한형이고 평탄하며 전사인 사상
$f : U \to \mathcal{X}$를 가지는 대수공간 $U$가 존재하여 당김
$f^*\mathcal{F}|_{U_\etale}$가 연접이다.
\end{enumerate}
\end{lemma}

\begin{proof}
$f : U \to \mathcal{X}$가 (4)에서와 같다고 하자.
그러면 $U$는 국소 뇌터이고(『스택의 사상』의 보조정리
\ref{stacks-morphisms-lemma-locally-finite-type-locally-noetherian}),
따라서 보조정리의 명제는 의미가 있다. 더욱이 『스택의 사상』의 보조정리
\ref{stacks-morphisms-lemma-noetherian-finite-type-finite-presentation}에 의해
$f$는 국소 유한 표시이다. $x$를 스킴 $V$ 위에 놓인 $\mathcal{X}$의
대상이라 하자. (2)를 증명하려면 $V$를 $V$의 어떤 fppf 피복의 원소들로 바꾼 뒤
제한 $x^*\mathcal{F}$가
$\mathcal{X}/x \cong (\Sch/V)_{fppf}$ 위에서 대역 유한 표시를 가짐을
보여야 한다. 사영 $W = U \times_\mathcal{X} V \to V$는 국소 유한 표시이고,
평탄하며, 전사이다. 따라서 $V$를 $W$의 어떤 에탈 피복의 스킴 원소들로
바꾸어 $f \circ h = x$를 만족하는 사상 $h : V \to U$가 있다고 가정해도
된다. $\mathcal{F}$가 준연접이므로, 『스택 위의 층』의 보조정리
\ref{stacks-sheaves-lemma-compare-quasi-coherent}에 의해 제한
$x^*\mathcal{F}$는
$h_{small}^*(f^*\mathcal{F})|_{U_\etale}$를 $\pi_V$로 당긴 것이다.
가정에 따라 $f^*\mathcal{F}|_{U_\etale}$는 에탈 위상에서 국소적으로
유한 표시를 가지므로 (4) $\Rightarrow$ (2)를 얻는다.

\medskip\noindent
(2)는 임의의 환 달린 토포스에서 (1)을 함의한다(정의에서 즉시 따른다).
성질 ``유한형''과 ``준연접''은 환 달린 토포스의 임의의 사상에 의한
당김 아래에서 보존된다. 『사이트 위의 가군』의 보조정리
\ref{sites-modules-lemma-local-pullback}을 보라. 따라서 (1)은 (3)을 함의한다.
『대수공간의 코호몰로지』의 보조정리
\ref{spaces-cohomology-lemma-coherent-Noetherian}을 보라.
마지막으로 (3)은 자명하게 (4)를 함의한다.
\end{proof}

\begin{definition}
\label{stacks-cohomology-definition-coherent}
$\mathcal{X}$를 국소 뇌터 대수적 스택이라 하자.
$\mathcal{O}_\mathcal{X}$-가군 $\mathcal{F}$를 {\it 연접}이라고 한다.
이는 $\mathcal{F}$가 보조정리 \ref{stacks-cohomology-lemma-coherent-Noetherian}의 동치 조건들
가운데 하나를(따라서 모두를) 만족한다는 뜻이다.
연접 $\mathcal{O}_\mathcal{X}$-가군들의 범주는
$\textit{Coh}(\mathcal{O}_\mathcal{X})$로 표기된다.
\end{definition}

\begin{lemma}
\label{stacks-cohomology-lemma-elementary-coherent}
$\mathcal{X}$를 국소 뇌터 대수적 스택이라 하자.
가군 $\mathcal{O}_\mathcal{X}$는 연접이고, 모든 가역
$\mathcal{O}_\mathcal{X}$-가군은 연접이며, 더 일반적으로 모든 유한 국소 자유
$\mathcal{O}_\mathcal{X}$-가군은 연접이다.
\end{lemma}

\begin{proof}
정의와 『대수공간의 코호몰로지』의 보조정리
\ref{spaces-cohomology-lemma-coherent-Noetherian}에서 따른다.
\end{proof}

\begin{lemma}
\label{stacks-cohomology-lemma-pullback-coherent}
$f : \mathcal{X} \to \mathcal{Y}$를 국소 뇌터 대수적 스택의 사상이라 하자.
그러면 $f^*$는 $\mathcal{Y}$ 위의 연접 가군을
$\mathcal{X}$ 위의 연접 가군으로 보낸다.
\end{lemma}

\begin{proof}
정의와, 환 달린 토포스의 임의의 사상에 대한 당김이 유한 표시 가군을
보존한다는 사실에서 즉시 따른다. 『사이트 위의 가군』의 보조정리
\ref{sites-modules-lemma-local-pullback}을 보라.
\end{proof}

\begin{lemma}
\label{stacks-cohomology-lemma-coherent-abelian-Noetherian}
$\mathcal{X}$를 국소 뇌터 대수적 스택이라 하자.
연접 $\mathcal{O}_\mathcal{X}$-가군들의 범주는 아벨 범주이다.
$\varphi : \mathcal{F} \to \mathcal{G}$가 연접
$\mathcal{O}_\mathcal{X}$-가군들의 사상이면 다음이 성립한다.
\begin{enumerate}
\item $\textit{Mod}(\mathcal{O}_\mathcal{X})$에서 계산한 여핵
$\Coker(\varphi)$는 연접 $\mathcal{O}_\mathcal{X}$-가군이다.
\item $\textit{Mod}(\mathcal{O}_\mathcal{X})$에서 계산한 상
$\Im(\varphi)$는 연접 $\mathcal{O}_\mathcal{X}$-가군이다.
\item $\textit{Mod}(\mathcal{O}_\mathcal{X})$에서 계산한 핵
$\Ker(\varphi)$는 연접이 아닐 수 있지만
$\textit{LQCoh}^{fbc}(\mathcal{O}_\mathcal{X})$에 속한다.
$Q(\Ker(\varphi))$는 연접이고
$\textit{Coh}(\mathcal{O}_\mathcal{X})$에서의 $\varphi$의 핵이다.
\end{enumerate}
포함 함자 $\textit{Coh}(\mathcal{O}_\mathcal{X}) \to
\QCoh(\mathcal{O}_\mathcal{X})$는 완전하다.
\end{lemma}

\begin{proof}
$\textit{Coh}(\mathcal{O}_\mathcal{X})$에서 핵, 상, 여핵을 취하는 규칙은
주 \ref{stacks-cohomology-remark-QCoh-abelian}의 준연접 가군에 대한 처방과 일치한다.
따라서 준연접 가군 $\Coker(\varphi)$, $\Im(\varphi)$,
$Q(\Ker(\varphi))$가 연접임을 보이면 보조정리가 따른다.
보조정리 \ref{stacks-cohomology-lemma-coherent-Noetherian}에 의해, 어떤 매끄러운 전사 사상
$f : U \to \mathcal{X}$에 대해 $U_\etale$로 제한한 뒤 이를 증명하면 충분하다.
함자 $\mathcal{F} \mapsto f^*\mathcal{F}|_{U_\etale}$는 완전하다.
따라서 $f^*\Coker(\varphi)$와 $f^*\Im(\varphi)$는 연접
$\mathcal{O}_U$-가군들 사이의 사상의 여핵과 상이고, 따라서 원하는 대로
연접이다. 보조정리 \ref{stacks-cohomology-lemma-parasitic}에 의해 함자
$\mathcal{F} \mapsto f^*\mathcal{F}|_{U_\etale}$는 기생 가군을 영으로 보낸다.
따라서 보조정리 \ref{stacks-cohomology-lemma-adjoint-kernel-parasitic}의 (2)에 의해
$f^*Q(\Ker(\varphi))|_{U_\etale} = f^*\Ker(\varphi)|_{U_\etale}$이다.
그러므로 같은 방식으로 $Q(\Ker(\varphi))$가 연접이라는 결론을 얻는다.
\end{proof}

\begin{lemma}
\label{stacks-cohomology-lemma-extension-coherent}
$\mathcal{X}$를 국소 뇌터 대수적 스택이라 하자.
$\textit{Mod}(\mathcal{O}_\mathcal{X})$ 안의 짧은 완전열
$0 \to \mathcal{F}_1 \to \mathcal{F}_2 \to \mathcal{F}_3 \to 0$이
주어지고 $\mathcal{F}_1$과 $\mathcal{F}_3$가 연접이면,
$\mathcal{F}_2$는 연접이다.
\end{lemma}

\begin{proof}
『스택 위의 층』의 보조정리
\ref{stacks-sheaves-lemma-quasi-coherent-algebraic-stack}의 (7)에 의해
$\mathcal{F}_2$가 준연접임을 알 수 있다. 그러면 어떤 매끄러운 전사 사상
$U \to \mathcal{X}$에 대해 $U_\etale$로 제한하여
$\mathcal{F}_2$가 연접임을 확인할 수 있다. 이는 『대수공간의 코호몰로지』의
보조정리 \ref{spaces-cohomology-lemma-coherent-abelian-Noetherian}에서 따른다.
일부 세부사항은 생략한다.
\end{proof}

\noindent
연접 가군들은 준연접 $\mathcal{O}_\mathcal{X}$-가군들의 범주의
Serre 부분범주를 이룬다. 이는 일반적인 환 달린 토포스 위의 가군에
대해서는 성립하지 않는다.

\begin{lemma}
\label{stacks-cohomology-lemma-coherent-Noetherian-quasi-coherent-sub-quotient}
$\mathcal{X}$를 국소 뇌터 대수적 스택이라 하자. 그러면
$\textit{Coh}(\mathcal{O}_\mathcal{X})$는
$\QCoh(\mathcal{O}_\mathcal{X})$의 Serre 부분범주이다.
$\varphi : \mathcal{F} \to \mathcal{G}$를 준연접
$\mathcal{O}_\mathcal{X}$-가군들의 사상이라 하자. 다음이 성립한다.
\begin{enumerate}
\item $\mathcal{F}$가 연접이고 $\varphi$가 전사이면 $\mathcal{G}$는 연접이다.
\item $\mathcal{F}$가 연접이면 $\Im(\varphi)$는 연접이다.
\item $\mathcal{G}$가 연접이고 $\Ker(\varphi)$가 기생이면
$\mathcal{F}$는 연접이다.
\end{enumerate}
\end{lemma}

\begin{proof}
스킴 $U$와 매끄러운 전사 사상 $f : U \to \mathcal{X}$를 택한다.
그러면 함자
$f^* : \QCoh(\mathcal{O}_\mathcal{X}) \to \QCoh(\mathcal{O}_U)$는
완전하고(보조정리 \ref{stacks-cohomology-lemma-flat-pullback-quasi-coherent}),
더욱이 정의에 따라 $\textit{Coh}(\mathcal{O}_\mathcal{X})$는
$f^*\mathcal{F}$가 $\textit{Coh}(\mathcal{O}_U)$에 속하는 대상
$\mathcal{F}$들로 이루어진
$\QCoh(\mathcal{O}_\mathcal{X})$의 충만한 부분범주이다.
$\textit{Coh}(\mathcal{O}_\mathcal{X})$가
$\QCoh(\mathcal{O}_\mathcal{X})$의 Serre 부분범주라는 명제는 이 사실과
$U$에 대한 대응하는 사실에서 즉시 따른다. 『대수공간의 코호몰로지』의
보조정리 \ref{spaces-cohomology-lemma-coherent-abelian-Noetherian}와
\ref{spaces-cohomology-lemma-coherent-Noetherian-quasi-coherent-sub-quotient}를
보라. (1), (2), (3)의 증명은 생략한다. 힌트: 보조정리
\ref{stacks-cohomology-lemma-coherent-abelian-Noetherian}의 증명과 비교하라.
\end{proof}

\noindent
$\mathcal{X}$를 국소 뇌터 대수적 스택이라 하자. $U$를 대수공간이라 하고
$f : U \to \mathcal{X}$를 전사이고 국소 유한 표시이며 평탄하다고 하자.
$U$는 국소 뇌터임에 유의하라(『스택의 사상』의 보조정리
\ref{stacks-morphisms-lemma-locally-finite-type-locally-noetherian}).
『대수적 스택』의 보조정리
\ref{algebraic-lemma-map-space-into-stack}과 주
\ref{algebraic-remark-flat-fp-presentation}에서 구성한 대수공간에서의 준군을
$(U, R, s, t, c)$라 하고 동형을
$f_{can} : [U/R] \to \mathcal{X}$라 하자.
『스택 위의 층』의 절
\ref{stacks-sheaves-section-quasi-coherent-algebraic-stacks}에서처럼 동치
$$
\QCoh(\mathcal{O}_\mathcal{X})
\cong
\QCoh(\mathcal{O}_{[U/R]})
\cong
\QCoh(U, R, s, t, c)
$$
를 얻는다. 여기서 두 번째 동치는 『스택 위의 층』의 명제
\ref{stacks-sheaves-proposition-quasi-coherent}이다.
『공간에서의 준군』의 절 \ref{spaces-groupoids-section-colimits}에서
{\it 연접 가군}들의 충만한 부분범주
$$
\textit{Coh}(U, R, s, t, c) \subset \QCoh(U, R, s, t, c)
$$
를, $\mathcal{G}$가 연접 $\mathcal{O}_U$-가군인
$(\mathcal{G}, \alpha)$들로 정의했음을 상기하라.

\begin{lemma}
\label{stacks-cohomology-lemma-coherent-presentation}
위에서 논의한 상황에서 동치
$\QCoh(\mathcal{O}_\mathcal{X}) \cong \QCoh(U, R, s, t, c)$는
연접층을 연접층으로 보내며 그 역도 성립한다. 즉, 동치
$\textit{Coh}(\mathcal{O}_\mathcal{X}) \cong
\textit{Coh}(U, R, s, t, c)$를 유도한다.
\end{lemma}

\begin{proof}
이는 연접 $\mathcal{O}_\mathcal{X}$-가군의 정의에서 즉시 따른다.
참고문헌을 기록해 두면, 위 내용은 『스택의 사상』의 보조정리
\ref{stacks-morphisms-lemma-locally-finite-type-locally-noetherian},
『대수적 스택』의 보조정리 \ref{algebraic-lemma-map-space-into-stack}과 주
\ref{algebraic-remark-flat-fp-presentation},
『스택 위의 층』의 절
\ref{stacks-sheaves-section-quasi-coherent-algebraic-stacks},
『스택 위의 층』의 명제 \ref{stacks-sheaves-proposition-quasi-coherent},
그리고 『공간에서의 준군』의 절
\ref{spaces-groupoids-section-colimits}을 사용한다.
\end{proof}

\begin{lemma}
\label{stacks-cohomology-lemma-coherent-hom}
$\mathcal{X}$를 국소 뇌터 대수적 스택이라 하자.
$\mathcal{F}$와 $\mathcal{G}$를 연접
$\mathcal{O}_\mathcal{X}$-가군이라 하자. 보조정리
\ref{stacks-cohomology-lemma-internal-hom-fp-into-qcoh}에서 구성한 내부 Hom
$hom(\mathcal{F}, \mathcal{G})$은 연접 $\mathcal{O}_\mathcal{X}$-가군이다.
\end{lemma}

\begin{proof}
$U \to \mathcal{X}$를 스킴에서 오는 매끄러운 전사 사상이라 하자.
항목 (\ref{stacks-cohomology-item-hom-restriction})(절
\ref{stacks-cohomology-section-further-remarks})에 의해
$hom(\mathcal{F}, \mathcal{G})$를 $U$에 제한한 것은 제한들의 Hom 층이다.
따라서 이 보조정리는 대수공간의 경우에서 따른다. 『대수공간의 코호몰로지』의
보조정리 \ref{spaces-cohomology-lemma-tensor-hom-coherent}를 보라.
\end{proof}

\section{뇌터 스택 위의 연접층}
\label{stacks-cohomology-section-coherent-on-noetherian}

\noindent
이 절은 『대수공간의 코호몰로지』의 절
\ref{spaces-cohomology-section-coherent-quasi-compact}에 대응한다.

\begin{lemma}
\label{stacks-cohomology-lemma-directed-colimit-coherent}
$\mathcal{X}$를 뇌터 대수적 스택이라 하자. 모든 준연접
$\mathcal{O}_\mathcal{X}$-가군은 그 연접 부분가군들의 여과 쌍대극한이다.
\end{lemma}

\begin{proof}
$\mathcal{F}$를 준연접 $\mathcal{O}_\mathcal{X}$-가군이라 하자.
$\mathcal{G}, \mathcal{H} \subset \mathcal{F}$가 연접
$\mathcal{O}_\mathcal{X}$-부분가군이면,
$\mathcal{G} \oplus \mathcal{H} \to \mathcal{F}$의 상은 둘을 모두
포함하는 또 다른 연접 $\mathcal{O}_\mathcal{X}$-부분가군이다.
보조정리 \ref{stacks-cohomology-lemma-coherent-Noetherian-quasi-coherent-sub-quotient}를 보라.
따라서 이 계가 유향임을 알 수 있다. 이제 $\mathcal{F}$가 연접 가군들의
여과 쌍대극한으로 쓰임을 보이면 충분하다. 그러면 이 가군들의
$\mathcal{F}$ 안에서의 상을 취하여 충분히 많이 존재한다는 결론을
내릴 수 있기 때문이다.

\medskip\noindent
$U$를 아핀 스킴이라 하고 $U \to \mathcal{X}$를 매끄러운 전사 사상이라 하자
(『스택의 성질』의 보조정리
\ref{stacks-properties-lemma-quasi-compact-stack}).
$R = U \times_\mathcal{X} U$라 놓으면 『대수적 스택』의 보조정리
\ref{algebraic-lemma-stack-presentation}에서처럼 $\mathcal{X} = [U/R]$이다.
보조정리 \ref{stacks-cohomology-lemma-coherent-presentation}에 의해
$\QCoh(\mathcal{O}_X) = \QCoh(U, R, s, t, c)$이고
$\textit{Coh}(\mathcal{O}_X) = \textit{Coh}(U, R, s, t, c)$이다.
이렇게 하여 $\QCoh(U, R, s, t, c)$에 대해 대응하는 사실을 증명하는 문제로
환원된다. 이는 『공간에서의 준군』의 보조정리
\ref{spaces-groupoids-lemma-colimit-coherent}이다. 다음 문단에서 그 가정들을
확인한다.

\medskip\noindent
독자에게 증명의 나머지는 건너뛰기를 권한다.
아핀 스킴 $U$는 뇌터이다. 이는 $\mathcal{X}$가 국소 뇌터라는 우리의
정의에서 따른다. 『스택의 성질』의 정의
\ref{stacks-properties-definition-type-property}와 주
\ref{stacks-properties-remark-list-properties-local-smooth-topology}를 보라.
사영 사상 $s, t : R \to U$는 매끄럽고(위에서 제시한 참고문헌을 보라),
준분리이며 준콤팩트이다(『스택의 사상』의 보조정리
\ref{stacks-morphisms-lemma-quasi-compact-quasi-separated-permanence}).
특히 $R$은 $U$ 위에서 매끄러운 준콤팩트 준분리 대수공간이고,
따라서 뇌터이다(『대수공간의 사상』의 보조정리
\ref{spaces-morphisms-lemma-finite-presentation-noetherian}).
\end{proof}
